% =====================================================================
%  RedVital - Documento de Arquitectura de Software (SAD) - Version 2.0
%  Arkhe Software S.A.S.
%  Compilar con pdfLaTeX (Overleaf: Menu > Compiler > pdfLaTeX)
%  Requiere DOS compilaciones para resolver indice y referencias cruzadas.
% =====================================================================
\documentclass[11pt,a4paper]{article}

% ---------- Codificacion e idioma ----------
\usepackage[utf8]{inputenc}
\usepackage[T1]{fontenc}
\usepackage[spanish,es-noshorthands,es-tabla]{babel}
\usepackage{lmodern}
\usepackage{amsmath}
\usepackage{microtype}

% ---------- Pagina ----------
\usepackage[a4paper,top=2.6cm,bottom=2.6cm,left=2.5cm,right=2.5cm]{geometry}
\usepackage{fancyhdr}
\setlength{\headheight}{14pt}
\usepackage{lastpage}

% ---------- Tablas ----------
\usepackage{array}
\usepackage{longtable}
\usepackage{booktabs}
\usepackage[table]{xcolor}
\usepackage{ragged2e}

% ---------- Graficos ----------
\usepackage{graphicx}
\usepackage{tikz}
\usetikzlibrary{positioning,arrows.meta,fit,backgrounds,calc,shapes.geometric}

% ---------- Varios ----------
\usepackage{enumitem}
\usepackage{needspace}
\usepackage{titlesec}
\usepackage{caption}
\usepackage[hidelinks]{hyperref}

% =====================================================================
%  Identidad visual
% =====================================================================
\definecolor{vinotinto}{RGB}{139,26,36}
\definecolor{vinoclaro}{RGB}{190,90,95}
\definecolor{grisTexto}{RGB}{90,90,90}

\definecolor{ringAdoptar}{RGB}{139,26,36}
\definecolor{ringProbar}{RGB}{200,120,40}
\definecolor{ringEvaluar}{RGB}{70,110,150}
\definecolor{ringEvitar}{RGB}{120,120,120}

\newcommand{\adoptar}{\textcolor{ringAdoptar}{\textbf{Adoptar}}}
\newcommand{\probar}{\textcolor{ringProbar}{\textbf{Probar}}}
\newcommand{\evaluar}{\textcolor{ringEvaluar}{\textbf{Evaluar}}}
\newcommand{\evitar}{\textcolor{ringEvitar}{\textbf{Evitar}}}

% Niveles de prioridad y estado, con el mismo criterio cromatico del radar
\newcommand{\alta}{\textcolor{ringAdoptar}{\textbf{Alta}}}
\newcommand{\media}{\textcolor{ringProbar}{\textbf{Media}}}
\newcommand{\baja}{\textcolor{ringEvitar}{\textbf{Baja}}}

% ---------- Encabezado y pie ----------
\pagestyle{fancy}
\fancyhf{}
\fancyhead[L]{\footnotesize\textcolor{grisTexto}{RedVital --- Documento de Arquitectura de Software (SAD) V2.0}}
\fancyhead[R]{\footnotesize\textcolor{grisTexto}{Arkhé Software S.A.S.}}
\fancyfoot[C]{\footnotesize Página \thepage{} de \pageref{LastPage}}
\renewcommand{\headrulewidth}{0.4pt}
\renewcommand{\headrule}{\hbox to\headwidth{\color{vinotinto}\leaders\hrule height \headrulewidth\hfill}}

% ---------- Titulos ----------
\titleformat{\section}{\normalfont\Large\bfseries\color{vinotinto}}{\thesection.}{0.6em}{}
\titleformat{\subsection}{\normalfont\large\bfseries\color{vinotinto}}{\thesubsection}{0.6em}{}
\titleformat{\subsubsection}{\normalfont\normalsize\bfseries}{\thesubsubsection}{0.6em}{}

\setlist[itemize]{leftmargin=1.2em,itemsep=2pt,topsep=3pt}
\setlist[enumerate]{leftmargin=1.4em,itemsep=2pt,topsep=3pt}

\captionsetup{font=small,labelfont={bf,color=vinotinto}}

% Columna de ancho fijo, alineada a la izquierda, sin justificar
\newcolumntype{L}[1]{>{\RaggedRight\arraybackslash}p{#1}}

\renewcommand{\arraystretch}{1.25}

% Pie de continuacion para tablas que se parten entre paginas
\newcommand{\contlinea}[1]{\multicolumn{#1}{r}{\footnotesize\itshape Continúa en la página siguiente}}

% Ficha compacta de escenario de calidad
\newenvironment{ficha}[3]%
{\par\needspace{3\baselineskip}\noindent{\bfseries\color{vinotinto}#1}\quad{\bfseries #2}\hfill{\footnotesize\color{grisTexto}#3}\par\vspace{0.08cm}%
\longtable{L{3.1cm} L{11.5cm}}\toprule}
{\bottomrule\endlongtable\addtocounter{table}{-1}\vspace{0.15cm}}

% Bandas del indice de prioridad
\newcommand{\bcritica}{\textcolor{vinotinto}{\textbf{Crítica}}}
\newcommand{\balta}{\textcolor{vinoclaro}{\textbf{Alta}}}
\newcommand{\bmedia}{\textcolor{ringProbar}{\textbf{Media}}}
\newcommand{\bbaja}{\textcolor{ringEvitar}{\textbf{Baja}}}

\begin{document}

% =====================================================================
%  PORTADA
% =====================================================================
\begin{titlepage}
\centering
\vspace*{2.0cm}

% Para incluir el isotipo, suba el archivo al proyecto de Overleaf
% y descomente la linea siguiente:
% \includegraphics[width=0.16\textwidth]{isotipo-redvital.png}\\[1.2cm]

{\Huge\bfseries\color{vinotinto} REDVITAL}\\[0.5cm]
{\large\color{vinotinto} Plataforma Nacional para la Gestión Integral del Ciclo de Vida\\ de la Donación de Sangre}\\[1.6cm]

{\color{vinotinto}\rule{0.45\textwidth}{1.2pt}}\\[1.6cm]

{\LARGE\bfseries DOCUMENTO DE ARQUITECTURA}\\[0.4cm]
{\LARGE\bfseries DE SOFTWARE --- SAD}\\[1.8cm]

{\large\bfseries Arkhé Software S.A.S.}\\[0.3cm]
{\itshape De la raíz a la red}\\[1.8cm]

\begin{tabular}{r l}
\textbf{Curso:} & Arquitectura de Software  \\
\textbf{Cliente:} & Profesor de la asignatura, en el rol simulado de \\
 & INS \\
\textbf{Entrega:} & Review Sprint 2 y Planning Sprint 3 --- Semana 8 \\
\textbf{Versión:} & 2.0 \\
\textbf{Estado:} & Aprobado para entrega \\
\textbf{Fecha:} & 22 de septiembre de 2026 \\
\textbf{Clasificación:} & Uso interno del proyecto \\
\end{tabular}

\vfill
{\small\color{grisTexto} Pontificia Universidad Javeriana --- Bogotá, Colombia}
\end{titlepage}

% =====================================================================
%  CONTROL DE VERSIONES
% =====================================================================
\section*{Control de versiones}
\addcontentsline{toc}{section}{Control de versiones}

\begin{longtable}{L{1.5cm} L{1.9cm} L{2.4cm} L{8.0cm}}
\toprule
\textbf{Versión} & \textbf{Fecha} & \textbf{Autor} & \textbf{Descripción del cambio} \\
\midrule
\endfirsthead
\toprule
\textbf{Versión} & \textbf{Fecha} & \textbf{Autor} & \textbf{Descripción del cambio} \\
\midrule
\endhead
\midrule
\contlinea{4} \\
\endfoot
\bottomrule
\endlastfoot
1.0 & 01/09/2026 & Equipo Arkhé & Versión inicial. Identificación de \emph{drivers} y \emph{killers} arquitectónicos, priorización de atributos de calidad conforme a ISO/IEC 25010:2023 con cobertura de las cuarenta subcaracterísticas de la norma, índice de prioridad cuantitativo con regla de desempate, catálogo de cuarenta y dos escenarios de calidad con medida de respuesta, registro de compromisos y decisiones arquitectónicas aprobadas (ADR-003 a ADR-006), arquitectura de alto nivel, arquitectura de negocio, arquitectura de datos y arquitectura de infraestructura. Matriz de trazabilidad entre requerimientos, escenarios, tácticas y componentes. \\
2.0 & 22/09/2026 & Equipo Arkhé & Incorporación de ADR-007 a ADR-016. La arquitectura se realiza con cinco servicios de grano grueso ---Identidad, Institucional, Campañas, Donación y Notificaciones--- conforme a ADR-014, que reemplaza a ADR-004, y con una topología de cinco zonas de confianza, dieciocho aristas permitidas y un proxy de borde con Caddy (ADR-015). La emisión de identidad pasa al Servicio de Identidad, con firma asimétrica, token de acceso de quince minutos y sesión de renovación de ocho horas (ADR-008). La bitácora de auditoría se distribuye en cuatro series con consulta compuesta en el Servicio Institucional (ADR-016, que reemplaza a ADR-010). El escalamiento en dos niveles pasa a ser un proceso del Servicio de Donación con alcance departamental (ADR-013). El límite de tasa se desdobla en escala anónima por origen y escala autenticada por sesión, ambas en el gateway (ADR-007). Los ambientes se fijan en Desarrollo, QA y Producción declarada. Los escenarios de calidad incorporan los hallazgos de la revisión cruzada del catálogo: EC-11 pasa de impacto 2 a impacto 3 porque su fallo retrasa la convocatoria ante una escasez; se cierran las medidas de EC-02, EC-06, EC-12, EC-13, EC-39 y EC-40; se define el conjunto cerrado de acciones irreversibles que citan EC-19 y EC-41; y se declara dónde se ejecutan las mediciones de carga. Corrección del conteo de transiciones de la máquina de estados y de la matriz de trazabilidad, que incorpora los requerimientos funcionales que no figuraban en ella. Las reglas de negocio invariantes pasan al prefijo RNI. Cierre de las cuestiones abiertas P-01 a P-08. Alineación con el Documento de Infraestructura V1: roles separados por base con tablas de solo anexado impuestas en la base, imágenes \emph{chiseled} y \emph{distroless}, frontera F5 de observabilidad y generación de carga con Gatling fuera de la máquina de QA. \\
\end{longtable}

\vspace{0.3cm}

\section*{Control de cambios}
\addcontentsline{toc}{section}{Control de cambios}

\noindent Esta versión sustituye a la V1.0 del 1 de septiembre de 2026. La arquitectura descrita sustituye a la registrada en ADR-002, \emph{Monolito modular}, retirada por la restricción de diseño RD-01, y realiza con cinco servicios la arquitectura distribuida de ADR-003, en lugar de los dos servicios de ADR-004. Los doce Registros de Decisión Arquitectónica vigentes que sustentan la estructura descrita ---ADR-003, ADR-005 a ADR-009 y ADR-011 a ADR-016--- fueron deliberados y aprobados en Mesa de Arquitectura antes de la emisión de esta versión. ADR-004 queda reemplazado por ADR-014, y ADR-010 por ADR-016.

\vspace{0.4cm}

\section*{Aprobaciones}
\addcontentsline{toc}{section}{Aprobaciones}

\begin{longtable}{L{4.6cm} L{4.4cm} L{4.6cm}}
\toprule
\textbf{Rol} & \textbf{Responsable} & \textbf{Firma / Fecha} \\
\midrule
\endhead
\bottomrule
\endlastfoot
Arquitecta de Software & Sara & \\
Scrum Master / Project Manager & Carlos Santiago Pinzón & \\
Product Owner & Alexander Aponte & \\
QA / Tester & Tomás & \\
DevOps & Veli & \\
\end{longtable}

\vspace{0.2cm}
\noindent{\footnotesize La aprobación de la Arquitecta de Software es constitutiva: es la autora del contenido técnico. La del Product Owner acredita que los \emph{drivers} de negocio corresponden a lo acordado con el cliente. La de QA acredita que los cuarenta y dos escenarios de la Sección \ref{sec:escenarios} tienen medida de respuesta verificable. La de DevOps acredita que la arquitectura de infraestructura de la Sección \ref{sec:infra} es realizable con el instrumental del catálogo de herramientas.}

\newpage

% =====================================================================
%  INDICES
% =====================================================================
\tableofcontents
\newpage
\listoffigures
\vspace{0.8cm}
\listoftables
\newpage

% =====================================================================
%  1. INTRODUCCION
% =====================================================================
\section{Introducción}

\subsection{Propósito}

Este documento describe la arquitectura de RedVital: qué fuerzas la determinan, qué estructura se adoptó en respuesta, por qué se descartó lo que se descartó, y cómo se comprueba que la estructura elegida sostiene los atributos de calidad que el sistema necesita.

No es un documento de diseño. No especifica clases, esquemas de tabla con atributos ni firmas de operación; eso corresponde al Documento de Diseño de Datos (DD) y al Documento de Diseño de Software (SDD). Aquí se decide dónde están las fronteras, quién es dueño de qué, y qué no puede ocurrir nunca.

El criterio de inclusión es estricto: entra en este documento aquello cuya modificación posterior sería costosa. Lo que puede cambiarse en un sprint sin arrastrar consecuencias no es arquitectura y no se documenta aquí.

\subsection{Alcance}

\subsubsection*{Este documento cubre}

\begin{itemize}
\item Los \emph{drivers} arquitectónicos y los \emph{killers} del proyecto (Sección \ref{sec:drivers}).
\item La priorización de los atributos de calidad conforme a ISO/IEC 25010:2023 (Sección \ref{sec:atributos}).
\item Los cuarenta y dos escenarios de calidad con su medida de respuesta, su índice de prioridad y su sprint de verificación (Sección \ref{sec:escenarios}).
\item Los compromisos evaluados y las decisiones arquitectónicas registradas (Sección \ref{sec:tradeoff}).
\item La arquitectura de alto nivel: estilo, componentes, fronteras y tácticas (Sección \ref{sec:hld}).
\item La arquitectura de negocio: capacidades, procesos centrales y reglas invariantes (Sección \ref{sec:negocio}).
\item La arquitectura de datos: propiedad, clasificación, consistencia y retención (Sección \ref{sec:datos}).
\item La arquitectura de infraestructura: topología, ambientes, seguridad de borde y observabilidad (Sección \ref{sec:infra}).
\end{itemize}

\subsubsection*{Este documento no cubre}

\begin{itemize}
\item El modelo entidad--relación con atributos y los contratos OpenAPI, que corresponden al DD V2 de esta misma entrega.
\item Las vistas 4+1, el modelo C4 y el resultado de la prueba de concepto del gateway con su \emph{benchmarking}, que corresponden al SDD V1 de esta misma entrega.
\item La configuración concreta de los ambientes de Desarrollo, QA y Producción, la matriz de alcanzabilidad verificada sobre la composición real, el modelo de amenazas y el endurecimiento de contenedores, que corresponden al Documento de Infraestructura V1 de esta misma entrega.
\item El catálogo de herramientas. Este documento solo puede usar tecnología que ya figure en anillo \adoptar{} o \probar{} del Documento de Herramientas, Políticas y Lineamientos V3.0; no introduce ninguna.
\end{itemize}

\subsection{Audiencia y momento de uso}

Está dirigido al cliente, que evalúa si la estructura propuesta responde a las restricciones que impuso, y al equipo, que la construye. Se consulta en cuatro momentos: antes de abrir un componente nuevo, antes de mover una responsabilidad de un servicio a otro, antes de introducir una llamada entre servicios, y ante cualquier duda sobre qué dato puede almacenarse.

Fuera de esos cuatro momentos, la referencia diaria es el SRS.

\subsection{Relación con los demás artefactos}

\begin{longtable}{L{4.2cm} L{10.4cm}}
\caption{Relación de este documento con los demás artefactos del proyecto.}\label{tab:artefactos}\\
\toprule
\textbf{Artefacto} & \textbf{Qué aporta a este documento, y qué toma de él} \\
\midrule
\endfirsthead
\toprule
\textbf{Artefacto} & \textbf{Qué aporta a este documento, y qué toma de él} \\
\midrule
\endhead
\midrule
\contlinea{2} \\
\endfoot
\bottomrule
\endlastfoot
SRS V3.0 & Aporta los requerimientos funcionales RF-01 a RF-24, veintitrés de ellos vigentes; los no funcionales RNF-01 a RNF-11; los siete perfiles de usuario y las restricciones. Es la fuente de los \emph{drivers} de la Sección \ref{sec:drivers}. Este documento no redefine ningún requerimiento: cuando necesita uno, lo cita. \\
Herramientas, Políticas y Lineamientos V3.0 & Aporta el catálogo de tecnología permitida y las restricciones de diseño transcritas del cliente. Toma de este documento la justificación estructural de las entradas en anillo \adoptar{}. \\
ADR-003 a ADR-016 & Los doce vigentes, aprobados en Mesa de Arquitectura. Aportan la deliberación de cada decisión estructural. Son la fuente; este documento es la síntesis. Ante discrepancia, gana el ADR. \\
Escenarios de calidad EC-01 a EC-42 & Aportan la forma verificable de cada atributo de calidad. Se consolidan y priorizan en la Sección \ref{sec:escenarios}, que es su fuente; la hoja de cálculo del catálogo es una vista filtrable derivada de ella. \\
DD V2 & Toma de este documento el reparto de entidades en las cuatro bases y la política de consistencia; desarrolla el modelo detallado, los contratos y la matriz de autorización. \\
SDD V1 & Toma de este documento la estructura de componentes y la desarrolla en vistas 4+1 y modelo C4. Documenta la prueba de concepto del gateway y su \emph{benchmarking}. \\
Infraestructura V1 & Toma de este documento las zonas de confianza, las aristas permitidas y los ambientes, y los realiza en configuración verificable. \\
Backlog & Toma de este documento las historias técnicas derivadas de tácticas arquitectónicas, conforme a RD-07. \\
\end{longtable}

\subsection{Definiciones}

Los términos del dominio están definidos en el SRS V3.0, Sección 1.3, y no se repiten. Aquí se definen únicamente los términos propios del discurso arquitectónico empleados en este documento.

\begin{longtable}{L{3.4cm} L{11.2cm}}
\caption{Glosario de términos arquitectónicos empleados en este documento.}\label{tab:glosarioarq}\\
\toprule
\textbf{Término} & \textbf{Definición} \\
\midrule
\endfirsthead
\toprule
\textbf{Término} & \textbf{Definición} \\
\midrule
\endhead
\midrule
\contlinea{2} \\
\endfoot
\bottomrule
\endlastfoot
Driver arquitectónico & Fuerza que determina la estructura del sistema y cuya omisión obligaría a rehacerla. Puede ser un requerimiento funcional de alto impacto, un atributo de calidad o una restricción. \\
Killer arquitectónico & Condición que, de materializarse, invalida la arquitectura o impide entregar el sistema. No es un riesgo cualquiera: es aquel cuya realización no admite corrección dentro del horizonte del proyecto. \\
Táctica & Decisión de diseño que produce una respuesta concreta ante un estímulo, y que sirve a un atributo de calidad determinado. \\
Punto de sensibilidad & Elemento de la arquitectura cuya modificación altera de forma apreciable la respuesta de un atributo de calidad. \\
Punto de compromiso & Punto de sensibilidad que afecta a dos o más atributos de calidad en direcciones opuestas. Es donde se toma una decisión, no donde se optimiza. \\
Contexto acotado & Frontera dentro de la cual un modelo de dominio es coherente y cada término tiene un único significado. Determina dónde se corta un servicio. \\
Servicio propietario & Servicio desplegable dueño exclusivo de los datos de un módulo y única vía de acceso a ellos. \\
Grano grueso & Granularidad de servicio en la que un despliegue aloja varios módulos funcionales que comparten lenguaje de dominio, por oposición al servicio por módulo. \\
Escenario de calidad & Situación concreta descrita en seis elementos ---fuente, estímulo, contexto, artefacto, respuesta y medida de respuesta--- que hace verificable un atributo de calidad. \\
Medida de respuesta & Valor numérico o condición binaria que permite afirmar si un escenario se cumple. Un adjetivo no es una medida de respuesta. \\
Índice de prioridad & Valor entre 1,0 y 5,0 que ordena los escenarios de calidad entre sí. Se calcula a partir del impacto del incumplimiento, la exposición y el costo de corrección tardía, conforme a la Sección \ref{sec:indice}. \\
Escenario de carga & Escenario cuya fuente de estímulo es una cantidad declarada de trabajo concurrente o de peticiones por unidad de tiempo, y no la acción de un usuario individual. \\
Escenario de sobrecarga & Escenario de carga cuya cantidad excede deliberadamente la capacidad nominal declarada, para observar cómo degrada el sistema en lugar de si la soporta. \\
Acción irreversible & Operación cuyo efecto el sistema no admite deshacer. Son tres, y solo tres: la disposición final de una unidad, la decisión del banco cedente sobre una transferencia y el despacho de una unidad. EC-19 y EC-41 citan este conjunto. \\
Zona de confianza & Agrupación de componentes con el mismo nivel de exposición, aislada de las demás por una red de contenedores propia. La arquitectura define cinco: borde, presentación, aplicación, datos y observabilidad (ADR-014). \\
Bandeja de salida & Tabla de eventos notificables que el servicio emisor escribe en la misma transacción que el hecho que los origina, y que el trabajador de Notificaciones consume por un contrato interno. \\
Pila de medición & Instancia de la composición de QA, separada de la que se muestra al cliente, sobre la que se ejecutan las pruebas de carga y de sobrecarga. \\
\end{longtable}

\subsection{Referencias}

\begin{longtable}{L{1.2cm} L{13.4cm}}
\caption{Documentos y normas que este documento aplica o de los que depende.}\label{tab:referencias}\\
\toprule
\textbf{Ref.} & \textbf{Fuente} \\
\midrule
\endfirsthead
\toprule
\textbf{Ref.} & \textbf{Fuente} \\
\midrule
\endhead
\midrule
\contlinea{2} \\
\endfoot
\bottomrule
\endlastfoot
{[1]} & Documento de Requerimientos SRS V3.0. Usuarios, módulos, requerimientos funcionales y no funcionales, restricciones y trazabilidad. \\
{[2]} & Documento de Herramientas, Políticas y Lineamientos V3.0. Tech Radar, restricciones de diseño y políticas de construcción. \\
{[3]} & Documento de Proyecto V2. Formulación, objetivos, alcance, limitaciones, cronograma y presupuesto. \\
{[4]} & ADR-003. Arquitectura distribuida de servicios de grano grueso por contexto acotado. Aprobada. \\
{[5]} & ADR-004. Reparto de componentes entre Java y .NET. Reemplazada por ADR-014. \\
{[6]} & ADR-005. Base de datos relacional con esquema independiente por servicio. Aprobada. \\
{[7]} & ADR-006. Sistema de estilos y componentes de la aplicación web. Aprobada. \\
{[8]} & Catálogo de escenarios de calidad EC-01 a EC-42, en hoja de cálculo filtrable derivada de la Sección \ref{sec:escenarios} de este documento. \\
{[9]} & ISO/IEC 25010:2023. \emph{Systems and software engineering --- SQuaRE --- Product quality model}. \\
{[10]} & ISO/IEC/IEEE 42010:2022. \emph{Software, systems and enterprise --- Architecture description}. \\
{[11]} & Bass, L., Clements, P. y Kazman, R. \emph{Software Architecture in Practice}. Escenarios de atributos de calidad y tácticas. \\
{[12]} & Evans, E. \emph{Domain-Driven Design}. Contexto acotado como criterio de corte de servicio. \\
{[13]} & Congreso de la República de Colombia. Ley 1581 de 2012, protección de datos personales. \\
{[14]} & Presidencia de la República de Colombia. Decreto 1571 de 1993, funcionamiento de bancos de sangre. \\
{[15]} & Ministerio de Salud. Resolución 901 de 1996, modificada por la Resolución 1676 de 2023. \\
{[16]} & Instituto Nacional de Salud. Circular Externa 026 de 2017, SIHEVI-INS. Informe Nacional de Bancos de Sangre 2024. \\
{[17]} & W3C. \emph{Web Content Accessibility Guidelines} (WCAG) 2.2, nivel AA. \\
{[18]} & ADR-007. API Gateway con YARP. Aprobada. \\
{[19]} & ADR-008. Token de sesión con firma asimétrica. Aprobada. \\
{[20]} & ADR-009. Línea de plataforma Java 25. Aprobada. \\
{[21]} & ADR-010. Almacenamiento distribuido de la bitácora de auditoría. Reemplazada por ADR-016. \\
{[22]} & ADR-011. Matriz de autorización por grupo de operaciones. Aprobada. \\
{[23]} & ADR-012. Clasificación de la extensibilidad del dominio. Aprobada. \\
{[24]} & ADR-013. Orquestación del escalamiento en dos niveles. Aprobada. \\
{[25]} & ADR-014. Descomposición en cinco servicios y topología por zonas de confianza. Aprobada. \\
{[26]} & ADR-015. Caddy como proxy de borde. Aprobada. \\
{[27]} & ADR-016. Bitácora de auditoría distribuida en cuatro contextos. Aprobada. \\
\end{longtable}

% =====================================================================
%  2. DRIVERS Y KILLERS
% =====================================================================
\section{Drivers y killers arquitectónicos}
\label{sec:drivers}

\subsection{Método}

Un \emph{driver} arquitectónico es una fuerza cuya omisión obliga a rehacer la estructura del sistema. No todo requerimiento es un \emph{driver}: la mayoría de los veintitrés requerimientos funcionales vigentes del SRS se implementan dentro de una estructura ya definida y no la condicionan. Se admiten como \emph{drivers} únicamente los elementos que satisfacen al menos una de estas tres condiciones:

\begin{enumerate}
\item \textbf{Determina una frontera.} Su cumplimiento decide dónde se corta un componente o dónde reside un dato.
\item \textbf{Es irreversible dentro del horizonte del proyecto.} Corregirlo después cuesta más que un sprint completo.
\item \textbf{Su incumplimiento tiene consecuencia externa.} Sanción legal, daño a una persona o rechazo de la entrega.
\end{enumerate}

Los \emph{drivers} se agrupan en tres familias, siguiendo la práctica habitual: funcionales primarios, de atributo de calidad y restricciones. Las restricciones no se discuten ni se ponderan: son entradas del problema.

\subsection{Drivers funcionales primarios}

De los veintitrés requerimientos funcionales vigentes, seis condicionan la estructura. El resto se aloja en ella.\footnote{Los identificadores llegan hasta RF-24. RF-16 se retiró del catálogo en el SRS V2.0 por no describir una capacidad ejecutable, y su contenido se especifica en RNF-05 y en RI-01.}

\begin{longtable}{L{1.5cm} L{5.0cm} L{8.0cm}}
\caption{Requerimientos funcionales con impacto estructural.}\label{tab:driversfun}\\
\toprule
\textbf{ID} & \textbf{Requerimiento} & \textbf{Por qué condiciona la estructura} \\
\midrule
\endfirsthead
\toprule
\textbf{ID} & \textbf{Requerimiento} & \textbf{Por qué condiciona la estructura} \\
\midrule
\endhead
\midrule
\contlinea{3} \\
\endfoot
\bottomrule
\endlastfoot
DF-01 & RF-03, RF-04, RF-05, RF-06 --- Registro de donación, trazabilidad, marcado de unidad no apta y disposición final & Forman una sola operación conceptual que debe ser atómica: una unidad no puede quedar registrada a medias ni en dos estados a la vez. Exige que las cuatro capacidades residan en un mismo servicio y una misma base de datos transaccional. Es el argumento central del corte de contextos. \\
DF-02 & RF-09 --- Administración de la jerarquía territorial & La jurisdicción es el eje del control de acceso de todo el sistema. Viaja en el token como ruta territorial, de modo que ningún servicio consulta la jerarquía para autorizar una petición. La jerarquía la mantiene el Servicio Institucional; la asignación de jurisdicción a cada usuario, el Servicio de Identidad (ADR-008, ADR-014). \\
DF-03 & RF-07 --- Consulta de existencias dentro del alcance de datos del usuario & Toda consulta de inventario está filtrada por jurisdicción. El filtro no es un adorno de la consulta: es una condición que atraviesa las fronteras entre servicios y obliga a transportar la jurisdicción en el token. \\
DF-04 & RF-10 y RF-14 --- Transferencia entre bancos y movilización de donantes & Es el único proceso que necesita, para una sola decisión, datos de inventario (Servicio de Donación) y de territorio de instituciones (Servicio Institucional). Obliga a declarar qué servicio orquesta, con qué alcance de datos evalúa y con qué garantía de consistencia. Resuelto en ADR-013. \\
DF-05 & RF-18 --- Exclusión de unidades vencidas sin intervención manual & Introduce un actor que no es un usuario: el paso del tiempo. Exige un ejecutor programado dentro del servicio propietario del inventario y un actor \emph{sistema} en la bitácora. \\
DF-06 & RF-21 --- Consulta de bitácora en modo de solo lectura & El único usuario cuyo trabajo es la bitácora no tiene módulo funcional propio. Obliga a declarar la auditoría como servicio transversal con dueño explícito, y no como efecto secundario de los demás módulos. \\
\end{longtable}

\subsection{Drivers de atributo de calidad}

Los atributos de calidad son la fuerza dominante de esta arquitectura. La razón es explícita y viene del cliente: la seguridad es atributo de primer orden (RD-08) porque el sistema trata datos sensibles de salud.

\begin{longtable}{L{1.5cm} L{4.2cm} L{8.8cm}}
\caption{Atributos de calidad con impacto estructural y su origen.}\label{tab:driverscal}\\
\toprule
\textbf{ID} & \textbf{Driver} & \textbf{Consecuencia estructural} \\
\midrule
\endfirsthead
\toprule
\textbf{ID} & \textbf{Driver} & \textbf{Consecuencia estructural} \\
\midrule
\endhead
\midrule
\contlinea{3} \\
\endfoot
\bottomrule
\endlastfoot
DQ-01 & Confidencialidad de la causa clínica (RNF-01, EC-02) & El dato no se modela. Determina el modelo de dominio, el contrato de la API, el contenido de los registros de aplicación y las etiquetas de las métricas. Es el \emph{driver} de mayor alcance transversal del sistema. \\
DQ-02 & Control de acceso por jurisdicción (RNF-02, EC-01) & Exige un punto único de entrada que resuelva identidad y jurisdicción, y filtrado a nivel de registro dentro de cada servicio propietario. Dos controles, no uno. \\
DQ-03 & Responsabilidad y no repudio (RNF-09, EC-05) & Exige una bitácora de solo anexado, sin operación de modificación ni de borrado expuesta por la aplicación, y un catálogo explícito de acciones sensibles. \\
DQ-04 & Bloqueo del despacho de unidad no apta o vencida (EC-12) & Exige comportamiento restrictivo por defecto ante estado indeterminado, aplicado en el servicio y en toda ruta que consulte disponibilidad. Es el único escenario cuyo incumplimiento daña a una persona. \\
DQ-05 & Aislamiento del fallo de un servicio par (EC-08) & Exige plazo de espera acotado en toda llamada entre servicios y degradación explícita en la interfaz. Sin él, distribuir añade puntos de fallo sin contrapartida. \\
DQ-06 & Extensibilidad a otros tipos de donación (RNF-05, EC-24) & Exige que el tipo de material donado sea un valor de atributo y no una estructura propia. Es la decisión más barata de tomar ahora y la más cara de revertir después. \\
DQ-07 & Puesta en marcha con un comando (RNF-11, EC-23) & Exige que toda diferencia entre ambientes se exprese como variable de entorno y que la composición del sistema completo esté versionada. \\
DQ-08 & Evolución del contrato entre servicios (EC-22) & Exige contrato independiente de lenguaje, versionado semántico y consumidores tolerantes a campos desconocidos. \\
\end{longtable}

\subsection{Restricciones}

Las ocho restricciones de diseño del cliente (RD-01 a RD-08), las cuatro normativas (Ley 1581 de 2012, Decreto 1571 de 1993, Resolución 901 de 1996 e ISO/IEC 25010:2023) y las cuatro internas (RI-01 a RI-04) están transcritas en el SRS V3.0, Sección 7, y en el Documento de Herramientas V3.0, Sección 2. No se reproducen aquí para no crear una segunda copia que pueda divergir.

Se recuerda únicamente la \textbf{regla de precedencia} que gobierna los conflictos, porque este documento la aplica varias veces:

\begin{center}
\textit{restricción de negocio} $\;\rightarrow\;$ \textit{restricción de diseño} $\;\rightarrow\;$ \textit{condición del proyecto}
\end{center}

\noindent Una condición del proyecto ---presupuesto cero, capacidad del equipo, horizonte de cinco sprints--- nunca justifica romper una restricción de negocio. Esta regla es la que resuelve, por ejemplo, que no se adopte un proveedor de identidad completo por costo de aprendizaje pero sí se implemente la verificación doble de sesión: lo primero es instrumental, lo segundo es la restricción.

\subsection{Killers arquitectónicos}

Un \emph{killer} no es un riesgo más. Es una condición que, si se materializa, no admite corrección dentro del horizonte del proyecto: o invalida la arquitectura, o impide entregar. Se distinguen de los riesgos de gestión ---que viven en el Documento de Proyecto--- en que su consecuencia es estructural.

Cada uno declara la señal que permite detectarlo antes de que sea tarde. Un \emph{killer} sin señal de detección es una advertencia decorativa.

\begin{longtable}{L{1.3cm} L{3.6cm} L{4.4cm} L{4.6cm}}
\caption{Killers arquitectónicos, señal de detección y contramedida.}\label{tab:killers}\\
\toprule
\textbf{ID} & \textbf{Killer} & \textbf{Señal de detección} & \textbf{Contramedida adoptada} \\
\midrule
\endfirsthead
\toprule
\textbf{ID} & \textbf{Killer} & \textbf{Señal de detección} & \textbf{Contramedida adoptada} \\
\midrule
\endhead
\midrule
\contlinea{4} \\
\endfoot
\bottomrule
\endlastfoot
K-01 & \textbf{Exposición de la causa clínica.} Cualquier ruta ---campo, mensaje de error, registro de aplicación, etiqueta de métrica--- por la que se pueda inferir el diagnóstico de un donante. & Aparición de un campo de motivo, causa o resultado de tamizaje en una migración, un contrato o un texto de interfaz. & El dato no se captura. RI-02 lo prohíbe en el modelo. La revisión de cada Pull Request verifica migraciones, contrato y textos. Verificado por EC-02. \\
K-02 & \textbf{Contraprestación económica al donante.} Un reconocimiento con valor directo o indirecto convierte la donación en remunerada. & Un elemento del catálogo de reconocimientos con canje, saldo, descuento, sorteo o prioridad de atención; o un texto que prometa beneficios. & El catálogo de reconocimientos es simbólico por definición y se revisa en la Definición de Terminado de toda historia de M1. Verificado por EC-14. \\
K-03 & \textbf{Despacho de una unidad no apta o vencida.} Único fallo del sistema cuya consecuencia es física y sobre una persona. & Una unidad no apta o vencida presente en cualquier respuesta de disponibilidad, incluida la ruta de sugerencia de transferencia. & Comportamiento restrictivo por defecto ante estado indeterminado, aplicado en el servicio y en toda ruta de disponibilidad. Verificado por EC-12 y EC-13. \\
K-04 & \textbf{Acoplamiento por la capa de persistencia.} Dos servicios que comparten base de datos dejan de ser distribuidos y RD-01 se incumple de hecho aunque haya dos despliegues. & Una cadena de conexión que apunte a la base de otro servicio, o una consulta que cruce dos esquemas. & Una base por servicio con base (ADR-005, ADR-014), credenciales distintas por base, y verificación de integración continua sobre las cadenas de conexión. Verificado por EC-22. \\
K-05 & \textbf{Fallo en cascada entre servicios.} Una llamada sin plazo de espera agota el pool de conexiones y convierte la caída de un servicio en la caída de los que lo consultan. & Ausencia de tiempo de espera en el cliente HTTP de cualquier servicio. & Tiempo de espera máximo de 3 segundos por llamada entre servicios, y degradación declarada en la interfaz. Verificado por EC-08. \\
K-06 & \textbf{Modelo de dominio acoplado a la sangre.} Si «sangre» aparece en el nombre de las tablas de trazabilidad e inventario, la extensibilidad se pierde y ya no se recupera barato. & Nombres de tabla, entidad o columna que nombren el material donado en lugar de tratarlo como valor de un atributo de tipo. & El tipo de donación es un valor de catálogo. Se revisa en la primera migración del Sprint 2, antes de que existan datos. Verificado por EC-24. \\
K-07 & \textbf{Autorización solo en la interfaz.} Ocultar una opción del menú no es un control cumplido; el atacante llama a la API. & Un endpoint que responda a un usuario fuera de jurisdicción cuando se le invoca sin pasar por la aplicación web. & Doble verificación: compuerta en el API Gateway y filtrado por jurisdicción en cada servicio propietario. Se prueba invocando el servicio sin pasar por el gateway. Verificado por EC-01 y EC-03. \\
K-08 & \textbf{Sobrecarga del Sprint 3.} Concentrar en un mismo sprint el gateway, la autorización territorial, el primer flujo funcional completo y la conexión del frontend produce cuatro cosas a medias y ninguna demostrable. & Compromiso del Sprint 3 por encima de la capacidad medida, o historias sin criterio de aceptación verificable. & Recorte por categoría MoSCoW completa, nunca reducción uniforme del alcance de cada historia. El gateway y la emisión de identidad no son recortables: sostienen K-07. \\
K-09 & \textbf{Decisión estructural sin ADR.} Una decisión tomada dentro del código no se puede discutir, revisar ni revertir de forma trazable. & Tecnología en ejecución que no figure en el catálogo de herramientas, o componente cuya elección no tenga ADR asociado. & Toda decisión estructural vigente tiene ADR aprobado (Tabla \ref{tab:adrs}), y las cuestiones de arquitectura de la V1.0 quedan cerradas en el Anexo \ref{sec:anexoB}. La entrada de una tecnología nueva pasa por la Mesa de Arquitectura, no por un Pull Request. \\
\end{longtable}

\subsection{Priorización de los drivers}

Los \emph{drivers} se ordenan por dos ejes: importancia para el cliente y dificultad técnica para el equipo. El cuadrante superior derecho ---alta importancia y alta dificultad--- es el que consume la atención de la arquitectura y el que justifica que los Sprints 1 y 2 sean fundacionales.

\begin{figure}[h!]
\centering
\begin{tikzpicture}[scale=0.95, every node/.style={font=\scriptsize}]
% Ejes
\draw[->, draw=grisTexto] (0,0) -- (11.6,0) node[anchor=north east, text=grisTexto] {Dificultad técnica};
\draw[->, draw=grisTexto] (0,0) -- (0,7.4) node[anchor=north west, rotate=90, text=grisTexto] {Importancia para el cliente};
\draw[dashed, draw=grisTexto!50] (5.8,0) -- (5.8,7.2);
\draw[dashed, draw=grisTexto!50] (0,3.7) -- (11.4,3.7);

% Etiquetas de cuadrante
\node[text=grisTexto!70, anchor=north west] at (0.15,7.25) {Resolver pronto y barato};
\node[text=vinotinto!75, anchor=north east] at (11.3,7.25) {Foco de la arquitectura};
\node[text=grisTexto!70, anchor=south west] at (0.15,0.15) {Rutina};
\node[text=grisTexto!70, anchor=south east] at (11.3,0.15) {Diferir o acotar};

% Puntos: cuadrante superior derecho
\fill[vinotinto] (7.6,6.6) circle (2.6pt); \node[anchor=west, text=vinotinto] at (7.75,6.6) {DQ-02 jurisdicción};
\fill[vinotinto] (9.2,6.0) circle (2.6pt); \node[anchor=east, text=vinotinto] at (9.05,6.25) {DQ-03 no repudio};
\fill[vinotinto] (6.6,5.4) circle (2.6pt); \node[anchor=west, text=vinotinto] at (6.75,5.4) {DF-04 transferencias};
\fill[vinotinto] (8.4,4.6) circle (2.6pt); \node[anchor=west, text=vinotinto] at (8.55,4.6) {DQ-05 aislamiento};
\fill[vinotinto] (6.3,4.1) circle (2.6pt); \node[anchor=west, text=vinotinto] at (6.45,4.1) {DQ-08 contrato};

% Superior izquierdo
\fill[vinoclaro] (1.5,6.9) circle (2.6pt); \node[anchor=west, text=vinoclaro] at (1.65,6.9) {DQ-01 confidencialidad};
\fill[vinoclaro] (2.4,6.2) circle (2.6pt); \node[anchor=west, text=vinoclaro] at (2.55,6.2) {DQ-04 no apta};
\fill[vinoclaro] (1.9,5.4) circle (2.6pt); \node[anchor=west, text=vinoclaro] at (2.05,5.4) {DQ-06 extensibilidad};
\fill[vinoclaro] (3.4,4.6) circle (2.6pt); \node[anchor=west, text=vinoclaro] at (3.55,4.6) {DF-01 atomicidad};
\fill[vinoclaro] (2.7,4.0) circle (2.6pt); \node[anchor=west, text=vinoclaro] at (2.85,4.0) {DF-05 vencimiento};

% Inferior derecho
\fill[ringEvaluar] (7.9,2.6) circle (2.6pt); \node[anchor=west, text=ringEvaluar] at (8.05,2.6) {DQ-07 puesta en marcha};
\fill[ringEvaluar] (6.9,1.6) circle (2.6pt); \node[anchor=west, text=ringEvaluar] at (7.05,1.6) {DF-06 bitácora};

% Inferior izquierdo
\fill[ringEvitar] (2.2,2.4) circle (2.6pt); \node[anchor=west, text=ringEvitar] at (2.35,2.4) {DF-02 jerarquía};
\fill[ringEvitar] (3.0,1.4) circle (2.6pt); \node[anchor=west, text=ringEvitar] at (3.15,1.4) {DF-03 existencias};
\end{tikzpicture}
\caption[Priorización de los drivers arquitectónicos]{Priorización de los \emph{drivers} arquitectónicos por importancia y dificultad. Los del cuadrante superior derecho concentran el esfuerzo de diseño y son los que no se recortan; los del superior izquierdo son igual de importantes pero se resuelven con una decisión de modelado tomada a tiempo, y por eso el Sprint 1 es el momento de tomarlas.}
\label{fig:cuadrantedrivers}
\end{figure}

\noindent Dos lecturas de la Figura \ref{fig:cuadrantedrivers} merecen declararse:

\begin{itemize}
\item \textbf{DQ-01 y DQ-06 son de alta importancia y baja dificultad, pero solo ahora.} No modelar la causa clínica y no acoplar el dominio a la sangre cuestan cero si se deciden antes de la primera migración. Después de que existan datos, ambas se convierten en migraciones de esquema con pérdida de trazabilidad. Su dificultad es baja únicamente en el Sprint 1, y por eso figuran en el cuadrante izquierdo con la advertencia de que no permanecerán allí.
\item \textbf{DF-04 es el \emph{driver} funcional más caro.} Es el único proceso que necesita datos de dos contextos para una sola decisión, y obliga a una decisión de orquestación que ADR-013 resuelve como proceso programado del Servicio de Donación (Sección \ref{sec:consistencia}).
\end{itemize}

% =====================================================================
%  3. ATRIBUTOS DE CALIDAD
% =====================================================================
\section{Atributos de calidad}
\label{sec:atributos}

\subsection{Marco de referencia}

Los atributos de calidad de este sistema se clasifican conforme a \textbf{ISO/IEC 25010:2023}, en su revisión de nueve características y cuarenta subcaracterísticas. La versión de la norma se declara aquí de forma expresa y \textbf{rige todo el documento}: no se mezclan categorías de la revisión de 2011 con las de 2023, y ningún escenario se clasifica en una subcaracterística que la revisión vigente no contemple.

La elección no es un detalle de cita. La revisión de 2023 introduce dos características que en este dominio no son ornamentales:

\begin{itemize}
\item \textbf{Capacidad de interacción} sustituye a \emph{usabilidad} y desagrega ocho subcaracterísticas, entre ellas \emph{inclusividad}, que es la que sustenta el criterio de accesibilidad de un sistema público de salud.
\item \textbf{Seguridad física} es característica nueva. Agrupa las cinco subcaracterísticas cuyo incumplimiento tiene consecuencia sobre la integridad de una persona, y en RedVital corresponde al único fallo que el sistema existe para impedir: que una unidad no apta llegue a transfusión.
\end{itemize}

\subsection{Cobertura de la norma}

Las nueve características y las cuarenta subcaracterísticas tienen al menos un escenario de calidad asociado. La cobertura completa es un criterio de este documento y no una aspiración: una subcaracterística sin escenario es un atributo enunciado como adjetivo, y un adjetivo no se puede verificar ni discutir.

\begin{longtable}{L{3.4cm} L{4.4cm} L{2.0cm} L{4.4cm}}
\caption{Cobertura de las nueve características y las cuarenta subcaracterísticas de ISO/IEC 25010:2023.}\label{tab:coberturacal}\\
\toprule
\textbf{Característica} & \textbf{Subcaracterística} & \textbf{Escenario} & \textbf{Denominación del escenario} \\
\midrule
\endfirsthead
\toprule
\textbf{Característica} & \textbf{Subcaracterística} & \textbf{Escenario} & \textbf{Denominación del escenario} \\
\midrule
\endhead
\midrule
\contlinea{4} \\
\endfoot
\bottomrule
\endlastfoot
Adecuación funcional & Completitud funcional & EC-25 & Cobertura del ciclo completo de la unidad \\
 & Corrección funcional & EC-13 & Exclusión automática de unidades vencidas \\
 & Pertinencia funcional & EC-14 & Reconocimiento sin contraprestación económica \\
\midrule
Eficiencia de desempeño & Comportamiento temporal & EC-15 & Tiempo de respuesta de la consulta de inventario \\
 & Utilización de recursos & EC-26 & Consumo bajo carga nominal sostenida \\
 & Capacidad & EC-16 & Concurrencia durante el pico de una campaña \\
\midrule
Compatibilidad & Coexistencia & EC-27 & Convivencia de los contenedores del sistema en un mismo anfitrión \\
 & Interoperabilidad & EC-22 & Evolución del contrato entre servicios \\
\midrule
Capacidad de interacción & Reconocibilidad de la adecuación & EC-28 & Identificación del propósito por un donante nuevo \\
 & Aprendizaje & EC-29 & Primer registro de donación por un operador nuevo \\
 & Operabilidad & EC-17 & Registro anónimo en tres pasos \\
 & Protección contra errores de usuario & EC-19 & Protección ante el desecho accidental de una unidad \\
 & Compromiso del usuario & EC-30 & Retorno del donante tras recuperar la elegibilidad \\
 & Inclusividad & EC-18 & Uso con lector de pantalla y sin ratón \\
 & Asistencia al usuario & EC-31 & Recuperación ante un rechazo de operación \\
 & Autodescripción & EC-32 & Legibilidad del estado de una unidad sin formación previa \\
\midrule
Fiabilidad & Ausencia de fallos & EC-10 & Registro de donación idempotente \\
 & Disponibilidad & EC-07 & Disponibilidad de la consulta de inventario \\
 & Tolerancia a fallos & EC-08 & Aislamiento del fallo del servicio par \\
 & Recuperabilidad & EC-09 & Recuperación tras pérdida de datos \\
 & Recuperabilidad & EC-11 & Reanudación de notificaciones acumuladas \\
\midrule
Seguridad & Confidencialidad & EC-01 & Control de acceso por jurisdicción \\
 & Confidencialidad & EC-02 & Confidencialidad de la causa clínica \\
 & Integridad & EC-04 & Integridad del ciclo de vida de la unidad \\
 & No repudio & EC-05 & No repudio de la bitácora de auditoría \\
 & Responsabilidad & EC-33 & Atribución de la acción automática al sistema \\
 & Autenticidad & EC-03 & Autenticidad de la sesión \\
 & Resistencia & EC-06 & Resistencia ante abuso automatizado \\
\midrule
Mantenibilidad & Modularidad & EC-34 & Aislamiento del cambio dentro de un contexto \\
 & Reusabilidad & EC-35 & Reutilización del componente de interfaz \\
 & Analizabilidad & EC-20 & Diagnóstico de un fallo en el ambiente de QA \\
 & Modificabilidad & EC-21 & Cambio de la regla de elegibilidad \\
 & Capacidad de ser probado & EC-36 & Verificación exhaustiva de la máquina de estados \\
\midrule
Flexibilidad & Adaptabilidad & EC-24 & Extensión a otro tipo de donación \\
 & Escalabilidad & EC-37 & Sobrecarga de cuatro veces la carga nominal \\
 & Instalabilidad & EC-23 & Puesta en marcha del sistema completo \\
 & Reemplazabilidad & EC-38 & Sustitución del API Gateway \\
\midrule
Seguridad física & Restricción operativa & EC-39 & Límite operativo del despacho por jurisdicción \\
 & Identificación de riesgos & EC-40 & Detección de unidades próximas al vencimiento \\
 & A prueba de fallos & EC-12 & Bloqueo del despacho de una unidad no apta o vencida \\
 & Advertencia de peligros & EC-41 & Advertencia previa a una acción irreversible \\
 & Integración segura & EC-42 & Incorporación de una institución nueva a la red \\
\end{longtable}

\begin{longtable}{L{4.2cm} L{2.6cm} L{2.6cm} L{4.6cm}}
\caption{Distribución de los escenarios por característica.}\label{tab:distribucion}\\
\toprule
\textbf{Característica} & \textbf{Subcaract.} & \textbf{Escenarios} & \textbf{Peso relativo del catálogo} \\
\midrule
\endfirsthead
\toprule
\textbf{Característica} & \textbf{Subcaract.} & \textbf{Escenarios} & \textbf{Peso relativo del catálogo} \\
\midrule
\endhead
\midrule
\contlinea{4} \\
\endfoot
\bottomrule
\endlastfoot
Capacidad de interacción & 8 & 8 & 19,0\,\% \\
Seguridad & 6 & 7 & 16,7\,\% \\
Fiabilidad & 4 & 5 & 11,9\,\% \\
Seguridad física & 5 & 5 & 11,9\,\% \\
Mantenibilidad & 5 & 5 & 11,9\,\% \\
Flexibilidad & 4 & 4 & 9,5\,\% \\
Adecuación funcional & 3 & 3 & 7,1\,\% \\
Eficiencia de desempeño & 3 & 3 & 7,1\,\% \\
Compatibilidad & 2 & 2 & 4,8\,\% \\
\midrule
\textbf{Total} & \textbf{40} & \textbf{42} & \textbf{100\,\%} \\
\end{longtable}

\noindent Dos subcaracterísticas tienen dos escenarios en lugar de uno. La confidencialidad los tiene porque las dos formas de fuga que este sistema debe impedir son de naturaleza distinta: una es de alcance ---ver datos de otra jurisdicción--- y la otra es de contenido ---ver la causa clínica de un rechazo---, y una sola medida de respuesta no puede cubrirlas. La recuperabilidad los tiene porque recuperar una base de datos y recuperar una cola de notificaciones pendientes son procedimientos distintos con objetivos de tiempo distintos.

El peso del catálogo no coincide con el peso de la arquitectura. La capacidad de interacción concentra el mayor número de escenarios por la sola razón de que la norma le asigna ocho subcaracterísticas, mientras que la seguridad y la seguridad física concentran los índices de prioridad más altos. El número de escenarios mide cobertura de la norma; el índice de prioridad de la sección siguiente mide importancia.

\subsection{Índice de prioridad}
\label{sec:indice}

Clasificar los escenarios en alta, media y baja produce empates masivos: en un catálogo de cuarenta y dos elementos, más de la mitad cae en la misma banda y la escala deja de ordenar. Cuando la capacidad del sprint obliga a decidir cuál de dos escenarios se verifica primero, una etiqueta compartida no resuelve nada.

Se adopta por ello un \textbf{índice de prioridad} numérico, calculado a partir de tres factores independientes, cada uno en escala de 1 a 5.

\begin{longtable}{L{1.0cm} L{3.2cm} L{10.4cm}}
\caption{Factores del índice de prioridad y significado de cada nivel.}\label{tab:factores}\\
\toprule
\textbf{F.} & \textbf{Factor} & \textbf{Escala} \\
\midrule
\endfirsthead
\toprule
\textbf{F.} & \textbf{Factor} & \textbf{Escala} \\
\midrule
\endhead
\midrule
\contlinea{3} \\
\endfoot
\bottomrule
\endlastfoot
\textbf{I} & Impacto del incumplimiento & \textbf{5} daño a la integridad física de una persona o incumplimiento normativo sancionable. \textbf{4} pérdida de confidencialidad o de integridad de datos sensibles. \textbf{3} pérdida de una capacidad esencial del servicio. \textbf{2} degradación perceptible sin pérdida de capacidad. \textbf{1} molestia o pérdida de comodidad. \\
\textbf{E} & Exposición & Frecuencia con que el estímulo se presenta. \textbf{5} en cada petición del sistema. \textbf{4} varias veces al día. \textbf{3} semanal. \textbf{2} mensual. \textbf{1} excepcional o de una sola vez. \\
\textbf{C} & Costo de corrección tardía & \textbf{5} obliga a rehacer el modelo de datos o la frontera entre servicios. \textbf{4} obliga a cambiar el contrato entre servicios. \textbf{3} varias historias de trabajo. \textbf{2} una historia. \textbf{1} un cambio de configuración. \\
\end{longtable}

\noindent El índice se calcula como media ponderada y se expresa con un decimal:

\begin{center}
\textbf{IP} $= 0{,}50 \cdot I + 0{,}30 \cdot E + 0{,}20 \cdot C$ \qquad con $\text{IP} \in [1{,}0;\ 5{,}0]$
\end{center}

Los pesos expresan una posición explícita del equipo. El impacto pesa la mitad del índice porque la consecuencia de un fallo es lo que el cliente y la norma juzgan. La exposición pesa menos porque un fallo raro y grave sigue siendo grave. El costo de corrección tardía pesa lo menos, pero pesa: es el factor que impide posponer indefinidamente una decisión barata hoy y cara en el Sprint 4.

\begin{longtable}{L{2.4cm} L{2.4cm} L{2.0cm} L{7.4cm}}
\caption{Bandas del índice de prioridad y consecuencia sobre la planeación.}\label{tab:bandas}\\
\toprule
\textbf{Banda} & \textbf{Rango} & \textbf{Escenarios} & \textbf{Consecuencia} \\
\midrule
\endhead
\bottomrule
\endlastfoot
\bcritica{} & 4,5 --- 5,0 & 2 & No se recorta ni se pospone bajo ninguna circunstancia. Su incumplimiento detiene la entrega. \\
\balta{} & 3,5 --- 4,4 & 13 & Se verifica en el sprint en que se vuelve verificable. Solo se aplaza con acta de la Mesa de Arquitectura. \\
\bmedia{} & 2,5 --- 3,4 & 18 & Se verifica en el sprint previsto; admite aplazamiento de un sprint si la capacidad obliga. \\
\bbaja{} & 1,0 --- 2,4 & 9 & Se verifica cuando la capacidad lo permita, sin comprometer fecha. \\
\end{longtable}

\subsubsection*{Regla de desempate}

Dos escenarios pueden compartir índice. La regla siguiente produce un \textbf{orden total}, de modo que la pregunta «cuál de estos dos primero» siempre tiene una respuesta y no queda a criterio de quien planea el sprint. Se aplica en orden y se detiene en el primer criterio que discrimine:

\begin{enumerate}
\item \textbf{Mayor impacto (I).} Entre dos escenarios de igual índice, gana el que hace más daño al incumplirse.
\item \textbf{Mayor costo de corrección tardía (C).} A igual impacto, gana el que resulta más caro de arreglar después.
\item \textbf{Mayor exposición (E).} A igual impacto y costo, gana el que ocurre más a menudo.
\item \textbf{Menor identificador de escenario.} Criterio de cierre: garantiza que el orden sea determinista y reproducible.
\end{enumerate}

\noindent La regla no es teórica en este catálogo. EC-10, EC-33 y EC-39 comparten índice 4,0: el primer criterio adelanta a EC-39, cuyo impacto es 5 porque su incumplimiento alcanza a una persona; el segundo y el tercero no discriminan entre EC-10 y EC-33, y es el cuarto el que ordena. Lo mismo ocurre entre EC-19 y EC-41, ambos en 3,8 con idénticos factores.

\subsection{Regla de precedencia entre atributos}

El índice ordena el trabajo. La precedencia resuelve el conflicto de diseño, que es un problema distinto: qué se cede cuando dos atributos no pueden satisfacerse a la vez.

\begin{enumerate}
\item \textbf{Seguridad física.} Ninguna consideración de facilidad de uso o de desempeño justifica despachar una unidad no apta.
\item \textbf{Seguridad y cumplimiento legal.} Son restricciones normativas, no objetivos negociables.
\item \textbf{Fiabilidad.}
\item \textbf{El resto}, ponderado por índice de prioridad. Si la ponderación tiene consecuencia estructural, se registra en un ADR.
\end{enumerate}

\noindent El conflicto más frecuente en este sistema es \textbf{seguridad contra facilidad de uso}: EC-17 exige un registro anónimo de tres pasos y EC-06 exige contener el abuso automatizado del mismo punto de entrada. La resolución adoptada no añade un paso de verificación al registro ---eso rompería EC-17, cuyo límite de tres pasos es duro--- sino que limita la tasa por origen en el API Gateway. La tensión se resuelve moviendo el control de componente, no cediendo en el atributo.

% =====================================================================
%  4. ESCENARIOS DE CALIDAD
% =====================================================================
\section{Escenarios de calidad}
\label{sec:escenarios}

\subsection{Estructura del escenario}

El escenario convierte un atributo en una situación concreta descrita en seis elementos, con una \textbf{medida de respuesta que es un número o una condición binaria}.

\begin{longtable}{L{3.6cm} L{11.0cm}}
\caption{Los seis elementos del escenario de atributo de calidad.}\label{tab:elementosec}\\
\toprule
\textbf{Elemento} & \textbf{Qué declara} \\
\midrule
\endfirsthead
\toprule
\textbf{Elemento} & \textbf{Qué declara} \\
\midrule
\endhead
\midrule
\contlinea{2} \\
\endfoot
\bottomrule
\endlastfoot
Fuente del estímulo & Quién o qué origina la situación. Puede no ser una persona: en EC-13 la fuente es el paso del tiempo, y en los escenarios de carga es una cantidad declarada de trabajo concurrente. \\
Estímulo & El suceso concreto al que el sistema debe responder. \\
Contexto & El estado del sistema cuando ocurre. Un mismo estímulo tiene respuestas distintas en operación normal, degradada o de sobrecarga. \\
Artefacto afectado & Qué parte del sistema responde. Es lo que ata el escenario a la arquitectura. \\
Respuesta & Qué hace el sistema, descrito de forma observable. \\
Medida de respuesta & El número o la condición binaria que permite afirmar si se cumplió. \\
\end{longtable}

\subsection{Escenarios de carga y de sobrecarga}

Seis escenarios del catálogo declaran como fuente de estímulo una \textbf{cantidad de carga}, y no la acción de un usuario individual. Se agrupan aquí porque comparten método de verificación, volumen de referencia y ambiente de medición, y porque su interpretación depende de una cifra que debe estar declarada antes de medir.

\begin{longtable}{L{1.4cm} L{4.2cm} L{4.2cm} L{4.6cm}}
\caption{Escenarios cuya fuente de estímulo es una carga declarada.}\label{tab:carga}\\
\toprule
\textbf{ID} & \textbf{Escenario} & \textbf{Carga declarada} & \textbf{Naturaleza} \\
\midrule
\endfirsthead
\toprule
\textbf{ID} & \textbf{Escenario} & \textbf{Carga declarada} & \textbf{Naturaleza} \\
\midrule
\endhead
\midrule
\contlinea{4} \\
\endfoot
\bottomrule
\endlastfoot
EC-06 & Resistencia ante abuso automatizado & 600 peticiones por minuto desde un mismo origen, diez veces la tasa nominal por origen & \textbf{Sobrecarga} dirigida a un punto de entrada público \\
EC-07 & Disponibilidad de la consulta de inventario & 20 usuarios concurrentes durante la ventana de evaluación & Carga nominal con fallo inducido \\
EC-15 & Tiempo de respuesta de la consulta de inventario & 20 usuarios concurrentes sobre el volumen sintético de referencia & Carga nominal \\
EC-16 & Concurrencia durante el pico de una campaña & 50 usuarios concurrentes durante 10 minutos & Carga de pico declarada \\
EC-26 & Consumo bajo carga nominal sostenida & 50 usuarios concurrentes durante 30 minutos & Carga nominal sostenida \\
EC-37 & Sobrecarga de cuatro veces la carga nominal & 200 usuarios concurrentes durante 5 minutos & \textbf{Sobrecarga} deliberada del sistema completo \\
\end{longtable}

\noindent\textbf{Volumen de referencia.} Toda medición se ejecuta sobre el conjunto sintético declarado y versionado en el repositorio: 5.000 donantes, 12.000 donaciones, 30.000 unidades y aproximadamente 95.000 eventos de trazabilidad. Una medida de desempeño sin volumen declarado no es comparable entre sprints y permitiría presentar como mejora lo que solo es un conjunto de datos más pequeño.

\textbf{Escala de la demostración.} Las cifras de 50 y 200 usuarios concurrentes corresponden a la escala de demostración del proyecto y no a la del sistema nacional real. Se declaran como tales ante el cliente y se ajustan si el cliente fija otra.

\textbf{Dónde se mide.} Las seis mediciones se ejecutan sobre la pila de medición: la misma composición de QA, levantada con otro nombre de proyecto sobre la máquina virtual institucional, con la instancia que se muestra al cliente detenida y fuera de las ventanas de demostración. La carga la genera Gatling desde el equipo de un integrante, fuera de la máquina virtual. La configuración de esa pila declara el origen de medición exento de la escala anónima del límite de tasa (ADR-007), sin el cual una prueba de cincuenta usuarios concurrentes se rechazaría a sí misma. EC-16 y EC-26 comparten carga y ambiente y se ejecutan como una sola campaña de medición: los diez primeros minutos de la carga sostenida de EC-26 son la medición de EC-16.

\textbf{Distinción entre carga y sobrecarga.} EC-16 y EC-26 comprueban que el sistema \emph{soporta} una carga declarada. EC-06 y EC-37 comprueban algo distinto y más informativo: cómo \emph{degrada} cuando la carga excede su capacidad. Un sistema que solo se prueba dentro de su capacidad nominal no ha demostrado nada sobre su comportamiento el día que la exceda, y ese día su modo de fallo es lo único que importa.

\subsection{Catálogo}

Los escenarios se presentan ordenados por índice de prioridad descendente, con la regla de desempate aplicada. El orden de esta tabla es el orden en que se atienden.

\begin{longtable}{L{1.2cm} L{4.4cm} L{2.8cm} L{1.7cm} L{0.9cm} L{1.6cm}}
\caption{Catálogo de los cuarenta y dos escenarios de calidad, ordenado por índice de prioridad.}\label{tab:catalogoec}\\
\toprule
\textbf{ID} & \textbf{Escenario} & \textbf{Subcaract.} & \textbf{I\,·\,E\,·\,C} & \textbf{IP} & \textbf{Banda} \\
\midrule
\endfirsthead
\toprule
\textbf{ID} & \textbf{Escenario} & \textbf{Subcaract.} & \textbf{I\,·\,E\,·\,C} & \textbf{IP} & \textbf{Banda} \\
\midrule
\endhead
\midrule
\contlinea{6} \\
\endfoot
\bottomrule
\endlastfoot
EC-02 & Confidencialidad de la causa clínica & Confidencialidad & 5·4·5 & \textbf{4,7} & \bcritica{} \\
EC-12 & Bloqueo del despacho de una unidad no apta o vencida & A prueba de fallos & 5·4·4 & \textbf{4,5} & \bcritica{} \\
EC-01 & Control de acceso por jurisdicción & Confidencialidad & 4·5·4 & 4,3 & \balta{} \\
EC-03 & Autenticidad de la sesión & Autenticidad & 4·5·4 & 4,3 & \balta{} \\
EC-04 & Integridad del ciclo de vida de la unidad & Integridad & 4·4·5 & 4,2 & \balta{} \\
EC-13 & Exclusión automática de unidades vencidas & Corrección funcional & 4·5·3 & 4,1 & \balta{} \\
EC-39 & Límite operativo del despacho por jurisdicción & Restricción operativa & 5·3·3 & 4,0 & \balta{} \\
EC-10 & Registro de donación idempotente & Ausencia de fallos & 4·4·4 & 4,0 & \balta{} \\
EC-33 & Atribución de la acción automática al sistema & Responsabilidad & 4·4·4 & 4,0 & \balta{} \\
EC-19 & Protección ante el desecho accidental de una unidad & Protección contra errores & 5·3·2 & 3,8 & \balta{} \\
EC-41 & Advertencia previa a una acción irreversible & Advertencia de peligros & 5·3·2 & 3,8 & \balta{} \\
EC-14 & Reconocimiento sin contraprestación económica & Pertinencia funcional & 5·2·3 & 3,7 & \balta{} \\
EC-05 & No repudio de la bitácora de auditoría & No repudio & 4·3·4 & 3,7 & \balta{} \\
EC-08 & Aislamiento del fallo del servicio par & Tolerancia a fallos & 3·4·4 & 3,5 & \balta{} \\
EC-36 & Verificación exhaustiva de la máquina de estados & Capacidad de ser probado & 3·4·4 & 3,5 & \balta{} \\
EC-42 & Incorporación de una institución nueva a la red & Integración segura & 4·2·4 & 3,4 & \bmedia{} \\
EC-22 & Evolución del contrato entre servicios & Interoperabilidad & 3·3·5 & 3,4 & \bmedia{} \\
EC-07 & Disponibilidad de la consulta de inventario & Disponibilidad & 3·5·2 & 3,4 & \bmedia{} \\
EC-23 & Puesta en marcha del sistema completo & Instalabilidad & 3·4·3 & 3,3 & \bmedia{} \\
EC-40 & Detección de unidades próximas al vencimiento & Identificación de riesgos & 4·2·3 & 3,2 & \bmedia{} \\
EC-15 & Tiempo de respuesta de la consulta de inventario & Comportamiento temporal & 2·5·3 & 3,1 & \bmedia{} \\
EC-18 & Uso con lector de pantalla y sin ratón & Inclusividad & 3·3·3 & 3,0 & \bmedia{} \\
EC-25 & Cobertura del ciclo completo de la unidad & Completitud funcional & 3·3·3 & 3,0 & \bmedia{} \\
EC-09 & Recuperación tras pérdida de datos & Recuperabilidad & 4·1·3 & 2,9 & \bmedia{} \\
EC-37 & Sobrecarga de cuatro veces la carga nominal & Escalabilidad & 3·2·4 & 2,9 & \bmedia{} \\
EC-17 & Registro anónimo en tres pasos & Operabilidad & 2·5·2 & 2,9 & \bmedia{} \\
EC-06 & Resistencia ante abuso automatizado & Resistencia & 3·3·2 & 2,8 & \bmedia{} \\
EC-11 & Reanudación de notificaciones acumuladas & Recuperabilidad & 3·3·2 & 2,8 & \bmedia{} \\
EC-20 & Diagnóstico de un fallo en el ambiente de QA & Analizabilidad & 2·4·3 & 2,8 & \bmedia{} \\
EC-34 & Aislamiento del cambio dentro de un contexto & Modularidad & 2·2·5 & 2,6 & \bmedia{} \\
EC-28 & Identificación del propósito por un donante nuevo & Reconocibilidad & 2·4·2 & 2,6 & \bmedia{} \\
EC-30 & Retorno del donante tras recuperar la elegibilidad & Compromiso del usuario & 3·2·2 & 2,5 & \bmedia{} \\
EC-16 & Concurrencia durante el pico de una campaña & Capacidad & 2·3·3 & 2,5 & \bmedia{} \\
EC-21 & Cambio de la regla de elegibilidad & Modificabilidad & 2·2·4 & 2,4 & \bbaja{} \\
EC-32 & Legibilidad del estado de una unidad sin formación previa & Autodescripción & 2·4·1 & 2,4 & \bbaja{} \\
EC-24 & Extensión a otro tipo de donación & Adaptabilidad & 2·1·5 & 2,3 & \bbaja{} \\
EC-26 & Consumo bajo carga nominal sostenida & Utilización de recursos & 2·3·2 & 2,3 & \bbaja{} \\
EC-29 & Primer registro de donación por un operador nuevo & Aprendizaje & 2·3·2 & 2,3 & \bbaja{} \\
EC-31 & Recuperación ante un rechazo de operación & Asistencia al usuario & 2·3·1 & 2,1 & \bbaja{} \\
EC-27 & Convivencia de los contenedores del sistema en un mismo anfitrión & Coexistencia & 2·2·2 & 2,0 & \bbaja{} \\
EC-35 & Reutilización del componente de interfaz & Reusabilidad & 1·3·3 & 2,0 & \bbaja{} \\
EC-38 & Sustitución del API Gateway & Reemplazabilidad & 1·1·3 & 1,4 & \bbaja{} \\
\end{longtable}

\noindent\textbf{Lectura del catálogo.} Los dos escenarios críticos no comparten característica: uno es de seguridad y el otro de seguridad física, y son precisamente los dos cuyo incumplimiento tiene consecuencia fuera del sistema ---una sanción por tratamiento indebido de datos sensibles y una transfusión de una unidad que no debía despacharse. Los quince escenarios de las bandas crítica y alta concentran cinco de las seis subcaracterísticas de seguridad ---todas salvo la resistencia---, tres de las cinco de seguridad física y la totalidad de las que sostienen la trazabilidad. Los nueve escenarios de banda baja no son prescindibles: son los que admiten aplazamiento sin que la entrega pierda valor demostrable.

\subsection{Cuándo se verifica cada escenario}

Un escenario solo se puede comprobar cuando existe aquello que lo realiza. Comprometer una verificación antes de tiempo produce un resultado falso, no un adelanto.

\begin{longtable}{L{1.4cm} L{7.4cm} L{5.6cm}}
\caption{Sprint en que cada escenario se vuelve verificable y qué lo habilita.}\label{tab:verificacionec}\\
\toprule
\textbf{Sprint} & \textbf{Escenarios verificables} & \textbf{Qué lo habilita} \\
\midrule
\endfirsthead
\toprule
\textbf{Sprint} & \textbf{Escenarios verificables} & \textbf{Qué lo habilita} \\
\midrule
\endhead
\midrule
\contlinea{3} \\
\endfoot
\bottomrule
\endlastfoot
1 & EC-24, EC-25, EC-34, EC-36 & Se revisan sobre el modelo de dominio y la matriz de estados; no requieren código. La medida de EC-34 se verifica aquí en su parte de modelo y se completa en el Sprint 3 sobre el código. \\
2 & EC-22, EC-23, EC-27, EC-35, EC-38 & Contratos publicados, componentes contenerizados y capa de interfaz compartida en construcción. \\
3 & EC-01, EC-02, EC-03, EC-04, EC-05, EC-10, EC-12, EC-13, EC-17, EC-19, EC-21, EC-28, EC-31, EC-32, EC-33, EC-39, EC-41 & Gateway, Servicio de Identidad, autorización territorial y primer incremento funcional desplegado en QA. EC-17, EC-28 y EC-32 se verifican en una misma sesión con personas ajenas al equipo. \\
4 & EC-06, EC-07, EC-08, EC-11, EC-15, EC-16, EC-20, EC-26, EC-37, EC-40, EC-42 & Observabilidad instalada, sin la cual no hay medida de respuesta para los escenarios de carga, y escalamiento y notificaciones en operación. EC-08 exige los casos del ciclo de escalamiento y de la bandeja de salida, que existen desde este sprint; EC-42 exige comprobar las transferencias sugeridas, y EC-20, el plazo de diagnóstico sobre los tableros. \\
5 & EC-09, EC-14, EC-18, EC-29, EC-30 & Ensayo de restauración sobre el sistema completo, catálogo de reconocimientos cerrado y auditoría formal de accesibilidad y de aprendizaje sobre la interfaz ya estable. \\
\end{longtable}

\noindent Diecisiete de los cuarenta y dos escenarios se concentran en el Sprint 3, y doce de ellos pertenecen a las bandas crítica y alta. Es la constatación cuantitativa del \emph{killer} K-08: ese sprint sostiene el grueso de la verificación de calidad del sistema y es también donde se construyen el gateway, la autorización territorial y el primer flujo funcional completo. La consecuencia sobre la planeación es que en el Sprint 3 no se compromete alcance funcional adicional.

\subsection{Fichas de escenario}

Las fichas se presentan agrupadas por característica de ISO/IEC 25010:2023. El encabezado de cada una indica su índice de prioridad y su banda.

\subsubsection{Seguridad}

\begin{ficha}{EC-01}{Control de acceso por jurisdicción}{Confidencialidad · IP 4,3 · Alta}
Fuente del estímulo & Usuario territorial autenticado: coordinador territorial (U5), administrador de banco (U4) u operador (U3). \\
Estímulo & Intenta consultar o modificar información de un banco, un territorio o una unidad que pertenece a una jurisdicción distinta de la suya. \\
Contexto & Sistema en operación normal, usuario autenticado con token vigente. \\
Artefacto & API Gateway, filtrado por jurisdicción en cada servicio propietario, bitácora de auditoría: serie del Servicio de Identidad para las denegaciones del gateway y serie del servicio propietario para las suyas. \\
Respuesta & El sistema resuelve la jurisdicción del usuario a partir del token, rechaza la solicitud, responde sin confirmar ni negar que el recurso exista, y deja el intento en la bitácora de auditoría. \\
Medida de respuesta & El 100\,\% de las solicitudes fuera de jurisdicción se rechaza. Ninguna respuesta de rechazo incluye datos del recurso ni permite deducir su existencia. El 100\,\% de los intentos queda registrado con usuario, rol, recurso, marca de tiempo y resultado. \\
\end{ficha}

\noindent\emph{Consideraciones.} La restricción alcanza a los usuarios \textbf{territoriales}. El administrador nacional (U6) tiene el territorio nacional como jurisdicción propia, de modo que una consulta agregada del país no es un acceso fuera de jurisdicción y no debe rechazarse. La interfaz tampoco insinúa la existencia de datos ajenos: no muestra totales nacionales ni comparativos a un usuario territorial. Se prueba con los siete perfiles contra la matriz de autorización por grupo de operaciones del DD, noventa y una combinaciones, que es la especificación normativa del control de acceso (ADR-011); la matriz de usuarios contra módulos del SRS es una vista derivada de ella. Toda denegación que produce un servicio queda en su propia serie de auditoría (ADR-016). La autorización se evalúa del lado del servidor: ocultar un elemento en la interfaz no cuenta como control cumplido.

\begin{ficha}{EC-02}{Confidencialidad de la causa clínica}{Confidencialidad · IP 4,7 · Crítica}
Fuente del estímulo & Cualquier usuario con acceso al ciclo de vida de la unidad ---U3, U4, U7--- o un cliente que consulte la API directamente. \\
Estímulo & Consulta una unidad marcada como no apta, o su historial de trazabilidad, por la interfaz web o por el contrato de la API. \\
Contexto & Operación normal, usuario autenticado y dentro de su jurisdicción. \\
Artefacto & Modelo de dominio del Servicio de Donación, base de datos, contrato de la API, aplicación web, bitácora de auditoría. \\
Respuesta & El sistema devuelve el estado de aptitud y la trazabilidad de la unidad. No devuelve la causa clínica porque el dato no existe en el modelo, ni en el esquema, ni en el contrato. \\
Medida de respuesta & Cero campos de causa clínica, diagnóstico o resultado de tamizaje en el esquema, en el contrato OpenAPI y en cualquier respuesta. Revisión del 100\,\% de los \emph{endpoints} de M3 y M4. Ningún registro de bitácora, mensaje de error o traza contiene información diagnóstica. Cero ocurrencias de información diagnóstica en el campo de observación de los eventos de trazabilidad, verificado sobre el 100\,\% de los eventos del conjunto sintético y de los generados durante las pruebas. \\
\end{ficha}

\noindent\emph{Consideraciones.} La confidencialidad se garantiza en el \textbf{modelo}, no por filtrado en la capa de presentación: un dato que no se almacena no puede filtrarse por un descuido posterior. Se revisan también los mensajes de error y las trazas de diagnóstico, porque un texto como «rechazada por reactividad» viola el escenario igual que un campo. El rechazo de despacho de EC-12 debe redactarse sin motivo clínico. El campo de observación de los eventos de trazabilidad es la única vía del modelo por la que un operador podría escribir un motivo clínico por descuido, y por eso la medida lo nombra. Fundamento: Ley 1581 de 2012 y restricción RD-08.

\begin{ficha}{EC-03}{Autenticidad de la sesión}{Autenticidad · IP 4,3 · Alta}
Fuente del estímulo & Cliente cualquiera: la aplicación web, una herramienta externa o un agente no identificado. \\
Estímulo & Envía una petición con el token ausente, expirado, alterado en su contenido o firmado por un emisor que el sistema no reconoce. \\
Contexto & Operación normal, \emph{endpoint} que requiere sesión. \\
Artefacto & API Gateway, Servicio de Identidad como emisor único, validación propia en cada servicio que recibe peticiones. \\
Respuesta & El gateway rechaza la petición antes de enrutarla. La lógica de negocio no se ejecuta y no se consulta la base de datos. El servicio destino vuelve a verificar el token por su cuenta y no confía en que el gateway ya lo hizo. \\
Medida de respuesta & Los cuatro casos ---ausente, expirado, alterado, emisor desconocido--- se rechazan en el 100\,\% de los intentos, junto con los de audiencia incorrecta y token de servicio en una ruta de usuario. Ninguna petición sin token válido alcanza un servicio, verificado en la traza de los registros. El token de acceso vence a los quince minutos y la sesión de renovación, a las ocho horas como máximo. El 100\,\% de las denegaciones del gateway queda en la serie de auditoría del Servicio de Identidad o en el contador de denegaciones no entregadas. \\
\end{ficha}

\noindent\emph{Consideraciones.} El emisor del token es el Servicio de Identidad, único componente que conoce la clave privada de firma (ADR-008). Los demás verifican la firma RS256 contra sus claves públicas, de modo que comprometer cualquier otro componente no permite fabricar tokens. La jurisdicción viaja en el token, lo que evita una llamada a Identidad en cada verificación. La doble verificación es deliberada, no redundante: un servicio alcanzable por otra ruta no queda desprotegido. Se prueba invocando un \emph{endpoint} interno saltándose el gateway; debe rechazar igual.

\begin{ficha}{EC-04}{Integridad del ciclo de vida de la unidad}{Integridad · IP 4,2 · Alta}
Fuente del estímulo & Operador de banco (U3), o un cliente que invoca la API sin pasar por la interfaz. \\
Estímulo & Intenta llevar una unidad a un estado que no corresponde a su estado actual: devolver a disponible una unidad ya desechada, transferir una unidad no apta, o alterar su origen ya registrado. \\
Contexto & Sistema en operación normal. \\
Artefacto & Máquina de estados de la unidad en el Servicio de Donación, base de datos, historial de trazabilidad. \\
Respuesta & El sistema rechaza la transición, conserva el estado anterior sin modificarlo y registra el intento. El historial solo admite anexar eventos: nunca se edita ni se borra lo ya escrito. \\
Medida de respuesta & Cero transiciones inválidas aceptadas sobre la matriz completa de pares estado origen--estado destino. Ninguna unidad queda en un estado no previsto. Cero operaciones de modificación o borrado expuestas sobre el historial. El origen y la fecha de captación no son modificables tras el registro. \\
\end{ficha}

\noindent\emph{Consideraciones.} La matriz de transiciones válidas se prueba de forma exhaustiva y no por muestreo: el número de pares es pequeño y comprobarlos todos es viable, según establece EC-36. La restricción se aplica en el servicio y no solo en la interfaz, porque el estímulo contempla la llamada directa a la API. El historial de solo anexado es la base de RF-04 y de la reconstrucción que hace el auditor en EC-05.

\begin{ficha}{EC-05}{No repudio de la bitácora de auditoría}{No repudio · IP 3,7 · Alta}
Fuente del estímulo & Auditor (U7), o una investigación posterior a un incidente. \\
Estímulo & Se requiere establecer quién ejecutó una acción sensible: marcar una unidad como no apta, ejecutar su disposición final, aprobar una transferencia o intentar un acceso fuera de jurisdicción. \\
Contexto & Sistema en operación, consulta posterior al hecho. \\
Artefacto & Servicio transversal de auditoría: series de los Servicios de Identidad, Institucional, Campañas y Donación, consulta compuesta del Servicio Institucional, pantalla del auditor. \\
Respuesta & Cada acción sensible quedó registrada con actor, rol, jurisdicción, marca de tiempo, recurso afectado y resultado. Ningún usuario de la aplicación, incluido el administrador nacional, puede borrar ni alterar un registro desde el sistema. \\
Medida de respuesta & El 100\,\% de las acciones del catálogo de acciones sensibles queda registrado. Cero operaciones de actualización o borrado sobre la bitácora expuestas por la aplicación. El auditor reconstruye la cadena completa de eventos de una unidad, de captación a disposición final, en menos de cinco minutos y sin apoyo del equipo técnico, sobre la consulta compuesta que define ADR-016. Ante indisponibilidad de uno o varios de los contextos consultados, la respuesta marca \texttt{bitacora\_parcial} y nombra los contextos ausentes: cero respuestas parciales presentadas como completas. \\
\end{ficha}

\noindent\emph{Consideraciones.} El catálogo de acciones sensibles se define de forma explícita en el diseño y se revisa cada sprint: una acción sensible que no esté en el catálogo no se registra, y el hueco no se nota hasta que se necesita. El auditor no escribe en ningún módulo (RI-04). La bitácora registra el intento fallido de EC-01 con el mismo detalle que una acción exitosa. Los registros no contienen causa clínica ni datos personales innecesarios para identificar la acción.

\begin{ficha}{EC-06}{Resistencia ante abuso automatizado}{Resistencia · IP 2,8 · Media}
Fuente del estímulo & Agente externo no autenticado que genera \textbf{600 peticiones por minuto desde un mismo origen}, diez veces la tasa nominal admitida por origen. \\
Estímulo & Envía peticiones repetidas contra el registro anónimo o el inicio de sesión, o recorre identificadores de unidad o códigos de intención de forma sistemática para enumerarlos. \\
Contexto & Sistema desplegado en QA, expuesto en red, sin sesión válida, con tráfico legítimo simultáneo de 20 usuarios concurrentes. \\
Artefacto & API Gateway, rutas públicas de los Servicios de Donación, Identidad y Campañas. \\
Respuesta & El gateway limita la tasa de peticiones por origen y rechaza por exceso de solicitudes antes de enrutar. El tráfico legítimo sigue siendo atendido. Los identificadores de unidad no son adivinables, de modo que recorrerlos no produce información. \\
Medida de respuesta & Superado el umbral de 60 peticiones por minuto y por origen, el 100\,\% de las peticiones excedentes a rutas públicas se rechaza con código 429 y cabecera \texttt{Retry-After} sin alcanzar ningún servicio. Los identificadores de unidad son valores no secuenciales. La consulta pública de una intención por su código queda sometida al mismo umbral, de modo que ningún origen prueba más de sesenta códigos por minuto. El tiempo de respuesta del tráfico legítimo no se degrada más de un 20\,\% durante la prueba. \\
\end{ficha}

\noindent\emph{Consideraciones.} Es un escenario de resistencia, no de disponibilidad ante ataque masivo: el objetivo es que el abuso simple no degrade el servicio ni permita inferir el tamaño ni el contenido del inventario. Los identificadores no secuenciales protegen además a EC-01, porque un identificador adivinable permitiría deducir la existencia de unidades de otra jurisdicción aunque la consulta se rechace. El umbral de 60 peticiones por minuto y por origen rige \textbf{únicamente sobre los endpoints públicos}, que son los que crean o consultan datos sin sesión previa. El tráfico autenticado se limita por sesión y no por origen, con un umbral independiente y más alto. La distinción no es un detalle de configuración: un banco cuyos operadores comparten una sola dirección pública de salida constituye un único origen, y un límite por origen aplicado al tráfico autenticado lo estrangularía como si fuera un cliente abusivo. Por la misma razón, las mediciones de EC-07, EC-15, EC-16, EC-26 y EC-37 se ejecutan desde un origen declarado y exento del límite anónimo, sin el cual una prueba de carga de cincuenta usuarios concurrentes se rechazaría a sí misma. Ambos umbrales residen en la configuración del gateway, no en el código, y ADR-007 los fija: 60 peticiones por minuto y por origen en las rutas públicas, y 300 por minuto y por sesión en las protegidas. El proxy de borde no limita la tasa (ADR-015).

\begin{ficha}{EC-33}{Atribución de la acción automática al sistema}{Responsabilidad · IP 4,0 · Alta}
Fuente del estímulo & Proceso programado de vencimiento, o cualquier ejecución automática que cambie el estado de una unidad sin usuario conectado. \\
Estímulo & El sistema ejecuta por su cuenta una acción sensible mientras hay o no hay sesiones activas. \\
Contexto & Operación normal, sin intervención humana en la acción. \\
Artefacto & Proceso programado del Servicio de Donación, servicio transversal de auditoría. \\
Respuesta & El registro de auditoría atribuye la acción al sistema como actor, con el identificador del proceso que la ejecutó, y no a ninguna persona. \\
Medida de respuesta & El 100\,\% de las acciones automáticas queda registrado con tipo de actor \emph{sistema}. Cero registros de acción automática atribuidos a un usuario. El auditor distingue en la bitácora, sin ambigüedad, qué acciones ejecutó una persona y cuáles el sistema. \\
\end{ficha}

\noindent\emph{Consideraciones.} La responsabilidad se distingue del no repudio de EC-05: aquel garantiza que la acción de una persona quede atribuida a esa persona; este garantiza que la acción del sistema no quede atribuida a ninguna. Atribuir a un operador el vencimiento automático de una unidad crea una responsabilidad falsa que un auditor tomaría por verdadera, y es el modo de fallo más probable porque la implementación sencilla reutiliza el usuario de la última sesión. El escenario alcanza también a las transferencias sugeridas por el ciclo de escalamiento (ADR-013), que el proceso crea con su token de servicio y que la bitácora atribuye al sistema.

\subsubsection{Seguridad física}

\begin{ficha}{EC-12}{Bloqueo del despacho de una unidad no apta o vencida}{A prueba de fallos · IP 4,5 · Crítica}
Fuente del estímulo & Operador de banco (U3) o coordinador que gestiona una transferencia entre bancos. \\
Estímulo & Intenta despachar, transferir o reservar una unidad marcada como no apta o cuya fecha de vencimiento ya se cumplió. \\
Contexto & Operación normal, o degradada porque el estado de la unidad no se puede confirmar en ese momento. \\
Artefacto & Servicio de Donación: inventario, transferencias y disposición final. \\
Respuesta & El sistema rechaza la operación. Ante estado indeterminado o dato faltante, opta siempre por la alternativa más restrictiva y trata la unidad como no disponible. La unidad no aparece en ningún listado de disponibles, tampoco en el que alimenta las sugerencias de transferencia. \\
Medida de respuesta & Cero unidades no aptas o vencidas presentes en cualquier respuesta de unidades disponibles, incluidas la ruta de transferencia entre bancos y las transferencias sugeridas por el escalamiento. Ante estado indeterminado, el 100\,\% de las unidades se trata como no disponible, verificado con un caso de prueba propio que fuerza la indisponibilidad del dato de estado. El mensaje de rechazo no revela causa clínica. \\
\end{ficha}

\noindent\emph{Consideraciones.} Es el único escenario del catálogo cuyo incumplimiento afecta a la integridad física de una persona: una unidad no apta que llega a transfusión es el fallo que el sistema existe para impedir. El comportamiento por defecto ante la duda es el rechazo; un sistema que asume «disponible» cuando no puede confirmar el estado falla del lado peligroso. La regla se aplica en toda ruta que consulte disponibilidad, no solo en la pantalla de despacho: la sugerencia automática de transferencia es una ruta distinta y suele olvidarse.

\begin{ficha}{EC-39}{Límite operativo del despacho por jurisdicción}{Restricción operativa · IP 4,0 · Alta}
Fuente del estímulo & Operador o administrador de un banco que intenta despachar o transferir fuera de las condiciones operativas admitidas. \\
Estímulo & Solicita el despacho de una unidad custodiada por otra institución, o una transferencia hacia una institución que no está habilitada como banco de sangre. \\
Contexto & Operación normal, usuario autenticado y dentro de su jurisdicción. \\
Artefacto & Servicio de Donación (transferencias), Servicio Institucional (registro de instituciones). \\
Respuesta & El sistema rechaza la operación por violar una condición operativa declarada, indica cuál fue la condición incumplida y registra el intento. \\
Medida de respuesta & Cero despachos ejecutados por una institución distinta de la custodia. Cero transferencias hacia instituciones cuyo tipo registrado no admite custodia de unidades. El catálogo de tipos de institución distingue, como mínimo, el banco de sangre, que admite custodia, del servicio de transfusión, que no la admite, conforme a la Resolución 901 de 1996. El 100\,\% de los rechazos identifica la condición operativa incumplida, sin revelar datos de la institución ajena. \\
\end{ficha}

\noindent\emph{Consideraciones.} La restricción operativa es distinta del control de acceso de EC-01: allí el usuario no tiene derecho a ver el recurso; aquí lo ve y tiene derecho a verlo, pero la operación que pide no es admisible sobre él. La distinción del tipo de institución ---banco de sangre frente a servicio de transfusión--- proviene de la Resolución 901 de 1996: no todo prestador adscrito puede custodiar unidades, y confundirlo movería sangre a un lugar sin condiciones de conservación.

\begin{ficha}{EC-40}{Detección de unidades próximas al vencimiento}{Identificación de riesgos · IP 3,2 · Media}
Fuente del estímulo & El paso del tiempo sobre el inventario de una institución, evaluado por el proceso programado. \\
Estímulo & Una o varias unidades alcanzan el umbral de días previos al vencimiento configurado por la institución. \\
Contexto & Operación normal, con el inventario cargado y umbrales configurados. \\
Artefacto & Servicio de Donación: inventario, proceso programado de alertas, servicio transversal de notificaciones. \\
Respuesta & El sistema identifica la situación de riesgo antes de que se materialice, genera la alerta de vencimiento próximo y la dirige al administrador del banco custodio. \\
Medida de respuesta & El 100\,\% de las unidades que cruzan el umbral genera alerta dentro de las 24 horas siguientes. Cero unidades que venzan sin haber generado antes su alerta. El umbral configurable admite valores de un día o más, y una unidad que queda disponible con una vida útil restante menor que el umbral genera su alerta en el mismo momento en que queda disponible. La alerta identifica componente, grupo sanguíneo y días restantes, y no contiene dato clínico alguno. \\
\end{ficha}

\noindent\emph{Consideraciones.} La identificación del riesgo es lo que hace útil el escalamiento en dos niveles: una unidad detectada con antelación puede transferirse a un banco con déficit, mientras que una detectada el día del vencimiento solo puede desecharse. El valor del escenario está en la antelación, no en la alerta: por eso la medida exige que ninguna unidad venza sin alerta previa, y no solo que las alertas se generen.

\begin{ficha}{EC-41}{Advertencia previa a una acción irreversible}{Advertencia de peligros · IP 3,8 · Alta}
Fuente del estímulo & Operador de banco (U3) o administrador (U4) a punto de ejecutar una acción sin retorno. \\
Estímulo & Inicia una de las tres acciones irreversibles del sistema: la disposición final de una unidad, la decisión sobre una transferencia o el despacho de una unidad. \\
Contexto & Operación normal, con varias unidades en pantalla. \\
Artefacto & Aplicación web, Servicio de Donación. \\
Respuesta & El sistema advierte de forma explícita, antes de ejecutar, que la acción es irreversible; nombra la unidad concreta y la consecuencia; y exige una confirmación deliberada. \\
Medida de respuesta & El 100\,\% de las acciones irreversibles presenta, antes de la confirmación, una advertencia que nombra la unidad, declara que la acción no tiene retorno y enuncia su consecuencia. Cero acciones irreversibles ejecutables sin advertencia. La advertencia es perceptible por lector de pantalla y no depende solo del color. \\
\end{ficha}

\noindent\emph{Consideraciones.} La advertencia de peligro precede a la protección contra el error de EC-19 y no la sustituye: aquella informa de la consecuencia antes de actuar, y esta impide que la acción se dispare por descuido. Ambas se aplican al mismo conjunto cerrado de tres acciones irreversibles definido en el glosario. Una interfaz puede pedir confirmación sin advertir de nada ---«¿está seguro?»--- y en ese caso cumple EC-19 e incumple este escenario.

\begin{ficha}{EC-42}{Incorporación de una institución nueva a la red}{Integración segura · IP 3,4 · Media}
Fuente del estímulo & Coordinador territorial (U5) o administrador nacional (U6) que da de alta una institución. \\
Estímulo & Registra una institución nueva y le asigna usuarios, jurisdicción y umbrales de inventario. \\
Contexto & Sistema en operación, con el resto de la red operando con normalidad. \\
Artefacto & Servicio Institucional (registro de instituciones), Servicio de Identidad (usuarios y jurisdicción), Servicio de Donación (inventario y transferencias). \\
Respuesta & La institución queda incorporada sin exponer datos de otras instituciones, sin unidades preexistentes atribuidas y sin alterar el inventario ni las alertas del resto de la red. \\
Medida de respuesta & Cero unidades atribuidas a la institución nueva en el momento del alta. Cero cambios en el inventario, las alertas o las sugerencias de transferencia de las instituciones existentes. Los usuarios creados no acceden a ningún dato anterior a su alta. El alta queda registrada en la bitácora con su responsable. \\
\end{ficha}

\noindent\emph{Consideraciones.} El riesgo que este escenario controla es que la incorporación de un participante nuevo altere el estado de los que ya operan, que es el modo de fallo característico de una red que crece. El caso concreto a vigilar es el de las sugerencias de transferencia: una institución nueva y vacía no debe aparecer como origen de excedente ni desplazar a un banco que sí lo tiene.

\subsubsection{Fiabilidad}

\begin{ficha}{EC-07}{Disponibilidad de la consulta de inventario}{Disponibilidad · IP 3,4 · Media}
Fuente del estímulo & Carga nominal de \textbf{20 usuarios concurrentes} ---administradores y operadores--- durante la jornada operativa. \\
Estímulo & Consultan el inventario de su institución mientras un contenedor se reinicia por un despliegue o por una caída inducida. \\
Contexto & Ambiente de QA en operación, con el conjunto sintético de referencia cargado. \\
Artefacto & Servicio de Donación, API Gateway, base de datos, orquestación de contenedores. \\
Respuesta & La consulta responde. Si el contenedor del servicio cae, la política de reinicio lo levanta y el sondeo de salud lo devuelve a servicio sin que nadie intervenga. \\
Medida de respuesta & Disponibilidad igual o superior al 99\,\% sobre el \emph{endpoint} de salud, medida en la ventana de evaluación del sprint. Un contenedor caído vuelve a estado saludable en 60 segundos o menos. Cero reinicios que requieran intervención manual. \\
\end{ficha}

\noindent\emph{Consideraciones.} La disponibilidad se mide, no se estima: la serie proviene de Prometheus y se publica en Grafana. Cada contenedor expone un sondeo de salud que verifica también su dependencia de base de datos, no solo que el proceso esté vivo; un servicio que responde pero no alcanza su base no está disponible. El objetivo del 99\,\% corresponde a un proyecto académico con infraestructura de nivel gratuito y así se presenta al cliente.

\begin{ficha}{EC-08}{Aislamiento del fallo del servicio par}{Tolerancia a fallos · IP 3,5 · Alta}
Fuente del estímulo & Fallo de infraestructura o despliegue en curso. \\
Estímulo & Un servicio par deja de responder mientras otro lo necesita: el Servicio Institucional o el de Campañas mientras el Servicio de Donación los consulta, el Institucional durante el ciclo de escalamiento, o cualquiera de los contextos que compone la consulta del auditor. \\
Contexto & Sistema en operación, un solo servicio caído. \\
Artefacto & Cliente HTTP de cada servicio que tiene arista hacia otro, API Gateway, aplicación web. \\
Respuesta & La llamada expira en un plazo acotado en lugar de quedar esperando. Las funciones que no dependen del servicio caído siguen operando. La interfaz indica con precisión qué quedó no disponible, sin exponer detalle técnico ni dejar la pantalla en blanco. \\
Medida de respuesta & Tiempo de espera máximo de 3 segundos por llamada entre servicios. Cero hilos o conexiones bloqueados más allá de ese plazo. Las operaciones que no requieren el servicio caído conservan el 100\,\% de su funcionalidad; en particular, las de M3 y M4 con el Institucional o Campañas caídos. Con el Institucional caído, el ciclo de escalamiento termina sin crear nada y se reintenta en el ciclo siguiente. Cero errores no controlados en la interfaz durante la indisponibilidad. \\
\end{ficha}

\noindent\emph{Consideraciones.} Este escenario es la contrapartida de la arquitectura distribuida: separar los servicios solo aporta si el fallo de uno no arrastra al otro; sin él, la distribución añade puntos de fallo sin beneficio. Un plazo de espera sin límite es el fallo más común y el más caro, porque agota el \emph{pool} de conexiones y convierte la caída de un servicio en la de todos los que lo consultan. Con cinco servicios, el escenario se prueba sobre cada arista entre servicios de la tabla de aristas permitidas (ADR-014). La respuesta degradada nunca inventa datos.

\begin{ficha}{EC-09}{Recuperación tras pérdida de datos}{Recuperabilidad · IP 2,9 · Media}
Fuente del estímulo & Fallo de infraestructura, error humano o corrupción del volumen de datos. \\
Estímulo & Se pierde o se corrompe la base de datos de uno de los cuatro servicios con base. \\
Contexto & Ambiente de QA desplegado, con datos sintéticos de referencia. \\
Artefacto & PostgreSQL de cada servicio, respaldos, trabajos de migración. \\
Respuesta & Se restaura el último respaldo disponible y se reaplican las migraciones pendientes. El servicio vuelve a operar con el esquema en la versión esperada, sin reconstruir nada a mano. \\
Medida de respuesta & Tiempo de recuperación igual o menor a 4 horas. Pérdida máxima de datos de 24 horas. Cero diferencias entre la versión de esquema esperada y la obtenida tras restaurar. El procedimiento está documentado y se ensaya en cada sprint desde el Sprint 2 sobre las bases que existen en él; los ensayos de los Sprints 2 a 4 validan el procedimiento, y el del Sprint 5 verifica la medida, porque es el primero que se ejecuta sobre el sistema completo con datos de todos los módulos. \\
\end{ficha}

\noindent\emph{Consideraciones.} La restauración se ensaya de verdad: un respaldo que nunca se ha restaurado no es un respaldo, es una suposición. El versionado del esquema por migraciones es lo que hace verificable la medida. La pérdida admitida de 24 horas es aceptable porque el ambiente opera con datos sintéticos; en un despliegue real el requisito sería otro y se declara así ante el cliente. Cada servicio se recupera de forma independiente.

\begin{ficha}{EC-10}{Registro de donación idempotente}{Ausencia de fallos · IP 4,0 · Alta}
Fuente del estímulo & Operador de banco (U3). \\
Estímulo & Envía dos veces el mismo registro de donación, por doble pulsación, por reintento del navegador o porque la red cortó la respuesta y el cliente reenvió. \\
Contexto & Jornada operativa con conectividad intermitente. \\
Artefacto & Servicio de Donación, base de datos, inventario. \\
Respuesta & El sistema crea una sola unidad. La segunda petición devuelve el mismo resultado que la primera en lugar de crear un duplicado o de fallar con un error que confunda al operador. \\
Medida de respuesta & Cero unidades duplicadas ante 100 reenvíos del mismo identificador de operación. El conteo del inventario no varía por efecto del reintento. La segunda respuesta es idéntica a la primera. \\
\end{ficha}

\noindent\emph{Consideraciones.} Una unidad duplicada es peor que una donación perdida: infla el inventario disponible y puede llevar a no escalar una transferencia que sí hacía falta. La idempotencia se resuelve con un identificador de operación que genera el cliente y que el servicio usa como clave, no con una comprobación de «ya existe algo parecido». La regla aplica a las seis operaciones reenviables del contrato: registro de intención, registro de donante, registro de donación, disposición final, creación de transferencia y decisión sobre una transferencia. Se prueba en el nivel de la API: deshabilitar el botón tras el primer clic ayuda al usuario pero no protege al servicio.

\begin{ficha}{EC-11}{Reanudación de notificaciones acumuladas}{Recuperabilidad · IP 2,8 · Media}
Fuente del estímulo & Acumulación de \textbf{500 eventos notificables} generados durante 30 minutos de indisponibilidad del canal de salida. \\
Estímulo & Los eventos se producen ---elegibilidad recuperada, campaña cercana, reconocimiento obtenido o convocatoria por escasez--- mientras el canal no está disponible. \\
Contexto & Operación degradada: el sistema funciona, el canal de salida o el trabajador de Notificaciones no. \\
Artefacto & Bandeja de salida del Servicio de Donación y trabajador de Notificaciones (ADR-014). \\
Respuesta & Los eventos se conservan y se reintentan con espera creciente. Al restablecerse el canal, la cola se drena sin pérdida y sin duplicar avisos. \\
Medida de respuesta & El 100\,\% de los 500 eventos se entrega dentro de los 15 minutos siguientes al restablecimiento del canal. Cero notificaciones duplicadas al mismo donante por el mismo evento. Agotados los reintentos, el evento queda en estado de revisión manual y visible en el tablero, nunca descartado sin registro. \\
\end{ficha}

\noindent\emph{Consideraciones.} El caso que importa es la convocatoria por escasez: un aviso que se pierde cuando falta un componente sanguíneo es precisamente el que no puede perderse. Por esa razón su impacto es 3 y no 2: la pérdida del aviso retrasa la cobertura de un déficit, que es una capacidad esencial del servicio. La notificación es asíncrona por diseño; el evento se escribe en la bandeja en la misma transacción que el hecho que lo origina, y que el canal o el trabajador fallen no bloquea el registro de la donación. El descarte silencioso es el peor resultado posible, peor que el reintento fallido visible. Ninguna notificación incluye causa clínica ni ofrece contraprestación económica.

\subsubsection{Adecuación funcional}

\begin{ficha}{EC-25}{Cobertura del ciclo completo de la unidad}{Completitud funcional · IP 3,0 · Media}
Fuente del estímulo & Cliente, durante la demostración de aceptación de la entrega. \\
Estímulo & Recorre el ciclo de vida completo de una unidad, de la captación a la disposición final, sin salir del sistema en ningún paso. \\
Contexto & Ambiente de QA con datos sintéticos y los perfiles necesarios disponibles. \\
Artefacto & Módulos M1, M3 y M4 del Servicio de Donación, aplicación web. \\
Respuesta & El sistema cubre todos los estados de la máquina y todas las transiciones válidas mediante alguna operación disponible en la interfaz, sin exigir intervención en la base de datos ni herramientas externas. \\
Medida de respuesta & Las 13 transiciones válidas de la máquina de estados, más el alta de la unidad, tienen operación asociada en la interfaz o en un proceso automático. Cero pasos del ciclo que requieran intervención fuera del sistema. Cero estados alcanzables desde los que no exista salida prevista. \\
\end{ficha}

\noindent\emph{Consideraciones.} La completitud se mide contra la máquina de estados y no contra el catálogo de requerimientos, porque un estado sin salida es un hueco funcional aunque todos los requerimientos estén implementados. El caso que este escenario detecta es el de una unidad que llega a un estado del que ninguna pantalla permite sacarla, y que solo se descubre en la demostración.

\begin{ficha}{EC-13}{Exclusión automática de unidades vencidas}{Corrección funcional · IP 4,1 · Alta}
Fuente del estímulo & El paso del tiempo. No hay usuario involucrado. \\
Estímulo & Una unidad alcanza su fecha de vencimiento mientras figura como disponible en el inventario. \\
Contexto & Operación normal, sin que ningún usuario esté conectado. \\
Artefacto & Servicio de Donación: inventario, proceso programado de vencimiento, bitácora. \\
Respuesta & La unidad sale del inventario disponible y entra al flujo de disposición final, sin que nadie ejecute una acción ni confirme nada. \\
Medida de respuesta & En el 100\,\% de los casos la unidad deja de figurar como disponible dentro de los 15 minutos siguientes al vencimiento. Cero unidades vencidas visibles como disponibles en cualquier consulta, listado o sugerencia de transferencia. El cambio queda en la bitácora con el sistema como actor. \\
\end{ficha}

\noindent\emph{Consideraciones.} El requisito es explícito en que la exclusión ocurre \textbf{sin intervención manual}: una pantalla que pida al operador «revisar vencimientos» no cumple el escenario aunque produzca el mismo resultado. El plazo de quince minutos fija la frecuencia del proceso programado de vencimiento: se ejecuta cada quince minutos o menos. La prueba se hace adelantando el reloj sobre datos sintéticos. La atribución del cambio al sistema y no al último usuario conectado es lo que verifica EC-33.

\begin{ficha}{EC-14}{Reconocimiento sin contraprestación económica}{Pertinencia funcional · IP 3,7 · Alta}
Fuente del estímulo & Donante registrado (U2) que alcanza un hito de donación recurrente, o el equipo al definir un reconocimiento nuevo. \\
Estímulo & El sistema debe otorgar y mostrar un reconocimiento al donante. \\
Contexto & Sistema en operación normal. \\
Artefacto & M1: catálogo de reconocimientos, motor de niveles e insignias, aplicación web, notificaciones. \\
Respuesta & El reconocimiento es simbólico: insignia, nivel o constancia de participación. El catálogo no admite dinero, canje, descuento, sorteo ni prioridad de atención en el servicio de salud. \\
Medida de respuesta & Cero elementos del catálogo con valor económico directo o indirecto, verificado contra el Decreto 1571 de 1993. Cero textos de la interfaz o de las notificaciones que sugieran contraprestación. El \emph{checklist} de Ready lo comprueba en el 100\,\% de las historias de M1 antes de iniciarlas. \\
\end{ficha}

\noindent\emph{Consideraciones.} La restricción es legal, no de producto. El riesgo real está en la redacción, no en el modelo de datos: una insignia llamada «premio» o un texto que prometa «beneficios» incumple aunque no exista ningún objeto de valor asociado. La prioridad de atención cuenta como contraprestación indirecta y por eso está enumerada de forma explícita.

\subsubsection{Eficiencia de desempeño}

\begin{ficha}{EC-15}{Tiempo de respuesta de la consulta de inventario}{Comportamiento temporal · IP 3,1 · Media}
Fuente del estímulo & Carga nominal de \textbf{20 usuarios concurrentes} entre administradores (U4), operadores (U3) y coordinadores territoriales (U5). \\
Estímulo & Solicitan el inventario de su institución o de su jurisdicción, con filtro por componente sanguíneo y tipo de sangre. \\
Contexto & Operación normal en QA, con el volumen sintético de referencia cargado. \\
Artefacto & Aplicación web, API Gateway, Servicio de Donación, base de datos. \\
Respuesta & El sistema entrega el resultado completo, con las alertas de umbral y de vencimiento incluidas. \\
Medida de respuesta & Percentil 95 igual o menor a 3 segundos, medido en el gateway sobre el conjunto sintético de referencia. Percentil 99 igual o menor a 5 segundos. La medición se toma de Prometheus y se publica en Grafana. \\
\end{ficha}

\noindent\emph{Consideraciones.} RNF-04 habla de «menos de 3 segundos en condiciones normales»; el escenario lo traduce a percentiles porque un promedio esconde justamente los casos lentos que el usuario recuerda. Se mide en el gateway y no dentro del servicio: lo que importa es el tiempo que percibe quien consulta, incluida la red y el enrutamiento. El volumen sintético de referencia se declara y se versiona, porque una medida de desempeño sin volumen declarado no es comparable entre sprints. El inventario se calcula como proyección sobre las unidades en cada consulta; si la medición no cumple este escenario, la contingencia es materializar esa proyección, y esa decisión se registra en un ADR antes de implementarse.

\begin{ficha}{EC-26}{Consumo bajo carga nominal sostenida}{Utilización de recursos · IP 2,3 · Baja}
Fuente del estímulo & Carga nominal sostenida de \textbf{50 usuarios concurrentes durante 30 minutos} sobre el conjunto sintético de referencia. \\
Estímulo & La carga se mantiene sin interrupción, combinando consulta de inventario, registro de donaciones y consulta de campañas. \\
Contexto & Pila de medición con los catorce contenedores de QA desplegados y límites de recurso asignados. \\
Artefacto & Los catorce contenedores de QA, con especial atención a las cuatro bases de datos. \\
Respuesta & El sistema sostiene la carga sin crecimiento acumulativo de memoria ni agotamiento de conexiones, y devuelve el consumo a su nivel de reposo al cesar la carga. \\
Medida de respuesta & Uso de CPU por debajo del 80\,\% y de memoria por debajo del 75\,\% del límite asignado a cada contenedor. El consumo de memoria al final de los 30 minutos no supera en más de un 10\,\% al de los primeros 5 minutos. Cero contenedores reiniciados por agotamiento de recursos. Ninguna base de datos supera el 80\,\% de su límite de conexiones. \\
\end{ficha}

\noindent\emph{Consideraciones.} La duración sostenida es el elemento esencial del escenario, no la cifra de concurrencia: una prueba corta no distingue un sistema estable de uno que acumula memoria o conexiones y caería al cabo de unas horas. La comparación entre el consumo del minuto 5 y el del minuto 30 es lo que detecta esa acumulación, y por eso la medida está expresada como diferencia y no como valor absoluto.

\begin{ficha}{EC-16}{Concurrencia durante el pico de una campaña}{Capacidad · IP 2,5 · Media}
Fuente del estímulo & Carga de pico de \textbf{50 usuarios concurrentes durante 10 minutos}, entre donantes y operadores de banco. \\
Estímulo & Consultan campañas, se registran y registran donaciones al mismo tiempo, durante el pico de una campaña publicada. \\
Contexto & Campaña activa, ambiente de QA con los cinco servicios desplegados y el volumen sintético cargado, sobre la pila de medición. \\
Artefacto & API Gateway, Servicio de Campañas, Servicio de Donación y sus bases de datos. \\
Respuesta & El sistema atiende la carga sin perder peticiones. Si la carga excede su capacidad, degrada de forma controlada rechazando con un código explícito, en lugar de agotar recursos y fallar de manera imprevisible. \\
Medida de respuesta & Tasa de error inferior al 1\,\% y percentil 95 igual o menor a 3 segundos durante los 10 minutos. El uso de CPU y memoria de cada contenedor se mantiene por debajo del 80\,\% de su límite asignado. Cero peticiones perdidas sin respuesta. \\
\end{ficha}

\noindent\emph{Consideraciones.} La cifra de 50 usuarios concurrentes es la capacidad nominal declarada del sistema en la escala de demostración del proyecto y se ajusta si el cliente fija otra. Perder una petición en silencio es peor que rechazarla: el donante cree haberse registrado y no lo está. La prueba de carga se ejecuta contra datos sintéticos y nunca contra el ambiente que se usa para la demostración al cliente.

\subsubsection{Capacidad de interacción}

\begin{ficha}{EC-17}{Registro anónimo en tres pasos}{Operabilidad · IP 2,9 · Media}
Fuente del estímulo & Donante anónimo (U1) que usa el sistema por primera vez y sin ayuda de nadie. \\
Estímulo & Quiere registrar su intención de donar sin entregar datos personales completos. \\
Contexto & Navegador de escritorio o móvil, sin sesión previa ni cuenta creada. \\
Artefacto & Aplicación web, M1. \\
Respuesta & Completa el registro en un máximo de tres pasos, sin ningún campo de identificación obligatorio, y recibe confirmación de que quedó registrado. \\
Medida de respuesta & Tres pasos como máximo, contados en el indicador de progreso de la propia pantalla. Cero campos de identificación marcados como obligatorios. En prueba con cinco personas ajenas al equipo, al menos cuatro completan el registro sin ayuda en dos minutos o menos. \\
\end{ficha}

\noindent\emph{Consideraciones.} El límite de tres pasos es duro: cualquier necesidad nueva se acomoda dentro de los tres pasos o va a la modalidad con perfil completo, que es voluntaria y no tiene ese límite. La modalidad anónima no se presenta como inferior a la registrada. La barrera de entrada baja es el mecanismo del sistema para atacar el problema de fondo: solo el 27,1\,\% de los donantes del país son habituales. Se prueba con personas ajenas al equipo.

\begin{ficha}{EC-18}{Uso con lector de pantalla y sin ratón}{Inclusividad · IP 3,0 · Media}
Fuente del estímulo & Donante o funcionario que usa lector de pantalla, que navega únicamente con teclado, o que tiene baja visión o daltonismo. \\
Estímulo & Recorre el registro anónimo y la consulta de inventario de principio a fin. \\
Contexto & Operación normal, navegador con lector de pantalla activo o sin ratón conectado. \\
Artefacto & Aplicación web. \\
Respuesta & Todo control es alcanzable con teclado y muestra el foco de forma visible. Toda información transmitida por color va acompañada de texto o de forma. Los mensajes de error se anuncian al lector de pantalla y no solo cambian de color. \\
Medida de respuesta & Conformidad con WCAG 2.2 nivel AA en las pantallas críticas. Cero incidencias de contraste por debajo de 4.5:1 en texto normal y 3:1 en texto grande. Objetivos interactivos de al menos 24 por 24 píxeles. Recorrido completo con teclado sin trampas de foco. \\
\end{ficha}

\noindent\emph{Consideraciones.} Un sistema público de salud que excluye a una parte de sus usuarios no cumple su función. La definición de terminado exige desde el Sprint 1 la verificación automática de accesibilidad en cada historia; este escenario es la auditoría formal, con lector de pantalla y recorrido manual, que la herramienta automática no sustituye. Las pantallas críticas son el registro anónimo, la consulta de inventario y la disposición final de una unidad. Los indicadores de estado del inventario no pueden distinguirse solo por color. La revisión combina herramienta automática y recorrido manual: la herramienta sola detecta contraste y etiquetas, pero no trampas de foco ni orden de lectura incoherente.

\begin{ficha}{EC-19}{Protección ante el desecho accidental de una unidad}{Protección contra errores · IP 3,8 · Alta}
Fuente del estímulo & Operador de banco (U3) en jornada operativa, con prisa o fatiga. \\
Estímulo & Ejecuta una de las tres acciones irreversibles del sistema: la disposición final de una unidad, la decisión sobre una transferencia o el despacho de una unidad. \\
Contexto & Operación normal, varias unidades en pantalla y una lista larga donde es fácil equivocarse de fila. \\
Artefacto & Aplicación web, Servicio de Donación. \\
Respuesta & La acción exige una confirmación explícita que identifica la unidad concreta sobre la que se va a actuar. No se dispara con una sola pulsación ni queda contigua a una acción frecuente e inofensiva. \\
Medida de respuesta & El 100\,\% de las acciones irreversibles exige una confirmación deliberada, separada del control que la inicia, que muestra el identificador de la unidad. Cero acciones irreversibles ejecutables con un solo clic. El 100\,\% queda registrado en la bitácora con su responsable. \\
\end{ficha}

\noindent\emph{Consideraciones.} La confirmación debe mostrar \textbf{qué} unidad, no preguntar «¿está seguro?»: un mensaje genérico no evita el error de fila, que es el error real que se quiere prevenir. La acción destructiva no se coloca junto a la más usada de la pantalla. Este escenario protege contra el error del usuario legítimo, mientras que EC-12 protege contra la operación inválida aunque el usuario la confirme.

\begin{ficha}{EC-28}{Identificación del propósito por un donante nuevo}{Reconocibilidad · IP 2,6 · Media}
Fuente del estímulo & Persona que llega a la aplicación web por primera vez, sin referencia previa del sistema. \\
Estímulo & Abre la página inicial y debe decidir en pocos segundos si el sistema sirve para lo que necesita. \\
Contexto & Sesión nueva, sin cuenta, en escritorio o en móvil. \\
Artefacto & Aplicación web: página inicial y punto de entrada al registro. \\
Respuesta & La página declara sin ambigüedad qué es el sistema, para quién es y cuál es la acción principal disponible, sin exigir lectura de texto extenso ni conocimiento previo del dominio. \\
Medida de respuesta & En prueba con cinco personas ajenas al equipo, al menos cuatro describen correctamente el propósito del sistema y localizan la acción de donar en menos de 30 segundos. Cero terminología técnica del dominio en la pantalla inicial sin explicación adyacente. \\
\end{ficha}

\noindent\emph{Consideraciones.} La reconocibilidad se prueba con personas que no conocen el proyecto; quien lo construyó no puede juzgar si se entiende. El vocabulario del dominio ---tamizaje, componente, unidad--- es correcto entre profesionales y opaco para un donante, y la pantalla inicial se dirige a un donante.

\begin{ficha}{EC-29}{Primer registro de donación por un operador nuevo}{Aprendizaje · IP 2,3 · Baja}
Fuente del estímulo & Cuatro personas ajenas al equipo, sin experiencia previa en el sistema, que reciben por escrito el procedimiento clínico del registro de donación en sustitución de la formación de un operador de banco (U3). \\
Estímulo & Debe registrar su primera donación completa, generar sus unidades y consultarlas en el inventario. \\
Contexto & Ambiente de QA con datos sintéticos, sin acompañamiento del equipo técnico. \\
Artefacto & Aplicación web, módulos M1, M3 y M4. \\
Respuesta & El operador completa el flujo apoyándose únicamente en la propia interfaz y en el material de ayuda incorporado. \\
Medida de respuesta & Al menos tres de las cuatro personas completan el flujo sin asistencia en 10 minutos o menos, y la segunda ejecución del mismo flujo no supera los 3 minutos. Cero pasos que requieran consultar documentación externa al sistema, salvo el procedimiento clínico entregado. \\
\end{ficha}

\noindent\emph{Consideraciones.} La medida contempla dos ejecuciones deliberadamente: el aprendizaje se manifiesta en la diferencia entre la primera y la segunda, no en el tiempo absoluto de la primera. Un sistema en el que la segunda ejecución cuesta lo mismo que la primera no se está aprendiendo, se está tolerando.

\begin{ficha}{EC-30}{Retorno del donante tras recuperar la elegibilidad}{Compromiso del usuario · IP 2,5 · Media}
Fuente del estímulo & Donante registrado (U2) que recupera su elegibilidad tras el periodo mínimo entre donaciones. \\
Estímulo & El sistema detecta la recuperación de elegibilidad y dispone de campañas publicadas en la jurisdicción del donante. \\
Contexto & Operación normal, donante con consentimiento de avisos otorgado y vigente. \\
Artefacto & M1 (elegibilidad y reconocimientos) y bandeja de salida en el Servicio de Donación, M2 en el Servicio de Campañas, trabajador de Notificaciones, aplicación web. \\
Respuesta & El sistema notifica la recuperación de elegibilidad, muestra el progreso de reconocimiento alcanzado y ofrece las campañas disponibles como acción inmediata, sin ofrecer contraprestación alguna. \\
Medida de respuesta & El 100\,\% de los donantes elegibles con consentimiento vigente recibe la notificación dentro de las 24 horas siguientes a recuperar la elegibilidad. La notificación contiene al menos una campaña accionable cuando exista en su jurisdicción y el Servicio de Campañas responda; si no responde, la notificación se entrega en el mismo plazo sin campaña, conforme a la degradación de EC-08. Cero notificaciones a donantes sin consentimiento de avisos. \\
\end{ficha}

\noindent\emph{Consideraciones.} Este escenario es la expresión verificable del objetivo de negocio del sistema: convertir donante ocasional en habitual sin incentivos económicos. El mecanismo es la oportunidad y el reconocimiento simbólico, y por eso la medida exige que la notificación llegue con una campaña accionable: avisar de que se recuperó la elegibilidad sin ofrecer dónde donar traslada el trabajo al donante.

\begin{ficha}{EC-31}{Recuperación ante un rechazo de operación}{Asistencia al usuario · IP 2,1 · Baja}
Fuente del estímulo & Usuario de cualquier perfil autenticado que recibe un rechazo del sistema. \\
Estímulo & Su operación es rechazada por regla de negocio, por conflicto de estado o por falta de alcance sobre el recurso. \\
Contexto & Operación normal. \\
Artefacto & Aplicación web, API Gateway, servicios. \\
Respuesta & El sistema explica qué ocurrió en términos del dominio y qué puede hacer el usuario a continuación, sin exponer detalle técnico, sin culpar al usuario y sin revelar información fuera de su alcance. \\
Medida de respuesta & El 100\,\% de los rechazos presenta un mensaje que nombra la causa en lenguaje de dominio y al menos una acción siguiente posible, con una única excepción declarada: la respuesta \emph{no encontrado} del catálogo de errores del DD, deliberadamente indistinguible del recurso que existe fuera del alcance del usuario, que no nombra causa ni ofrece acción, conforme a EC-01. Cero mensajes que expongan trazas, nombres de componente o códigos internos. Cero mensajes que revelen la existencia de recursos fuera de la jurisdicción del usuario. \\
\end{ficha}

\noindent\emph{Consideraciones.} La asistencia al usuario tiene aquí una restricción que no tiene en otros sistemas: el mensaje más útil sería el más explícito, y el más explícito puede violar EC-01 o EC-02. Un rechazo de despacho no puede decir por qué la unidad no es apta, y un rechazo por jurisdicción no puede confirmar que el recurso exista. La redacción de estos mensajes se revisa con el mismo criterio que un campo del modelo.

\begin{ficha}{EC-32}{Legibilidad del estado de una unidad sin formación previa}{Autodescripción · IP 2,4 · Baja}
Fuente del estímulo & Auditor (U7) o coordinador territorial (U5) que consulta unidades sin conocer la máquina de estados. \\
Estímulo & Abre el historial de trazabilidad de una unidad y debe interpretar en qué situación se encuentra y qué le ocurrió. \\
Contexto & Operación normal, usuario dentro de su jurisdicción. \\
Artefacto & Aplicación web: vista de trazabilidad de la unidad. \\
Respuesta & Cada estado y cada transición se presentan con una denominación comprensible y con la fecha y el responsable, de modo que la secuencia se entienda sin consultar documentación del sistema. \\
Medida de respuesta & Cero códigos internos de estado visibles en la interfaz. El 100\,\% de los estados muestra denominación en lenguaje natural. En prueba con tres personas ajenas al equipo, las tres reconstruyen correctamente la historia de una unidad sin explicación previa. \\
\end{ficha}

\noindent\emph{Consideraciones.} La autodescripción es la condición que hace utilizable la bitácora: el auditor es el usuario que menos contexto técnico tiene y el que más necesita interpretar la secuencia completa. La medida de EC-05 exige que reconstruya la cadena en menos de cinco minutos, y ese objetivo depende de que los estados sean legibles. Este escenario se verifica sobre la vista de trazabilidad del perfil auditor, que se construye en el Sprint 3 sobre la consulta compuesta del Servicio Institucional (ADR-016).

\subsubsection{Mantenibilidad}

\begin{ficha}{EC-34}{Aislamiento del cambio dentro de un contexto}{Modularidad · IP 2,6 · Media}
Fuente del estímulo & Equipo de desarrollo que modifica una capacidad interna de un módulo. \\
Estímulo & Cambia la lógica de cálculo de alertas de inventario, que pertenece por completo a M4. \\
Contexto & Sistema en operación, los cinco servicios desplegados con la versión vigente de sus contratos. \\
Artefacto & Servicio de Donación (M4), demás servicios, contratos entre servicios. \\
Respuesta & El cambio queda contenido dentro del módulo propietario. No obliga a modificar otro módulo, ni los contratos, ni otro servicio, ni a desplegarlos. \\
Medida de respuesta & En el Sprint 1, sobre el modelo: cero entidades propiedad de M4 referenciadas desde otro módulo, cero operaciones del contrato entre servicios afectadas por el cálculo de alertas, y la regla de alertas residente en una sola entidad del modelo lógico. En el Sprint 3, sobre el código: cero archivos modificados fuera de M4, cero cambios en los contratos OpenAPI, cero despliegues de otros servicios requeridos, y las pruebas de los demás módulos pasando sin modificación. \\
\end{ficha}

\noindent\emph{Consideraciones.} La modularidad se mide por lo que \emph{no} cambia. Si un cambio interno de M4 obliga a tocar M3 o el contrato, la frontera entre módulos está mal trazada y el problema es estructural, no de implementación. El cálculo de alertas se elige como sonda porque es lógica de negocio con probabilidad real de cambiar y sin ninguna razón legítima para atravesar la frontera.

\begin{ficha}{EC-35}{Reutilización del componente de interfaz}{Reusabilidad · IP 2,0 · Baja}
Fuente del estímulo & Equipo de desarrollo frontend que construye una pantalla nueva. \\
Estímulo & Necesita un componente ya existente ---tabla paginada, diálogo de confirmación, indicador de estado--- en un módulo distinto de aquel donde se creó. \\
Contexto & Aplicación web en construcción, con la capa de componentes compartidos establecida. \\
Artefacto & Capa de componentes compartidos de la aplicación web. \\
Respuesta & El componente se consume desde la capa compartida sin copiarlo, sin bifurcarlo y sin añadirle lógica propia del módulo consumidor. \\
Medida de respuesta & Cero componentes duplicados entre módulos. El 100\,\% de los diálogos de confirmación y de las tablas paginadas proviene de la capa compartida. Ningún componente compartido contiene reglas de negocio de un módulo concreto. \\
\end{ficha}

\noindent\emph{Consideraciones.} La consecuencia práctica de este escenario es la consistencia visible: si cada módulo escribe su propio diálogo de confirmación, la protección contra errores de EC-19 se cumple de forma desigual y el usuario aprende comportamientos distintos para la misma acción. La reusabilidad aquí no es economía de código sino uniformidad de una garantía de seguridad.

\begin{ficha}{EC-20}{Diagnóstico de un fallo en el ambiente de QA}{Analizabilidad · IP 2,8 · Media}
Fuente del estímulo & Desarrollador o responsable de infraestructura, ante un error reportado por un compañero o por el cliente. \\
Estímulo & Una petición falló en QA y hay que determinar en qué componente y por qué, sin poder reproducir el caso. \\
Contexto & Sistema desplegado, sin acceso a la máquina de quien reportó el fallo. \\
Artefacto & Registros del proxy de borde, del gateway y de los cinco servicios, Prometheus, Grafana. \\
Respuesta & La petición se sigue de extremo a extremo mediante un identificador de correlación que aparece en el registro de todos los componentes que la atendieron. \\
Medida de respuesta & El 100\,\% de las peticiones que cruzan el gateway lleva identificador de correlación propagado a los servicios. El componente donde ocurrió el fallo se identifica en 15 minutos o menos sin reproducir el error. Ningún registro contiene datos personales, causa clínica ni contenido de token. \\
\end{ficha}

\noindent\emph{Consideraciones.} En una arquitectura distribuida, un registro por componente sin identificador común obliga a adivinar qué línea de un servicio corresponde a qué línea de otro. El identificador se genera en el gateway si el cliente no lo trae y se propaga en la cabecera de toda llamada entre servicios. Los registros son parte de la superficie de exposición de datos: un volcado de la petición completa filtra exactamente lo que EC-02 y EC-03 protegen.

\begin{ficha}{EC-21}{Cambio de la regla de elegibilidad}{Modificabilidad · IP 2,4 · Baja}
Fuente del estímulo & El cliente, o un cambio en la normativa de captación. \\
Estímulo & Cambia el periodo mínimo entre donaciones, o el criterio que determina si un donante es elegible. \\
Contexto & Sistema en operación, con donantes ya registrados y con historial. \\
Artefacto & Servicio de Donación: dominio de donantes y elegibilidad. \\
Respuesta & El cambio se aplica en un único punto del modelo de dominio. No obliga a modificar la aplicación web, ningún otro servicio ni ningún contrato entre servicios. \\
Medida de respuesta & El cambio afecta a dos archivos o menos del Servicio de Donación. Cero cambios en los demás servicios y cero en los contratos. Las pruebas existentes siguen pasando y la regla nueva queda cubierta por prueba propia. Esfuerzo estimado de un punto de historia o menos, convertido a horas con la equivalencia que publica cada sprint el Documento de Herramientas. \\
\end{ficha}

\noindent\emph{Consideraciones.} La regla de elegibilidad es la que más probabilidades tiene de cambiar durante la vida del sistema, porque depende de normativa externa; por eso es la que se elige para medir modificabilidad. La medida en número de archivos es un indicador de acoplamiento, no un objetivo de estilo: si el cambio se dispersa, la regla está duplicada y ese es el defecto que el escenario detecta. La interfaz muestra el resultado del cálculo, nunca lo replica.

\begin{ficha}{EC-36}{Verificación exhaustiva de la máquina de estados}{Capacidad de ser probado · IP 3,5 · Alta}
Fuente del estímulo & Equipo de QA que debe demostrar el cumplimiento de EC-04 sin recurrir al muestreo. \\
Estímulo & Necesita ejercitar todas las combinaciones de estado origen y estado destino de la unidad, válidas e inválidas. \\
Contexto & Ambiente de integración continua, sin interfaz gráfica y sin intervención manual. \\
Artefacto & Máquina de estados y catálogo de transiciones del Servicio de Donación, suite de pruebas automatizadas. \\
Respuesta & El sistema expone la transición como operación invocable y verificable de forma independiente, de modo que la matriz completa se pueda recorrer por prueba automatizada. \\
Medida de respuesta & Las 81 combinaciones de la matriz de 9 por 9 estados quedan cubiertas por prueba automatizada: 13 aceptadas y 68 rechazadas. El alta de la unidad se prueba como caso propio, por no ser un par de estados: 82 casos en total, conforme al DD. La suite se ejecuta en menos de 2 minutos y sin base de datos externa a la prueba. Cero casos que requieran preparación manual de datos. \\
\end{ficha}

\noindent\emph{Consideraciones.} La capacidad de ser probado es la condición de posibilidad de EC-04: sin transiciones invocables de forma aislada, la única verificación disponible sería el recorrido manual por la interfaz, que no cubre 81 casos en ningún presupuesto razonable. Que la matriz de transiciones viva como dato de catálogo y no como condiciones repartidas por el código es lo que permite generar los casos en lugar de escribirlos uno a uno.

\subsubsection{Compatibilidad}

\begin{ficha}{EC-22}{Evolución del contrato entre servicios}{Interoperabilidad · IP 3,4 · Media}
Fuente del estímulo & El equipo de cualquiera de los servicios. \\
Estímulo & Publica una versión nueva del contrato que añade un campo a una respuesta existente. \\
Contexto & Los servicios consumidores y la aplicación web están en ejecución con la versión anterior del contrato. \\
Artefacto & Contratos OpenAPI versionados de los cinco servicios, internos y publicados, y aplicación web. \\
Respuesta & El consumidor sigue operando sin cambio alguno. Un cambio que sí rompa la compatibilidad obliga a versión mayor y a un periodo en el que ambas versiones conviven. \\
Medida de respuesta & Cero fallos del consumidor ante un cambio compatible. El 100\,\% de los cambios incompatibles se publica con incremento de versión mayor, conforme a versionado semántico. El contrato vive versionado en el repositorio y la verificación de integración continua lo comprueba en cada Pull Request. Cero llamadas entre servicios fuera del contrato. \\
\end{ficha}

\noindent\emph{Consideraciones.} La restricción de fondo es RD-01: los servicios no comparten base de datos ni se llaman por fuera del contrato. El contrato es independiente del lenguaje porque los servicios están escritos en Java y en .NET. Los contratos internos, con prefijo \texttt{/internal/}, siguen las mismas reglas de versionado aunque el gateway no los enrute. El consumidor ignora los campos que no conoce; un cliente que falla ante un campo desconocido convierte cualquier adición en un cambio incompatible.

\begin{ficha}{EC-27}{Convivencia de los contenedores del sistema en un mismo anfitrión}{Coexistencia · IP 2,0 · Baja}
Fuente del estímulo & Integrante del equipo que levanta el sistema completo en su máquina de trabajo, con otras aplicaciones en ejecución. \\
Estímulo & Los doce contenedores de Desarrollo arrancan y operan simultáneamente sobre un anfitrión compartido con el entorno de desarrollo del usuario. \\
Contexto & Máquina de escritorio con recursos limitados, sin dedicación exclusiva al proyecto. \\
Artefacto & Composición de contenedores de Desarrollo: proxy de borde, aplicación web, gateway, cinco servicios y cuatro bases de datos. \\
Respuesta & Los doce contenedores conviven sin competir por el mismo puerto, sin sobrescribir volúmenes ajenos y sin agotar los recursos del anfitrión hasta impedir su uso normal. \\
Medida de respuesta & Cero conflictos de puerto con servicios habituales de desarrollo: solo el proxy de borde publica puertos al anfitrión. Cero volúmenes compartidos entre las cuatro bases. Consumo agregado en reposo por debajo de 4~GB de memoria con los doce contenedores. El sistema convive con un entorno de desarrollo abierto sin que el anfitrión deje de responder. \\
\end{ficha}

\noindent\emph{Consideraciones.} La coexistencia importa aquí por una razón de proyecto: el equipo trabaja en máquinas personales sin dedicación exclusiva, y un sistema que exige apagar todo lo demás para arrancar se ejecuta menos veces de las necesarias. La ausencia de volúmenes compartidos entre las bases es además la comprobación física de la separación que exigen ADR-005 y ADR-014.

\subsubsection{Flexibilidad}

\begin{ficha}{EC-23}{Puesta en marcha del sistema completo}{Instalabilidad · IP 3,3 · Media}
Fuente del estímulo & Integrante nuevo del equipo, o el cliente al evaluar la entrega. \\
Estímulo & Clona los repositorios en una máquina limpia y quiere levantar el sistema completo para verlo funcionar. \\
Contexto & Máquina sin ninguna dependencia del proyecto instalada: sin Java, sin .NET, sin Node y sin PostgreSQL. Solo Docker. \\
Artefacto & Definición de contenedores y de la composición del sistema. \\
Respuesta & Un único comando levanta los doce contenedores de Desarrollo. Los cuatro trabajos de migración crean el esquema y cargan los datos sintéticos, y el sistema queda operativo sin ningún paso manual adicional. \\
Medida de respuesta & Un comando. Diez minutos o menos hasta que responden correctamente todos los sondeos de salud: los doce contenedores de Desarrollo desde el Sprint 2, después de que los cuatro trabajos de migración terminan con éxito, y los catorce de QA desde el Sprint 4, cuando se incorporan las dos piezas de observabilidad conforme al plan de adopción del catálogo de herramientas. Cero pasos manuales fuera de los documentados. Funciona igual en Windows, macOS y Linux. \\
\end{ficha}

\noindent\emph{Consideraciones.} Es la comprobación operativa de RD-03: que cada componente tenga imagen propia no basta si levantar el conjunto exige conocimiento que solo tiene quien lo construyó. La prueba la ejecuta alguien que no participó en la configuración, siguiendo únicamente el README. Los datos sintéticos se cargan como parte del arranque y son visiblemente ficticios.

\begin{ficha}{EC-24}{Extensión a otro tipo de donación}{Adaptabilidad · IP 2,3 · Baja}
Fuente del estímulo & El cliente, en una fase posterior del producto. \\
Estímulo & Pide administrar un tipo de donación distinto de la sangre, con su propio ciclo de vida y su propia caducidad. \\
Contexto & Sistema en operación con el dominio de sangre ya construido. \\
Artefacto & Modelo de dominio del Servicio de Donación, esquema de base de datos, inventario y trazabilidad. \\
Respuesta & El tipo nuevo se incorpora como dato del modelo. La trazabilidad, el inventario y la disposición final siguen aplicando sin duplicar código ni crear una estructura paralela. \\
Medida de respuesta & Cero cambios estructurales en las tablas de trazabilidad y de inventario. El tipo nuevo se agrega por registro en catálogo o por configuración, no creando tablas propias. Ninguna funcionalidad existente deja de operar. \\
\end{ficha}

\noindent\emph{Consideraciones.} La funcionalidad correspondiente está clasificada como \emph{Won't} para la primera versión. El escenario no la compromete: verifica que el modelo de datos no la impida, que es lo que se decide en el Sprint 1 y ya no se puede revertir barato después. La señal de que el escenario falla es un modelo donde «sangre» aparece en el nombre de las tablas de trazabilidad e inventario, en lugar de como valor de un atributo de tipo.

\begin{ficha}{EC-37}{Sobrecarga de cuatro veces la carga nominal}{Escalabilidad · IP 2,9 · Media}
Fuente del estímulo & Sobrecarga deliberada de \textbf{200 usuarios concurrentes durante 5 minutos}, cuatro veces la capacidad nominal declarada de 50. \\
Estímulo & La carga excede la capacidad para la que el sistema fue dimensionado, sobre las operaciones de consulta de inventario y registro de donación. \\
Contexto & Pila de medición con una sola instancia de cada servicio y los límites de recurso asignados. \\
Artefacto & API Gateway, Servicio de Donación y su base de datos. \\
Respuesta & El sistema degrada de forma controlada: rechaza el exceso con un código explícito antes de agotar recursos, no pierde peticiones en silencio y recupera su comportamiento nominal al cesar la sobrecarga. Levantar una segunda instancia del Servicio de Donación detrás del gateway, con la orden de escalado de réplicas de la herramienta de composición, restablece el desempeño sin modificar código. \\
Medida de respuesta & Cero peticiones perdidas sin respuesta durante la sobrecarga: toda petición se atiende o se rechaza con código 429 o 503. Cero contenedores caídos por agotamiento de recursos. El sistema vuelve al percentil 95 nominal de 3 segundos dentro de los 2 minutos siguientes al cese de la sobrecarga. Con una segunda instancia del Servicio de Donación, la tasa de error desciende por debajo del 1\,\% sin cambio de código ni de imagen. \\
\end{ficha}

\noindent\emph{Consideraciones.} El objeto del escenario no es demostrar que el sistema soporta 200 usuarios ---no está dimensionado para ello--- sino cómo se comporta cuando no puede. Un sistema que ante la sobrecarga rechaza explícitamente es predecible; uno que agota conexiones y cae arrastra a los servicios que lo consultan y convierte un pico en una caída completa. La escalabilidad horizontal sin cambio de código es una consecuencia directa de la arquitectura distribuida y de la ausencia de estado en el servicio, y conviene demostrarla aunque la capacidad nominal no lo exija.

\begin{ficha}{EC-38}{Sustitución del API Gateway}{Reemplazabilidad · IP 1,4 · Baja}
Fuente del estímulo & Equipo de arquitectura, ante una limitación del componente elegido o el fin de su soporte. \\
Estímulo & Se evalúa sustituir el API Gateway ante el fin de soporte del componente elegido o ante una limitación que impida cumplir RD-08. \\
Contexto & Sistema desplegado y operativo, con los cinco servicios en ejecución. \\
Artefacto & API Gateway, definición de rutas, aplicación web, servicios. \\
Respuesta & La sustitución afecta únicamente a la configuración del borde. Los servicios y la aplicación web no cambian, porque no dependen de ninguna capacidad exclusiva del producto sustituido. \\
Medida de respuesta & Cero cambios en el código de los servicios y de la aplicación web. Las rutas, la verificación de sesión y los dos umbrales del límite de tasa se expresan al 100\,\% como configuración portable y no como extensión propietaria, verificado por inspección. El esfuerzo de sustitución, estimado sobre esa configuración, no supera una historia de trabajo. La verificación no requiere ejecutar la sustitución, porque el catálogo no contempla hoy ninguna alternativa en anillo \emph{Adoptar} o \emph{Probar}. \\
\end{ficha}

\noindent\emph{Consideraciones.} Es el escenario de menor índice del catálogo y su valor es preventivo: la doble verificación de sesión de EC-03 ---en el gateway y en cada servicio--- ya garantiza que los servicios no dependan del borde para su seguridad, y este escenario comprueba que esa independencia se conserve en el resto de responsabilidades del gateway. Su verificación es de inspección y no requiere ejecutar la sustitución.

% =====================================================================
%  5. TRADE-OFF Y ADR
% =====================================================================
\section{Compromisos y decisiones arquitectónicas}
\label{sec:tradeoff}

\subsection{Cómo se decide en este proyecto}

Toda decisión estructural se delibera en la \textbf{Mesa de Arquitectura} y se registra en un ADR antes de implementarse. La regla tiene una consecuencia práctica que conviene declarar: una tecnología en ejecución sin ADR asociado es deuda, no una decisión. El \emph{killer} K-09 existe precisamente para eso.

Un ADR de este proyecto declara, como mínimo: el contexto que obliga a decidir, las opciones consideradas con sus consecuencias, el compromiso evaluado, la decisión, su justificación, y las consecuencias aceptadas y habilitadas. El ADR no oculta lo que se pierde: si una decisión no tiene costo, probablemente no era una decisión.

\subsection{Compromisos evaluados}

Un compromiso arquitectónico aparece cuando dos atributos de calidad tiran en direcciones opuestas y no existe una opción que mejore ambos. Se registran siete.

\begin{longtable}{L{1.2cm} L{3.4cm} L{4.6cm} L{4.7cm}}
\caption{Compromisos arquitectónicos evaluados y criterio de resolución.}\label{tab:tradeoffs}\\
\toprule
\textbf{ID} & \textbf{Tensión} & \textbf{Qué se gana y qué se pierde} & \textbf{Resolución y criterio} \\
\midrule
\endfirsthead
\toprule
\textbf{ID} & \textbf{Tensión} & \textbf{Qué se gana y qué se pierde} & \textbf{Resolución y criterio} \\
\midrule
\endhead
\midrule
\contlinea{4} \\
\endfoot
\bottomrule
\endlastfoot
TO-01 & Independencia de despliegue \emph{frente a} consistencia de datos y costo operativo & Microservicios de grano fino dan independencia real por módulo, a cambio de transacciones distribuidas con compensación y de multiplicar el costo de despliegue, observabilidad y contratos. & \textbf{Servicios de grano grueso por contexto acotado.} Criterio: en un dominio donde una unidad no puede quedar en estado ambiguo, y con un equipo sin experiencia operando sistemas distribuidos, la independencia de despliegue no compensa la pérdida de atomicidad. \textbf{ADR-003}. \\
TO-02 & Riesgo de aprendizaje \emph{frente a} criticidad del contexto & Situar .NET en el contexto crítico traslada la curva de aprendizaje a los requerimientos con menos margen de error; situarlo en un contexto marginal arriesga que ese componente nunca se construya e incumplir RD-02 de hecho. & \textbf{Java en el contexto del ciclo de vida de la sangre y en su trabajador de notificaciones; .NET en los contextos administrativos de identidad, territorio y campañas, y en el gateway.} Criterio: el riesgo se coloca donde el error cuesta menos, sobre contextos con volumen suficiente para que los componentes existan con certeza. \textbf{ADR-014}, que conserva el criterio de ADR-004. \\
TO-03 & Flexibilidad de esquema \emph{frente a} garantías de consistencia & Una base documental favorecería la extensibilidad a otros tipos de donación, a cambio de perder atomicidad e integridad referencial en un dominio clínico. & \textbf{Base relacional, una por servicio.} Criterio: el beneficio que justificaría NoSQL ---extensibilidad--- se obtiene igual modelando el tipo de donación como atributo, mientras que su costo no tiene compensación equivalente. \textbf{ADR-005}. \\
TO-04 & Velocidad inicial \emph{frente a} control de la identidad visual y tamaño de la superficie mantenida & Una librería de componentes completa entrega mucho hecho y cobra después en re-tematización; su paleta semántica usa rojo saturado para error, que es exactamente el registro que la marca decidió evitar. & \textbf{Tailwind con primitivos accesibles y capa de componentes propia.} Criterio: la accesibilidad de un diálogo modal es fácil de implementar mal sin que ninguna prueba manual lo note, y se delega; la identidad visual no se negocia. \textbf{ADR-006}. \\
TO-05 & Seguridad \emph{frente a} facilidad de uso en el punto de entrada público & EC-17 exige registro anónimo en tres pasos; EC-06 exige contener el abuso automatizado del mismo \emph{endpoint}. Añadir verificación al registro satisface uno y rompe el otro. & \textbf{Límite de tasa por origen en el API Gateway, sin paso adicional en el registro.} Criterio: el control se mueve de componente en lugar de ceder en el atributo. El límite de tres pasos es duro. Los umbrales se fijan en \textbf{ADR-007}. \\
TO-06 & Desacople de las capacidades transversales \emph{frente a} costo operativo y margen de puesta en marcha & Con la identidad dentro del Servicio Institucional y las notificaciones dentro del de Donación, el inicio de sesión dependía de la administración territorial y los reintentos compartían proceso con el flujo transaccional crítico. Extraerlos, junto con las campañas, sube los contenedores de ocho a doce en Desarrollo y a catorce en QA, y los repositorios de seis a nueve. & \textbf{Cinco servicios de grano grueso, con el flujo núcleo intacto en el Servicio de Donación.} Criterio: ninguna capacidad transversal vive dentro de un contexto funcional, y ninguna invariante local se convierte en distribuida. \textbf{ADR-014}. \\
TO-07 & Disponibilidad \emph{frente a} revocación inmediata de la sesión & Validar cada petición contra un registro central de sesiones daría revocación inmediata a cambio de acoplar la disponibilidad de todo el sistema a la del emisor. Un token sin estado desacopla la disponibilidad, pero no se revoca antes de vencer. & \textbf{Token de acceso de quince minutos sin estado, con sesión de renovación de ocho horas con estado en el Servicio de Identidad.} Criterio: el estado existe, pero solo se consulta al renovar; la ventana de revocación queda acotada a quince minutos. \textbf{ADR-008}. \\
\end{longtable}

\subsection{Puntos de sensibilidad y de compromiso}

Un punto de sensibilidad es un elemento cuya modificación altera de forma apreciable la respuesta de un atributo. Un punto de compromiso es un punto de sensibilidad que afecta a dos o más atributos en direcciones opuestas: es donde se decide, no donde se optimiza.

\begin{longtable}{L{1.2cm} L{3.6cm} L{2.0cm} L{7.4cm}}
\caption{Puntos de sensibilidad y de compromiso de la arquitectura.}\label{tab:sensibilidad}\\
\toprule
\textbf{ID} & \textbf{Elemento} & \textbf{Tipo} & \textbf{Atributos afectados y consecuencia} \\
\midrule
\endfirsthead
\toprule
\textbf{ID} & \textbf{Elemento} & \textbf{Tipo} & \textbf{Atributos afectados y consecuencia} \\
\midrule
\endhead
\midrule
\contlinea{4} \\
\endfoot
\bottomrule
\endlastfoot
PS-01 & API Gateway & Compromiso & Concentra autenticación, compuerta por ruta y rol, límite de tasa en dos escalas, identificador de correlación y entrega de las denegaciones a la bitácora: sostiene EC-01, EC-03, EC-06 y EC-20. A cambio es punto común de fallo de todo acceso externo (EC-07) y punto único de latencia añadida (EC-15). Es el componente cuya caída cuesta más. \\
PS-02 & Fronteras entre servicios & Compromiso & Colocadas donde están, mantienen la atomicidad del flujo núcleo dentro del Servicio de Donación (EC-04, EC-10) y obligan a que la evaluación del escalamiento, los avisos de campaña y la consulta del auditor consulten otros contextos por red (EC-08). Desde ADR-013 el escalamiento ocurre en un proceso programado y no en una petición interactiva, de modo que deja de competir con el objetivo de tiempo de respuesta de EC-15. Mover una frontera para reducir llamadas rompería la atomicidad o reintroduciría una capacidad transversal dentro de un contexto funcional. \\
PS-03 & Emisor del token en el Servicio de Identidad & Sensibilidad & La jurisdicción viaja en el token y ningún servicio consulta a Identidad para autorizar, lo que favorece EC-15. El inicio y la renovación de sesión dependen solo de Identidad: si cae, no se emiten sesiones nuevas ni se renuevan, pero las peticiones con token vigente se siguen atendiendo durante un máximo de quince minutos. Ya no dependen de la administración territorial. \\
PS-04 & Ausencia de la causa clínica en el modelo & Sensibilidad & Hace inviolable EC-02 por construcción. El costo es que el sistema no puede producir estadística clínica de motivos de rechazo, capacidad que sí tendría un sistema que capturara el dato. Se acepta: la restricción legal precede a la utilidad analítica. \\
PS-05 & Tipo de donación como valor de atributo & Sensibilidad & Habilita EC-24 sin costo hoy. Introduce una indirección en el modelo que no aporta nada mientras solo exista sangre, y que resulta gratuita solo si se decide antes de la primera migración. \\
PS-06 & Plazo de espera de 3 segundos entre servicios & Compromiso & Protege contra el fallo en cascada (EC-08) y acota el peor caso de EC-15. A cambio, una llamada legítima y lenta se aborta: se prefiere una respuesta degradada explícita a una espera indefinida. \\
PS-07 & Vigencia de quince minutos del token de acceso & Compromiso & Acota la exposición de un token robado y el tiempo en que un cambio de jurisdicción surte efecto (EC-03, RF-22). A cambio, cada sesión renueva su token cuatro veces por hora contra Identidad. Acortarla aumenta la carga sobre Identidad; alargarla amplía la ventana de revocación. \\
\end{longtable}

\subsection{Registro de decisiones arquitectónicas}

\begin{longtable}{L{1.6cm} L{5.6cm} L{3.2cm} L{3.4cm}}
\caption{Registro de decisiones arquitectónicas vigentes.}\label{tab:adrs}\\
\toprule
\textbf{ADR} & \textbf{Decisión} & \textbf{Estado} & \textbf{Próxima revisión} \\
\midrule
\endfirsthead
\toprule
\textbf{ADR} & \textbf{Decisión} & \textbf{Estado} & \textbf{Próxima revisión} \\
\midrule
\endhead
\midrule
\contlinea{4} \\
\endfoot
\bottomrule
\endlastfoot
ADR-001 & Stack tecnológico único en Java & Reemplazada por ADR-004 & --- \\
ADR-002 & Monolito modular & Reemplazada por ADR-003 & --- \\
ADR-003 & Arquitectura distribuida de servicios de grano grueso por contexto acotado & \textbf{Aprobada} & Cierre del Sprint 4 \\
ADR-004 & Reparto de componentes entre Java y .NET por frontera de contexto & Reemplazada por ADR-014 & --- \\
ADR-005 & Base de datos relacional con esquema independiente por servicio & \textbf{Aprobada} & Cierre del Sprint 3 \\
ADR-006 & Sistema de estilos y componentes de la aplicación web & \textbf{Aprobada} & Cierre del Sprint 3 \\
ADR-007 & API Gateway con YARP: compuerta, límite de tasa en dos escalas y registro de denegaciones & \textbf{Aprobada} & Cierre del Sprint 3 \\
ADR-008 & Token de sesión con firma asimétrica emitido por el Servicio de Identidad & \textbf{Aprobada} & Cierre del Sprint 3 \\
ADR-009 & Línea de plataforma Java 25 con Temurin y Spring Boot 4.1 & \textbf{Aprobada} & Cierre del Sprint 3 \\
ADR-010 & Almacenamiento distribuido de la bitácora de auditoría con consulta compuesta & Reemplazada por ADR-016 & --- \\
ADR-011 & Matriz de autorización por grupo de operaciones como especificación normativa del control de acceso & \textbf{Aprobada} & Cierre del Sprint 3 \\
ADR-012 & Clasificación de la extensibilidad del dominio en ISO/IEC 25010:2023 & \textbf{Aprobada} & Cierre del Sprint 2 \\
ADR-013 & Orquestación del escalamiento en dos niveles como proceso del Servicio de Donación & \textbf{Aprobada} & Cierre del Sprint 4 \\
ADR-014 & Descomposición en cinco servicios de grano grueso y topología por zonas de confianza & \textbf{Aprobada} & Cierre del Sprint 3 \\
ADR-015 & Caddy como proxy de borde & \textbf{Aprobada} & Cierre del Sprint 3 \\
ADR-016 & Bitácora de auditoría distribuida en cuatro contextos con consulta compuesta & \textbf{Aprobada} & Cierre del Sprint 3 \\
\end{longtable}

\noindent Los doce ADR vigentes fueron deliberados y aprobados en Mesa de Arquitectura, y la arquitectura descrita en este documento se deriva de ellos. La columna de próxima revisión no expresa provisionalidad: fija el momento en que cada decisión se contrasta con la evidencia que la construcción haya producido, conforme a la política de gobierno de la Sección \ref{sec:tradeoff}. Una decisión aprobada solo se modifica emitiendo un ADR que la reemplace.

ADR-001, ADR-002, ADR-004 y ADR-010 se conservan marcados como reemplazados y no se eliminan. Una decisión retirada en silencio impide auditar por qué cambió.

\subsection{Decisiones estructurales cerradas en el Sprint 2}

Tres elementos de la arquitectura se apoyaban en la V1.0 en el catálogo de herramientas y no en un ADR propio, y la revisión de la topología con el cliente abrió otros tres. Un catálogo registra la consecuencia de una decisión; no la sustituye. Los seis quedaron registrados en la Mesa de Arquitectura del 22 de septiembre de 2026.

\begin{longtable}{L{1.4cm} L{5.4cm} L{7.4cm}}
\caption{Decisiones estructurales cerradas en el Sprint 2 y registro que las sustenta.}\label{tab:pendientes}\\
\toprule
\textbf{ADR} & \textbf{Decisión} & \textbf{Qué fija} \\
\midrule
\endfirsthead
\toprule
\textbf{ADR} & \textbf{Decisión} & \textbf{Qué fija} \\
\midrule
\endhead
\midrule
\contlinea{3} \\
\endfoot
\bottomrule
\endlastfoot
ADR-007 & API Gateway con YARP. Cierra P-01. & Tabla de rutas por servicio, rutas públicas, compuerta derivada de la matriz de ADR-011, rechazo con 503 ante la ausencia de claves públicas, límite de tasa de 60 peticiones por minuto y origen en las rutas públicas y de 300 por minuto y sesión en las protegidas, plazo de 5 segundos hacia el servicio destino e identificador de correlación. \\
ADR-008 & Token de sesión con firma asimétrica. Cierra P-02. & Emisor único en el Servicio de Identidad; RS256 con claves de 3072 bits; ocho reivindicaciones sin dato personal; token de acceso de quince minutos; sesión de renovación de ocho horas con rotación, guardada como resumen; credenciales de servicio; publicación de claves y rotación sin invalidar sesiones. \\
ADR-009 & Línea de plataforma Java 25. Cierra P-03. & Java 25 con distribución Eclipse Temurin y Spring Boot 4.1, con sus ventanas de soporte verificadas en fuente oficial. \\
ADR-014 & Descomposición en cinco servicios. & Servicios, propiedad de los datos, contenedores, zonas de confianza, dieciocho aristas permitidas, bandeja de salida expuesta por contrato, orden de arranque y nueve repositorios. Reemplaza a ADR-004. \\
ADR-015 & Caddy como proxy de borde. & Terminación TLS, certificados por ambiente, cabeceras de seguridad, tamaño máximo de cuerpo y registro de acceso sin credenciales, sin limitación de tasa en el borde. \\
ADR-016 & Bitácora de auditoría en cuatro contextos. & Cuatro series, entrega asíncrona de las denegaciones del gateway al Servicio de Identidad y consulta compuesta del auditor en el Servicio Institucional con degradación explícita. Reemplaza a ADR-010. \\
\end{longtable}

% =====================================================================
%  6. ARQUITECTURA DE ALTO NIVEL
% =====================================================================
\section{Arquitectura de alto nivel}
\label{sec:hld}

\subsection{Estilo arquitectónico}

RedVital adopta una \textbf{arquitectura distribuida de servicios de grano grueso, delimitados por contexto acotado}. No es un monolito modular, prohibido por RD-01, y no es una arquitectura de microservicios de grano fino, que RD-01 tampoco exige y que este proyecto descartó por TO-01.

La frontera entre servicios no se traza por capa técnica ni por conveniencia de implementación, sino por \textbf{coherencia de lenguaje de dominio}. Dentro de un contexto acotado, cada término tiene un único significado; el corte se hace donde el significado cambia. En este sistema el corte natural es evidente: el ciclo de vida de la sangre es transaccional, de alta criticidad y alta volatilidad de estado; la gestión institucional y territorial es administrativa, de baja volatilidad y sin exigencias de atomicidad; las campañas son un contexto de lectura pública masiva con lenguaje propio. A esos tres contextos se suman dos capacidades transversales ---la identidad y la entrega de notificaciones--- que ningún contexto funcional debe alojar. El resultado son cinco servicios de grano grueso (ADR-014).

Tres reglas de frontera gobiernan la estructura y no admiten excepción:

\begin{enumerate}
\item \textbf{Ningún servicio accede a la base de datos de otro.} La única vía es el contrato publicado. Es lo que impide el \emph{killer} K-04.
\item \textbf{Ningún cliente externo alcanza un servicio sin pasar por el proxy de borde y el API Gateway}, y ningún servicio confía en que el gateway ya verificó por él. Es lo que impide K-07.
\item \textbf{Toda llamada entre servicios tiene plazo de espera acotado} y degradación declarada. Es lo que impide K-05.
\end{enumerate}

\subsection{Vista de contexto}

El sistema tiene cuatro clases de actores humanos, un actor no humano ---el tiempo, que dispara el vencimiento de unidades--- y un sistema externo con el que la interoperabilidad está especificada pero fuera de la línea base de la V1.

\begin{figure}[h!]
\centering
\begin{tikzpicture}[
  node distance=0.6cm,
  actor/.style={draw=grisTexto!70, fill=white, rounded corners=2pt, align=center, font=\scriptsize, inner sep=4pt, minimum height=0.95cm, text width=2.5cm},
  sistema/.style={draw=vinotinto, fill=vinotinto!10, rounded corners=3pt, align=center, font=\small\bfseries, inner sep=6pt, minimum height=2.2cm, text width=5.0cm, text=vinotinto},
  externo/.style={draw=ringEvaluar!80, fill=ringEvaluar!8, rounded corners=2pt, align=center, font=\scriptsize, inner sep=4pt, minimum height=0.95cm, text width=2.6cm},
  flecha/.style={-{Stealth}, draw=vinotinto!70, thick},
  flechagris/.style={-{Stealth}, draw=grisTexto!70, thick},
  flechapunt/.style={-{Stealth}, draw=ringEvaluar!80, thick, dashed}
]

\node[sistema] (sys) at (0,0) {REDVITAL\\[2pt]\normalfont\scriptsize Plataforma nacional de gestión\\ del ciclo de vida de la donación};

\node[actor] (donante) at (-6.2,2.4) {\textbf{Donantes}\\ U1 anónimo\\ U2 registrado};
\node[actor] (banco) at (-6.2,-0.3) {\textbf{Banco de sangre}\\ U3 operador\\ U4 administrador};
\node[actor] (territorio) at (6.2,2.4) {\textbf{Autoridad sanitaria}\\ U5 coordinador\\ U6 nacional};
\node[actor] (auditor) at (6.2,-0.3) {\textbf{U7 Auditor}\\ verificación de\\ cumplimiento};
\node[actor] (tiempo) at (-6.2,-3.0) {\textbf{El tiempo}\\ actor no humano};
\node[externo] (sihevi) at (6.2,-3.0) {\textbf{SIHEVI-INS}\\ Instituto Nacional\\ de Salud};

\draw[flechagris] (donante) -- node[above, font=\tiny, text=grisTexto, sloped, pos=0.58] {se registra} (sys);
\draw[flechagris] (banco) -- node[above, font=\tiny, text=grisTexto, pos=0.5] {registra donaciones} (sys);
\draw[flechagris] (sys) -- node[above, font=\tiny, text=grisTexto, sloped, pos=0.42] {supervisa} (territorio);
\draw[flechagris] (sys) -- node[above, font=\tiny, text=grisTexto, pos=0.5] {consulta bitácora} (auditor);
\draw[flechagris] (tiempo) -- node[above, font=\tiny, text=grisTexto, sloped, pos=0.58] {vence unidades} (sys);
\draw[flechapunt] (sys) -- node[above, font=\tiny, text=ringEvaluar, sloped, pos=0.5] {reporte (fuera de la V1)} (sihevi);
\end{tikzpicture}
\caption[Vista de contexto del sistema]{Vista de contexto. El tiempo figura como actor porque RF-18 exige que la unidad vencida salga del inventario sin intervención humana: si no se declara como fuente de estímulo, el proceso programado que lo realiza queda sin justificación en la arquitectura. La integración con SIHEVI-INS está especificada en RNF-07 pero fuera de la línea base de la V1.}
\label{fig:contexto}
\end{figure}

\subsection{Vista de contenedores}

El sistema se compone de cinco servicios, un API Gateway, una aplicación web, un proxy de borde, cuatro bases de datos y, en QA, dos piezas de observabilidad (ADR-014). Son doce contenedores de larga vida en Desarrollo y catorce en QA, más cuatro trabajos de migración que crean el esquema de cada base, cargan su porción del conjunto sintético y terminan. Es la cifra que EC-23 compromete levantar con un solo comando.

Los cinco servicios responden a dos reglas de corte que se suman al contexto acotado. La primera: ninguna capacidad transversal vive dentro de un contexto funcional, y por eso la identidad y la entrega de notificaciones tienen servicio propio. La segunda: ninguna invariante local se convierte en distribuida, y por eso el flujo núcleo ---donante, donación, unidad y su evento--- permanece en un solo servicio y una sola base.

\begin{figure}[h!]
\centering
\begin{tikzpicture}[scale=0.92, every node/.style={transform shape},
  comp/.style={draw=vinotinto!70, fill=vinotinto!7, rounded corners=3pt, align=center, font=\scriptsize, inner sep=3pt, minimum height=1.25cm, text width=2.55cm},
  borde/.style={draw=vinotinto, fill=vinotinto!16, rounded corners=3pt, align=center, font=\scriptsize, inner sep=3pt, minimum height=0.95cm, text width=3.0cm},
  bd/.style={draw=ringEvaluar!85, fill=ringEvaluar!8, rounded corners=3pt, align=center, font=\scriptsize, inner sep=3pt, minimum height=0.85cm, text width=2.55cm},
  obs/.style={draw=ringProbar!85, fill=ringProbar!8, rounded corners=3pt, align=center, font=\scriptsize, inner sep=3pt, minimum height=0.75cm, text width=2.55cm},
  ext/.style={draw=grisTexto!70, fill=white, rounded corners=2pt, align=center, font=\scriptsize, inner sep=3pt, minimum height=0.7cm, text width=2.4cm},
  flecha/.style={-{Stealth}, draw=vinotinto!75, thick},
  flechabd/.style={-{Stealth}, draw=ringEvaluar!85, thick},
  flechaint/.style={-{Stealth}, draw=grisTexto!80, thick},
  flechaobs/.style={-{Stealth}, draw=ringProbar!85, thick, dashed}
]
\node[ext] (nav) at (0,9.3) {Navegador};
\node[borde] (px) at (0,8.0) {\textbf{Proxy de borde}\\ Caddy · TLS · 443};
\node[comp, minimum height=0.95cm] (web) at (-3.4,6.5) {\textbf{Aplicación Web}\\ React · estática};
\node[borde] (gw) at (3.4,6.5) {\textbf{API Gateway}\\ YARP · compuerta · tasa};

\node[comp] (id)  at (-6.4,3.9) {\textbf{Identidad}\\ .NET\\ ST3 emisión · sesión};
\node[comp] (ins) at (-3.2,3.9) {\textbf{Institucional}\\ .NET\\ M6 · M7 · consulta ST2};
\node[comp] (cam) at (0,3.9)    {\textbf{Campañas}\\ .NET\\ M2};
\node[comp] (don) at (3.2,3.9)  {\textbf{Donación}\\ Java\\ M1 · M3 · M4 · M5 · bandeja};
\node[comp] (not) at (6.4,3.9)  {\textbf{Notificaciones}\\ Java · trabajador\\ ST1 · sin base};

\node[bd] (dbi) at (-6.4,0.9) {PostgreSQL\\ \texttt{db\_identidad}};
\node[bd] (dbn) at (-3.2,0.9) {PostgreSQL\\ \texttt{db\_institucional}};
\node[bd] (dbc) at (0,0.9)    {PostgreSQL\\ \texttt{db\_campana}};
\node[bd] (dbd) at (3.2,0.9)  {PostgreSQL\\ \texttt{db\_donacion}};
\node[obs] (prom) at (6.4,1.3) {Prometheus · QA};
\node[obs] (graf) at (6.4,0.2) {Grafana · QA};

\draw[flecha] (nav) -- node[right, font=\tiny, text=vinotinto] {HTTPS} (px);
\draw[flecha] (px) -- (web);
\draw[flecha] (px) -- node[above right, font=\tiny, text=vinotinto] {/api/*} (gw);
\draw[flecha] (gw.south) -- (id.north);
\draw[flecha] (gw.south) -- (ins.north);
\draw[flecha] (gw.south) -- (cam.north);
\draw[flecha] (gw.south) -- (don.north);

\draw[flechaint] (cam.west) -- (ins.east);
\draw[flechaint] (don.west) -- (cam.east);
\draw[flechaint] (not.west) -- node[below, font=\tiny, text=grisTexto, yshift=-2pt] {bandeja} (don.east);
\draw[flechaint, {Stealth}-{Stealth}] ([xshift=-6pt]don.south) .. controls +(0,-1.1) and +(0,-1.1) .. node[below, pos=0.5, font=\tiny, text=grisTexto, fill=white, inner sep=1pt] {instituciones · indicadores · auditoría} ([xshift=6pt]ins.south);

\draw[flechabd] ([xshift=-14pt]id.south) -- ([xshift=-14pt]dbi.north);
\draw[flechabd] ([xshift=-14pt]ins.south) -- ([xshift=-14pt]dbn.north);
\draw[flechabd] ([xshift=-14pt]cam.south) -- ([xshift=-14pt]dbc.north);
\draw[flechabd] ([xshift=14pt]don.south) -- ([xshift=14pt]dbd.north);
\draw[flechaobs] (prom) -- (graf);
\end{tikzpicture}
\caption[Vista de contenedores del sistema]{Vista de contenedores. El proxy de borde es el único componente alcanzable desde fuera; el gateway enruta a los cuatro servicios con rutas publicadas y nunca al trabajador de Notificaciones. Las flechas grises son contratos entre servicios con plazo de espera de 3 segundos; la figura muestra los de negocio, y la Tabla \ref{tab:aristas} enumera las dieciocho aristas permitidas, incluidas la obtención de claves públicas en Identidad por el gateway y los servicios, el canje de tokens de servicio y la consulta de contactos del trabajador. Ninguna base se comunica con otra ni es alcanzable desde un servicio ajeno.}
\label{fig:contenedores}
\end{figure}

\subsection{Responsabilidad de cada componente}

\begin{longtable}{L{3.0cm} L{5.6cm} L{5.6cm}}
\caption{Responsabilidad de cada componente y frontera de lo que no le corresponde.}\label{tab:componentes}\\
\toprule
\textbf{Componente} & \textbf{Responsabilidad} & \textbf{Lo que explícitamente no hace} \\
\midrule
\endfirsthead
\toprule
\textbf{Componente} & \textbf{Responsabilidad} & \textbf{Lo que explícitamente no hace} \\
\midrule
\endhead
\midrule
\contlinea{3} \\
\endfoot
\bottomrule
\endlastfoot
Proxy de borde (Caddy) & Único componente con puerto publicado. Terminación TLS 1.2 y 1.3, redirección a HTTPS, cabeceras de seguridad, tamaño máximo de cuerpo y reparto hacia la aplicación web, el gateway y, en QA, Grafana (ADR-015). & No valida tokens, no limita la tasa y no aplica ninguna regla del dominio. \\
Aplicación Web & Presentación de los siete perfiles. Recorridos, estados de excepción y accesibilidad. Muestra el resultado de las reglas del dominio. & No replica ninguna regla de negocio ---la elegibilidad se calcula en el servicio, nunca en el navegador--- y no es una barrera de seguridad: ocultar una opción del menú no autoriza nada. \\
API Gateway (YARP) & Punto único de entrada a los servicios. Enrutamiento, validación del token, compuerta por ruta y rol, límite de tasa en dos escalas, generación y propagación del identificador de correlación, y entrega de las denegaciones al Servicio de Identidad (ADR-007). & No contiene lógica de dominio ni accede a ninguna base de datos. No enruta nunca rutas \texttt{/internal/}. No es la única barrera de autorización: el filtrado por jurisdicción a nivel de registro vive en cada servicio. \\
Servicio de Identidad (.NET) & Usuarios, roles, asignación de jurisdicción, inicio, renovación y cierre de sesión, emisión de tokens de usuario y de servicio, publicación de claves públicas y serie de auditoría de identidad y de denegaciones del gateway (ADR-008, ADR-016). & No conoce la jerarquía territorial más allá de la ruta que asigna, ni ningún dato del dominio de la sangre. \\
Servicio Institucional (.NET) & Jerarquía nacional, departamental y municipal, registro de instituciones, indicadores agregados de solo lectura y consulta compuesta de la bitácora para el auditor. & No accede a unidades ni a inventario: los consulta por contrato y en solo lectura. No escribe nunca en otro contexto. No emite identidad. \\
Servicio de Campañas (.NET) & Planificación, publicación, cupos, reservas de cupo con expiración y seguimiento de campañas. Guarda con cada campaña el territorio y la ruta territorial de su sede. & No conoce donantes ni notifica: el Servicio de Donación consulta las campañas publicadas y genera los avisos. \\
Servicio de Donación (Java) & Contexto del ciclo de vida de la sangre: donantes, elegibilidad, reconocimientos, donaciones, unidades, trazabilidad, inventario, alertas, transferencias y proceso de escalamiento. Escribe la bitácora de sus entidades y la bandeja de salida de notificaciones. & No conoce la jerarquía territorial, el catálogo de instituciones ni las campañas: los consulta por contrato. No emite identidad ni entrega notificaciones. \\
Trabajador de Notificaciones (Java) & Consume la bandeja de salida del Servicio de Donación por contrato interno, resuelve los contactos institucionales en Identidad y entrega por el canal de salida con reintento de espera creciente. & No tiene base de datos ni rutas publicadas. No lee la base de ningún servicio. \\
PostgreSQL por servicio & Persistencia transaccional del contexto de su servicio, con red propia y roles separados: propietario para migrar y respaldar, y de servicio sin definición de esquema. & Ninguna instancia es alcanzable desde un servicio ajeno. \\
Trabajos de migración & Uno por base: esquema versionado y porción del conjunto sintético. Terminan y no permanecen en ejecución. & No alcanzan ninguna base distinta de la suya. \\
Prometheus y Grafana (QA) & Recolección y presentación de las series que sostienen las medidas de respuesta de EC-07, EC-15, EC-16 y EC-20. & No almacenan datos personales, causa clínica ni contenido de token en ninguna etiqueta de métrica. Prometheus no se expone. \\
\end{longtable}

\subsection{Reparto de módulos y servicios transversales}

\begin{longtable}{L{1.2cm} L{4.2cm} L{3.6cm} L{5.2cm}}
\caption{Servicio propietario de cada módulo y dependencia entre contextos.}\label{tab:repartomodulos}\\
\toprule
\textbf{ID} & \textbf{Módulo} & \textbf{Servicio propietario} & \textbf{Dependencia por contrato} \\
\midrule
\endfirsthead
\toprule
\textbf{ID} & \textbf{Módulo} & \textbf{Servicio propietario} & \textbf{Dependencia por contrato} \\
\midrule
\endhead
\midrule
\contlinea{4} \\
\endfoot
\bottomrule
\endlastfoot
M1 & Gestión de Donantes & Donación (Java) & Consulta a M2, en Campañas, para las campañas disponibles y publicadas. \\
M2 & Gestión de Campañas & Campañas (.NET) & Consulta a M6, en Institucional, la sede, la institución organizadora y su jurisdicción. \\
M3 & Trazabilidad y Ciclo de Vida de la Unidad & Donación (Java) & Consulta a M6 para la institución custodia. \\
M4 & Inventario y Alertas & Donación (Java) & Consulta a M6 para la institución propietaria de las existencias. \\
M5 & Red de Transferencias entre Bancos & Donación (Java) & Consulta a M6 las instituciones cedente y receptora y su departamento, para el escalamiento (ADR-013). \\
M6 & Administración Institucional y Territorial & Institucional (.NET) & Ninguna. Es el contexto de referencia territorial del sistema. \\
M7 & Analítica e Indicadores & Institucional (.NET) & Consulta a M3 y M4, en Donación, para las series que agrega. No escribe en ellos (RI-03). \\
ST1 & Notificaciones & Trabajador de Notificaciones (Java) & Consume la bandeja de salida del Servicio de Donación, alimentada por M1, M2, M4 y M5, y resuelve en Identidad los contactos institucionales (ADR-014). \\
ST2 & Auditoría & Distribuido en cuatro series, con consulta compuesta en Institucional & Cada servicio con base escribe los registros de las entidades que posee; el gateway entrega sus denegaciones a Identidad. La consulta del auditor compone las cuatro series (ADR-016). \\
ST3 & Autenticación y autorización & Identidad (.NET) + Gateway + cada servicio & La emisión reside en Identidad; la compuerta, en el gateway; el filtrado por registro, en cada servicio propietario (ADR-008). \\
\end{longtable}

\noindent\textbf{Sobre ST2.} La bitácora de auditoría es el único elemento del sistema cuyo dueño no es un módulo funcional. Es una consecuencia legítima de su naturaleza ---la escriben todos los componentes--- pero deja al perfil U7, cuyo trabajo \emph{es} la bitácora, sin módulo de referencia. La resolución adoptada en ADR-016 es declararla servicio transversal con \textbf{almacenamiento en cuatro contextos}: cada servicio con base escribe los registros de las entidades que posee, en su propia base y en la misma transacción que la acción que registra. Las denegaciones que produce un servicio quedan en su propia serie. Las que produce el gateway se entregan de forma asíncrona al Servicio de Identidad, porque un token ausente o inválido es un hecho de identidad; el gateway no retrasa el rechazo para registrarlo, y un evento que no puede entregarse se contabiliza en una métrica propia.

\vspace{0.2cm}
\noindent La consulta del auditor se sirve desde el Servicio Institucional, que lee su propia serie, invoca en paralelo las de Donación, Campañas e Identidad con el token del auditor y un plazo de 3 segundos por llamada, y ordena el conjunto por marca temporal en tiempo universal coordinado agrupándolo por identificador de correlación. El motivo de no concentrar el almacenamiento en un solo contexto es de integridad transaccional: la regla que exige insertar el evento de trazabilidad en la misma transacción que actualiza el estado de la unidad solo es imponible si el evento y la unidad residen en la misma base. Ante indisponibilidad de uno o varios contextos, la respuesta declara cuáles faltan y nunca se presenta como completa.

\subsection{Tácticas por atributo de calidad}

Una táctica es una decisión de diseño que produce una respuesta concreta ante un estímulo. La tabla siguiente es el puente entre los escenarios de la Sección \ref{sec:escenarios} y los componentes de esta sección: dice qué se hizo, dónde, y qué escenario lo comprueba.

\begin{longtable}{L{2.4cm} L{5.6cm} L{3.4cm} L{2.4cm}}
\caption{Tácticas arquitectónicas, componente donde residen y escenario que las verifica.}\label{tab:tacticas}\\
\toprule
\textbf{Atributo} & \textbf{Táctica} & \textbf{Dónde reside} & \textbf{Verifica} \\
\midrule
\endfirsthead
\toprule
\textbf{Atributo} & \textbf{Táctica} & \textbf{Dónde reside} & \textbf{Verifica} \\
\midrule
\endhead
\midrule
\contlinea{4} \\
\endfoot
\bottomrule
\endlastfoot
Seguridad & No capturar el dato sensible. La causa clínica no existe en el modelo, el contrato, los registros ni las métricas. & Modelo de dominio del Servicio de Donación & EC-02 \\
Seguridad & Doble verificación de autorización: compuerta por ruta y rol en el gateway, filtrado por jurisdicción a nivel de registro en el propietario. & Gateway y cada servicio & EC-01, EC-03 \\
Seguridad & Respuesta de rechazo que no confirma ni niega la existencia del recurso. & Gateway y cada servicio & EC-01 \\
Seguridad & Firma asimétrica con emisor único, token de acceso de quince minutos, rotación de la clave publicada con antelación y rechazo ante la ausencia de claves públicas. & Servicio de Identidad, gateway y cada servicio & EC-03 \\
Seguridad & Zonas de confianza con redes propias y aristas enumeradas; el proxy de borde como única superficie publicada; rutas internas nunca enrutadas por el gateway. & Composición de contenedores, proxy de borde, gateway & EC-01, EC-03 \\
Seguridad & Registro de solo anexado, sin operación de modificación ni borrado expuesta, con catálogo explícito de acciones sensibles. & ST2 auditoría & EC-05 \\
Seguridad & Máquina de estados explícita con matriz cerrada de transiciones válidas. & Servicio de Donación & EC-04 \\
Seguridad & Límite de tasa en dos escalas ---por origen en rutas públicas y por sesión en rutas protegidas--- e identificadores de unidad no adivinables. & Gateway y modelo de dominio & EC-06 \\
Seguridad física & Comportamiento restrictivo por defecto: ante estado indeterminado, la unidad se trata como no disponible en toda ruta de disponibilidad. & Servicio de Donación & EC-12 \\
Fiabilidad & Plazo de espera acotado en toda llamada entre servicios y degradación explícita en la interfaz. & Cliente HTTP de cada servicio & EC-08 \\
Fiabilidad & Sondeo de salud que verifica la dependencia de base de datos, con política de reinicio automático. & Composición de contenedores & EC-07 \\
Fiabilidad & Idempotencia por identificador de operación generado por el cliente. & Servicio de Donación & EC-10 \\
Fiabilidad & Bandeja de salida transaccional expuesta por contrato, con arrendamiento, reintento de espera creciente y estado de revisión manual visible. & Bandeja del Servicio de Donación y trabajador de Notificaciones & EC-11 \\
Fiabilidad & Migraciones versionadas y reversibles, aplicadas por un trabajo efímero por base, como base del ensayo de restauración. & Las cuatro bases de datos & EC-09 \\
Adecuación funcional & Ejecutor programado dentro del servicio propietario del inventario, con el sistema como actor de auditoría. & Servicio de Donación & EC-13 \\
Eficiencia & Medición en el gateway, no dentro del servicio, sobre volumen sintético declarado y versionado. & Gateway, Prometheus & EC-15, EC-16 \\
Mantenibilidad & Identificador de correlación generado en el gateway y propagado en cada llamada entre servicios. & Gateway y cada servicio & EC-20 \\
Mantenibilidad & Regla de elegibilidad en un único punto del dominio; la interfaz muestra el resultado, no lo replica. & Servicio de Donación & EC-21 \\
Compatibilidad & Contrato independiente de lenguaje, versionado semántico y consumidor tolerante a campos desconocidos. & Contratos OpenAPI & EC-22 \\
Flexibilidad & Toda diferencia entre ambientes expresada como variable de entorno; ninguna en la imagen. & Composición de contenedores & EC-23 \\
Flexibilidad & Tipo de donación como valor de catálogo, no como estructura propia. & Modelo de dominio del Servicio de Donación & EC-24 \\
Capacidad de interacción & Navegación derivada de la matriz usuario × módulo: la entrada que no corresponde no se renderiza, en lugar de aparecer deshabilitada. & Aplicación web & EC-01, EC-17 \\
Capacidad de interacción & Confirmación de acción irreversible que nombra la unidad concreta, separada físicamente de la acción frecuente, precedida de advertencia explícita de irreversibilidad. & Aplicación web & EC-19, EC-41 \\
Capacidad de interacción & Capa única de componentes compartidos: todo diálogo, tabla y estado se consume de allí y no se duplica por módulo. & Aplicación web & EC-35, EC-32 \\
Capacidad de interacción & Mensaje de rechazo en lenguaje de dominio con acción siguiente, sin trazas ni información fuera del alcance del usuario. & Gateway y cada servicio & EC-31 \\
Seguridad & Actor \emph{sistema} en el registro de toda acción automática, distinguible del actor persona. & ST2 auditoría & EC-33 \\
Seguridad física & Condiciones operativas de despacho verificadas contra el tipo y la custodia declarados de la institución. & Servicio de Donación & EC-39, EC-42 \\
Seguridad física & Umbral de vencimiento próximo configurable por institución, evaluado por proceso programado con antelación suficiente. & Servicio de Donación & EC-40 \\
Mantenibilidad & Matriz de transiciones como dato de catálogo, que permite generar los 82 casos de prueba en lugar de escribirlos. & Servicio de Donación & EC-36 \\
Eficiencia & Límites de recurso declarados por contenedor y rechazo explícito con código antes del agotamiento. & Composición de contenedores, Gateway & EC-26, EC-37 \\
Flexibilidad & Ausencia de estado en memoria en el Servicio de Donación, que admite una segunda instancia tras el gateway sin cambio de código ni de imagen; sus procesos programados toman un bloqueo en la base antes de cada ciclo, de modo que una sola instancia lo ejecuta. & Servicio de Donación & EC-37 \\
Flexibilidad & Responsabilidades del borde expresadas como configuración portable, sin extensiones propietarias del producto elegido. & API Gateway & EC-38 \\
Compatibilidad & Puertos, volúmenes y límites declarados por contenedor, un único puerto publicado al anfitrión y ningún volumen compartido entre bases. & Composición de contenedores & EC-27 \\
\end{longtable}

\noindent\textbf{Sobre la última táctica de navegación.} Ocultar en lugar de deshabilitar tiene una razón que no es estética: un elemento deshabilitado o con candado confirma que existe algo detrás, y eso ya es información sobre otra jurisdicción. La táctica sirve a la usabilidad y a la confidencialidad a la vez, pero no sustituye al control del servidor: es presentación, y el control vive en el gateway y en el servicio.

% =====================================================================
%  7. ARQUITECTURA DE NEGOCIO
% =====================================================================
\section{Arquitectura de negocio}
\label{sec:negocio}

\subsection{El problema que la arquitectura debe resolver}

Colombia recibió 830.000 unidades de sangre en 2023 frente a una necesidad anual estimada de 1,6 millones. En 2024 hubo 997.115 donaciones, un 0,3\,\% menos que el año anterior; el 93,8\,\% fueron voluntarias, pero solo el 27,1\,\% de los donantes son habituales. El déficit se agrava de forma estacional y predecible: las donaciones caen un 8\,\% en diciembre y un 12\,\% en enero.

De ese cuadro se derivan dos pérdidas concretas que la arquitectura debe atacar, y conviene nombrarlas porque explican por qué el sistema tiene la forma que tiene:

\begin{enumerate}
\item \textbf{Unidades que vencen sin usarse en un banco mientras otro tiene escasez.} Es un problema de visibilidad del inventario entre instituciones, no de captación. Su respuesta arquitectónica es el inventario con propiedad de datos declarada y la red de transferencias.
\item \textbf{Donantes que no vuelven porque nadie los convoca.} Es un problema de continuidad de la relación con el donante. Su respuesta arquitectónica es el perfil con elegibilidad calculada y el servicio de notificaciones.
\end{enumerate}

\subsection{Principio de escalamiento en dos niveles}

Es la regla de negocio que más condiciona el diseño. Ante una necesidad de un componente sanguíneo, el sistema privilegia primero la redistribución entre bancos de la red con excedente, y solo escala a movilizar donantes cuando la red interna no alcanza.

\begin{figure}[h!]
\centering
\begin{tikzpicture}[
  node distance=0.5cm,
  paso/.style={draw=vinotinto!70, fill=vinotinto!7, rounded corners=3pt, align=center, font=\scriptsize, inner sep=4pt, minimum height=1.0cm, text width=2.9cm},
  dec/.style={draw=ringProbar!90, fill=ringProbar!8, rounded corners=3pt, align=center, font=\scriptsize, inner sep=4pt, minimum height=1.0cm, text width=2.9cm},
  fin/.style={draw=ringEvaluar!85, fill=ringEvaluar!8, rounded corners=3pt, align=center, font=\scriptsize, inner sep=4pt, minimum height=0.95cm, text width=2.9cm},
  flecha/.style={-{Stealth}, draw=vinotinto!75, thick}
]
\node[paso] (a) at (0,0) {Alerta de escasez\\ RF-19\\ \emph{dispara el proceso}};
\node[dec] (b) at (3.6,0) {¿Hay excedente\\ compatible en el departamento\\ de la institución deficitaria?};
\node[paso] (c) at (7.5,1.35) {\textbf{Nivel 1}\\ Transferencia sugerida,\\ aprobada por el cedente\\ RF-10};
\node[paso] (d) at (7.5,-1.35) {\textbf{Nivel 2}\\ Convocar donantes\\ compatibles y elegibles\\ RF-14};
\node[fin] (e) at (11.6,1.35) {Unidad redistribuida};
\node[fin] (f) at (11.6,-1.35) {Campaña o convocatoria};

\draw[flecha] (a) -- (b);
\draw[flecha] (b.north) .. controls +(0,1.0) and +(-1.4,0) .. node[above, font=\tiny, text=vinotinto] {sí} (c.west);
\draw[flecha] (b.south) .. controls +(0,-1.0) and +(-1.4,0) .. node[below, font=\tiny, text=vinotinto] {no} (d.west);
\draw[flecha] (c) -- (e);
\draw[flecha] (d) -- (f);
\end{tikzpicture}
\caption[Escalamiento en dos niveles ante un déficit]{Escalamiento en dos niveles. La evaluación la ejecuta un proceso del Servicio de Donación con autoridad de sistema, no una consulta de usuario, y su alcance es el departamento de la \emph{institución deficitaria}, no la jurisdicción de quien consulta. Es lo que hace ejecutable el nivel~1: un administrador de banco, cuyo alcance de datos es su propia institución, nunca encontraría excedente ajeno si la búsqueda se filtrara por su jurisdicción. Ninguna institución conoce las existencias de otra: el sistema propone y cada administrador decide sobre lo suyo (ADR-013).}
\label{fig:escalamiento}
\end{figure}

La evaluación del nivel~1 necesita inventario ---Servicio de Donación--- y territorio de instituciones ---Servicio Institucional---. ADR-013 sitúa esa evaluación en un proceso programado del Servicio de Donación, que se dispara al crearse la alerta de escasez y se repite cada quince minutos sobre las alertas abiertas, y no en una petición interactiva. Una institución tiene excedente cuando sus unidades disponibles del componente y de un grupo compatible superan su propio umbral mínimo y conservan al menos cuarenta y ocho horas de vida útil. Si ninguna del departamento lo tiene, el proceso crea una convocatoria a los donantes compatibles y elegibles, que genera los avisos en la bandeja de salida. La consecuencia sobre la estructura es doble. El cruce de la frontera entre contextos sale del camino crítico de las peticiones de usuario, de modo que el plazo de espera acotado degrada sin afectar a ningún tiempo de respuesta percibido, y el compromiso PS-02 conserva una sola de sus dos tensiones. Y el \emph{driver} DF-04 deja de describirse como una acción de usuario: es el único proceso del sistema que necesita datos de dos contextos para decidir.

\vspace{0.2cm}
\noindent\textbf{Confidencialidad entre instituciones.} El resultado del nivel~1 se materializa como una transferencia en estado sugerida. El administrador del banco receptor ve que existe una transferencia sugerida para cubrir su déficit; el del banco cedente ve la solicitud sobre sus propias unidades y la aprueba o la rechaza. En ningún momento una institución consulta el inventario de otra, de modo que el escalamiento opera sin abrir excepción alguna a EC-01.

\subsection{Capacidades de negocio y su realización}

\begin{longtable}{L{3.2cm} L{5.4cm} L{2.4cm} L{2.8cm}}
\caption{Capacidades de negocio, módulo que las realiza y servicio propietario.}\label{tab:capacidades}\\
\toprule
\textbf{Capacidad} & \textbf{Qué resuelve en el negocio} & \textbf{Módulo} & \textbf{Servicio} \\
\midrule
\endfirsthead
\toprule
\textbf{Capacidad} & \textbf{Qué resuelve en el negocio} & \textbf{Módulo} & \textbf{Servicio} \\
\midrule
\endhead
\midrule
\contlinea{4} \\
\endfoot
\bottomrule
\endlastfoot
Captación del donante & Baja la barrera de entrada para convertir a la persona que quiere donar en un registro utilizable, sin exigir identificación. & M1 & Donación \\
Fidelización del donante & Calcula elegibilidad, mantiene historial y reconoce la donación recurrente sin remunerarla. Ataca el 27,1\,\% de donantes habituales. & M1 & Donación \\
Convocatoria y campañas & Planifica, publica y da seguimiento a campañas con cupo y sede, visibles solo en la jurisdicción correspondiente. & M2 & Campañas \\
Trazabilidad de la unidad & Conserva origen, estados y responsables de cada unidad, de la captación a la disposición final. Es la capacidad que sostiene la seguridad del paciente. & M3 & Donación \\
Visibilidad del inventario & Hace conocible la existencia real por componente y tipo de sangre, con umbrales configurables por institución. & M4 & Donación \\
Redistribución en red & Convierte un excedente ocioso de un banco en la cobertura del déficit de otro, antes de pedir sangre nueva. & M5 & Donación \\
Gobierno territorial & Mantiene la jerarquía nacional, departamental y municipal y el registro de instituciones. Es el eje del control de acceso. & M6 & Institucional \\
Supervisión e indicadores & Da a la autoridad sanitaria una vista agregada al nivel de su jurisdicción, hoy inexistente. & M7 & Institucional \\
Rendición de cuentas & Permite establecer quién hizo qué sobre una unidad, y qué intentos de acceso indebido hubo. & ST2 & Transversal \\
\end{longtable}

\subsection{Proceso central: ciclo de vida de la unidad}

Es el proceso que da nombre al sistema y el que concentra las restricciones críticas. La figura muestra los estados y las transiciones admitidas; el conjunto de pares no dibujados es tan importante como los dibujados, porque EC-04 exige que la matriz sea cerrada.

\begin{figure}[h!]
\centering
\begin{tikzpicture}[
  estado/.style={draw=vinotinto!75, fill=vinotinto!7, rounded corners=3pt, align=center, font=\scriptsize, inner sep=3pt, minimum height=0.85cm, text width=2.4cm},
  terminal/.style={draw=ringEvitar!90, fill=ringEvitar!10, rounded corners=3pt, align=center, font=\scriptsize, inner sep=3pt, minimum height=0.85cm, text width=2.4cm},
  flecha/.style={-{Stealth}, draw=vinotinto!75, thick},
  flechasis/.style={-{Stealth}, draw=ringProbar!90, thick}
]
\node[estado] (cap)  at (0,2.7)   {Captada\\ \tiny donación registrada};
\node[estado] (frac) at (3.4,2.7) {Fraccionada\\ \tiny unidad generada};
\node[estado] (tam)  at (6.8,2.7) {En tamizaje};
\node[estado] (disp) at (10.2,2.7){Disponible};

\node[terminal] (noapta) at (3.4,0.6) {No apta};
\node[terminal] (venc)   at (6.8,0.6) {Vencida};
\node[estado]   (res)    at (10.2,0.6){Reservada o\\ en transferencia};

\node[terminal] (des)  at (5.1,-1.5)  {Desechada\\ \tiny disposición final};
\node[terminal] (desp) at (10.2,-1.5) {Despachada};

\draw[flecha] (cap) -- (frac);
\draw[flecha] (frac) -- (tam);
\draw[flecha] (tam) -- node[above, font=\tiny, text=vinotinto] {apta} (disp);
\draw[flecha] (tam.south west) -- node[above, font=\tiny, text=vinotinto, sloped] {no apta} (noapta.north east);
\draw[flechasis] (disp.south west) -- node[above, font=\tiny, text=ringProbar, sloped] {\textsc{sistema}} (venc.north east);
\draw[flecha] ([xshift=-14pt]disp.south) -- node[left, font=\tiny, text=vinotinto] {reserva} ([xshift=-14pt]res.north);
\draw[flecha] ([xshift=14pt]res.north) -- node[right, font=\tiny, text=vinotinto] {liberada} ([xshift=14pt]disp.south);
\draw[flecha] (res) -- (desp);
\draw[flecha] (noapta.south east) -- (des.north west);
\draw[flecha] (venc.south) -- (des.north east);
\end{tikzpicture}
\caption[Ciclo de vida de la unidad]{Ciclo de vida de la unidad. La única transición cuyo actor no es una persona es la de vencimiento, que ejecuta el proceso programado y que la bitácora atribuye al sistema. \emph{No apta}, \emph{Vencida} y \emph{Desechada} no tienen retorno: ninguna transición devuelve una unidad a \emph{Disponible} desde ellas, y esa ausencia es la que EC-04 comprueba de forma exhaustiva.}
\label{fig:ciclovida}
\end{figure}

\subsection{Reglas de negocio invariantes}

Son las que no dependen de una versión del producto ni de una decisión de alcance. Una historia que las contradiga no se estima: se rechaza. Se identifican con el prefijo RNI, regla de negocio invariante, para no colisionar con el prefijo RN que el Documento de Herramientas, Políticas y Lineamientos reserva a las restricciones normativas de su Sección~2.2.

\begin{longtable}{L{1.3cm} L{8.4cm} L{4.6cm}}
\caption{Reglas de negocio invariantes del sistema.}\label{tab:invariantes}\\
\toprule
\textbf{ID} & \textbf{Regla} & \textbf{Fundamento y verificación} \\
\midrule
\endfirsthead
\toprule
\textbf{ID} & \textbf{Regla} & \textbf{Fundamento y verificación} \\
\midrule
\endhead
\midrule
\contlinea{3} \\
\endfoot
\bottomrule
\endlastfoot
RNI-01 & Ninguna respuesta del sistema expone la causa clínica por la que una unidad fue marcada como no apta. Se cumple no capturando el dato. & Ley 1581 de 2012; RI-02; EC-02. \\
RNI-02 & Ningún reconocimiento al donante tiene valor económico, directo ni indirecto. La prioridad de atención cuenta como valor indirecto. & Decreto 1571 de 1993; EC-14. \\
RNI-03 & Un usuario territorial nunca ve datos fuera de su jurisdicción, y la interfaz tampoco insinúa que existan. El administrador nacional tiene el país como jurisdicción propia y queda fuera del alcance de esta regla. & RNF-02; EC-01. \\
RNI-04 & Una unidad no apta o vencida no puede despacharse, transferirse ni reservarse. Ante estado indeterminado se trata como no disponible. & Resolución 901 de 1996; EC-12. \\
RNI-05 & El historial de trazabilidad de una unidad solo admite anexar eventos. Nada de lo escrito se edita ni se borra. & RF-04; EC-04, EC-05. \\
RNI-06 & El auditor no dispone de ninguna operación de escritura en el sistema. & RI-04; EC-05. \\
RNI-07 & El sistema opera exclusivamente sobre datos sintéticos, y su origen ficticio es evidente en la propia interfaz. & Restricción de proyecto; EC-23. \\
\end{longtable}

% =====================================================================
%  8. ARQUITECTURA DE DATOS
% =====================================================================
\section{Arquitectura de datos}
\label{sec:datos}

Esta sección define la propiedad, la clasificación, la consistencia y la retención de los datos. El modelo entidad--relación con atributos, las claves y los contratos de operación corresponden al DD V2 de esta misma entrega y no se anticipan aquí.

\subsection{Principios}

\begin{enumerate}
\item \textbf{Un dato tiene un único dueño.} El servicio propietario es la única vía de acceso; ningún otro componente lo lee de la base.
\item \textbf{Ninguna base es compartida.} Bases separadas, credenciales separadas. Es lo que impide K-04.
\item \textbf{Lo que no debe exponerse, no se modela.} La confidencialidad se implementa en el esquema, no en la capa de presentación.
\item \textbf{El dominio no nombra el material donado.} El tipo de donación es un valor de catálogo, no una estructura.
\item \textbf{Todo cambio de esquema es una migración versionada y reversible.} Es la condición que hace verificable el ensayo de restauración de EC-09.
\end{enumerate}

\subsection{Reparto de entidades por servicio}

\begin{longtable}{L{3.4cm} L{5.0cm} L{5.8cm}}
\caption{Entidades principales, servicio propietario y regla de acceso.}\label{tab:entidades}\\
\toprule
\textbf{Entidad} & \textbf{Servicio propietario} & \textbf{Quién más la necesita y cómo la obtiene} \\
\midrule
\endfirsthead
\toprule
\textbf{Entidad} & \textbf{Servicio propietario} & \textbf{Quién más la necesita y cómo la obtiene} \\
\midrule
\endhead
\midrule
\contlinea{3} \\
\endfoot
\bottomrule
\endlastfoot
Donante & Donación & Nadie más. Ningún otro servicio conoce donantes. \\
Donación & Donación & M7 para indicadores agregados, por contrato y en solo lectura. \\
Unidad & Donación & M7 para indicadores agregados, por contrato y en solo lectura. \\
Evento de trazabilidad & Donación & Nadie más. Es interno al contexto y de solo anexado. \\
Existencia de inventario & Donación & M7 para indicadores; el proceso de escalamiento, dentro del mismo servicio. \\
Transferencia & Donación & Requiere identidad y departamento de las instituciones cedente y receptora, obtenidos de M6 por contrato. \\
Reconocimiento & Donación & Nadie más. Catálogo cerrado y simbólico. \\
Evento de notificación & Donación & El trabajador de Notificaciones, por el contrato interno de la bandeja de salida; nunca por acceso a la base. \\
Institución (banco) & Institucional & Referenciada por unidad, inventario, transferencia y campaña mediante identificador; nunca por copia del registro completo. \\
Territorio (jerarquía) & Institucional & Es el eje de la jurisdicción. Viaja en el token de sesión como ruta territorial y no se consulta en cada petición. \\
Campaña & Campañas & M1 la consulta para mostrar campañas al donante y para generar los avisos de campañas publicadas. \\
Reserva de cupo & Campañas & Nadie más. El cupo y sus reservas viven en la misma base, de modo que reservar no exige transacción distribuida. \\
Usuario, rol y jurisdicción asignada & Identidad & Es el emisor de la identidad. Ningún otro componente almacena credenciales. El trabajador de Notificaciones consulta los contactos de los administradores de una institución por contrato interno. \\
Sesión y credencial de servicio & Identidad & Nadie más. Guardan resúmenes, nunca el secreto ni el token de renovación en claro. \\
Registro de auditoría & Cada servicio con base, el suyo & Escrito por cada servicio sobre sus propias entidades, y por Identidad sobre las denegaciones que le entrega el gateway; consultado por U7 en solo lectura mediante la composición que sirve el Servicio Institucional (ADR-016). \\
\end{longtable}

\subsection{Datos que cruzan la frontera y política de consistencia}
\label{sec:consistencia}

Distribuir tiene un costo que este documento no disimula: los datos institucionales y de campañas que el Servicio de Donación necesita no están en su base, y ningún servicio conserva la jerarquía territorial completa salvo el Institucional. Se declara para cada caso qué garantía se ofrece.

\begin{longtable}{L{3.6cm} L{3.4cm} L{7.2cm}}
\caption{Datos que cruzan la frontera entre servicios y garantía declarada.}\label{tab:consistencia}\\
\toprule
\textbf{Dato} & \textbf{Garantía} & \textbf{Justificación y modo de fallo aceptado} \\
\midrule
\endfirsthead
\toprule
\textbf{Dato} & \textbf{Garantía} & \textbf{Justificación y modo de fallo aceptado} \\
\midrule
\endhead
\midrule
\contlinea{3} \\
\endfoot
\bottomrule
\endlastfoot
Identidad, rol y jurisdicción del usuario & Fuerte durante la vigencia del token de acceso, máximo quince minutos & El token transporta el dato para evitar una llamada entre servicios en cada verificación. El costo aceptado es que un cambio de jurisdicción tarda hasta quince minutos en surtir efecto (RF-22, ADR-008); es admisible porque un coordinador no cambia de departamento en mitad de una jornada. \\
Institución custodia de una unidad & Eventual & El Servicio de Donación guarda el identificador, no el registro. Si el nombre de la institución no puede resolverse, la interfaz lo declara no disponible; nunca muestra un valor por defecto. \\
Jerarquía territorial para filtrar inventario & Fuerte en la lectura & El filtro se aplica con la jurisdicción del token, no con una copia local de la jerarquía. \\
Territorio y ruta de la sede de una campaña & Fijada al crear la campaña & El Servicio de Campañas los obtiene del Institucional con el token del administrador que la crea y los guarda con la campaña, de modo que la consulta pública de campañas se filtra por jurisdicción sin llamar a otro servicio. \\
Campaña asociada a una donación & Eventual & Si el Servicio de Campañas no responde, la donación se registra igual: la operación principal nunca se bloquea por un dato accesorio (EC-08). \\
Eventos notificables & Fuerte en la escritura, eventual en la entrega & El evento se escribe en la bandeja en la misma transacción que el hecho que lo origina; la entrega ocurre después, con reintento, y nunca duplica ni inventa un aviso. \\
Series de M3 y M4 para indicadores & Eventual & M7 es de solo lectura sobre datos ajenos (RI-03). Un indicador con retraso es aceptable; un indicador que escriba en el contexto de origen no lo es. \\
\end{longtable}

\noindent\textbf{Orquestación del escalamiento (ADR-013).} La evaluación cruza dos contextos, de modo que uno debe orquestarla. La orquesta el \textbf{Servicio de Donación}, por tres razones: la decisión depende de existencias, que son suyas; el estado de la transferencia es un estado de la unidad, que también es suyo; y la información institucional que necesita ---identidad de las instituciones y su territorio--- es de consulta y no de escritura, de modo que la llamada cruzada es de solo lectura y puede degradarse sin comprometer la operación. La alternativa, orquestar desde el Institucional, obligaría a escribir en el contexto de donación desde fuera, lo que contradice la propiedad exclusiva de datos.

\vspace{0.2cm}
\noindent El proceso se ejecuta con autoridad de sistema sobre el inventario completo, no bajo la jurisdicción de ningún usuario, y su alcance de búsqueda es el departamento de la institución deficitaria. Es legítimo porque el Servicio de Donación es el propietario exclusivo de todas las existencias y porque el proceso no tiene sesión asociada: invoca al Servicio Institucional con su propio token de servicio (ADR-008). Si el Institucional no responde en el plazo de 3 segundos, el ciclo termina sin crear nada y la alerta se evalúa de nuevo en el ciclo siguiente. La creación de la sugerencia se registra en la bitácora con tipo de actor sistema, como toda acción automática, y por tanto queda cubierta por EC-33.

\subsection{Clasificación de datos y datos prohibidos}

\begin{longtable}{L{3.0cm} L{4.6cm} L{6.6cm}}
\caption{Clasificación de los datos del sistema y tratamiento exigido.}\label{tab:clasificacion}\\
\toprule
\textbf{Clasificación} & \textbf{Datos} & \textbf{Tratamiento exigido} \\
\midrule
\endfirsthead
\toprule
\textbf{Clasificación} & \textbf{Datos} & \textbf{Tratamiento exigido} \\
\midrule
\endhead
\midrule
\contlinea{3} \\
\endfoot
\bottomrule
\endlastfoot
Sensible & Tipo de sangre y estado de aptitud asociado a una persona identificada. & Autorización explícita, previa e informada, en casilla propia y separada de la de avisos y campañas. Nunca en registros de aplicación ni en etiquetas de métricas. \\
Personal & Identificación, contacto y historial del donante registrado (U2). & Accesible únicamente al propio titular y a los perfiles que la matriz del SRS autoriza. Revocable por el titular. \\
Operacional restringido & Inventario, unidades y transferencias de una institución. & Filtrado por jurisdicción a nivel de registro. Nunca visible fuera del ámbito del usuario. \\
Operacional agregado & Indicadores por jurisdicción. & Agregados al nivel del usuario que consulta. El nivel nacional es jurisdicción propia de U6, no una excepción a la regla. \\
Auditoría & Actor, rol, jurisdicción, operación, recurso, resultado y marca temporal. & Solo anexado. Sin causa clínica ni datos personales innecesarios para identificar la acción. \\
\textbf{Prohibido} & \textbf{Causa clínica, diagnóstico o resultado individual de tamizaje.} & \textbf{No se captura, no se almacena, no se transmite y no se registra. No existe campo que lo admita en ningún componente del sistema.} \\
\end{longtable}

\subsection{Retención, respaldo y evolución del esquema}

\begin{itemize}
\item \textbf{Trazabilidad y auditoría no se depuran} dentro del horizonte del proyecto. Depurar el historial de una unidad destruye la capacidad que el sistema existe para ofrecer.
\item \textbf{Respaldo diario} de las cuatro bases, con ensayo de restauración en cada sprint desde el Sprint 2. La pérdida máxima admitida es de 24 horas y el tiempo de recuperación objetivo es de 4 horas, cifras válidas para un ambiente con datos sintéticos y que deberían revisarse en un despliegue real.
\item \textbf{Migraciones versionadas y reversibles} en los cuatro servicios con base, aplicadas por un trabajo de migración efímero por base y gestionadas de forma independiente. La versión de esquema es dato observable, no una suposición: sin ella, la medida de respuesta de EC-09 no se puede afirmar.
\item \textbf{Solo datos sintéticos}, cargados como parte del arranque del sistema y visiblemente ficticios en la propia demostración.
\end{itemize}

% =====================================================================
%  9. ARQUITECTURA DE INFRAESTRUCTURA
% =====================================================================
\section{Arquitectura de infraestructura}
\label{sec:infra}

\subsection{Principio de contenerización}

RD-03 exige que todo esté contenerizado. La restricción pide contenerización, no orquestación en clúster: Kubernetes está en anillo \evitar{} del catálogo, y esa exclusión es deliberada. Un clúster añadiría una superficie operativa que el equipo no puede sostener en cinco sprints y que ningún requerimiento pide.

La consecuencia práctica de RD-03 es la que mide EC-23: el sistema completo se levanta con un comando, en menos de diez minutos, en cualquiera de los tres sistemas operativos que usa el equipo, por alguien que no participó en la configuración.

\subsection{Topología de despliegue y zonas de confianza}

La topología agrupa los componentes en cinco zonas de confianza, cada una con su propia red de contenedores (ADR-014). Una zona reúne componentes con el mismo nivel de exposición, y entre dos zonas contiguas hay una frontera con controles propios. Solo el proxy de borde publica puertos al anfitrión: 443/tcp, y 80/tcp únicamente para redirigir a HTTPS. Todas las redes, salvo la de borde, se declaran internas y sin salida al exterior.

\begin{figure}[h!]
\centering
\begin{tikzpicture}[
  zona/.style={draw=vinotinto!45, fill=#1, rounded corners=4pt, minimum width=11.2cm, inner sep=0pt},
  c/.style={draw=vinotinto!70, fill=white, rounded corners=2pt, align=center, font=\scriptsize, inner sep=2pt, minimum height=0.65cm, text width=1.9cm},
  cb/.style={draw=ringEvaluar!85, fill=white, rounded corners=2pt, align=center, font=\tiny, inner sep=2pt, minimum height=0.65cm, text width=1.9cm},
  co/.style={draw=ringProbar!85, fill=white, rounded corners=2pt, align=center, font=\scriptsize, inner sep=2pt, minimum height=0.6cm, text width=2.2cm},
  zl/.style={anchor=north west, font=\scriptsize\bfseries, text=vinotinto},
  fr/.style={anchor=west, font=\tiny, text=grisTexto, fill=white, inner sep=1pt}
]
\node[zona=vinotinto!14, minimum height=1.25cm] (z1) at (-2.0,7.0) {};
\node[zona=vinotinto!9, minimum height=1.25cm] (z2) at (-2.0,5.3) {};
\node[zona=vinotinto!5, minimum height=1.45cm] (z3) at (-2.0,3.5) {};
\node[zona=ringEvaluar!7, minimum height=1.95cm] (z4) at (-2.0,1.4) {};
\node[zl] at (z1.north west) {\;Zona 1 · Borde};
\node[zl] at (z2.north west) {\;Zona 2 · Presentación};
\node[zl] at (z3.north west) {\;Zona 3 · Aplicación};
\node[zl] at (z4.north west) {\;Zona 4 · Datos};

\node[c, text width=3.4cm] at (-2.0,6.85) {Proxy de borde · Caddy\\ 443 · 80 → 443};
\node[c] at (-4.2,5.15) {Aplicación web};
\node[c] at (0.2,5.15) {API Gateway};
\node[c] at (-6.4,3.3) {Identidad};
\node[c] at (-4.2,3.3) {Institucional};
\node[c] at (-2.0,3.3) {Campañas};
\node[c] at (0.2,3.3) {Donación};
\node[c] at (2.4,3.3) {Notificaciones};
\node[cb] at (-6.4,1.6) {\texttt{db\_identidad}};
\node[cb] at (-4.2,1.6) {\texttt{db\_institucional}};
\node[cb] at (-2.0,1.6) {\texttt{db\_campana}};
\node[cb] at (0.2,1.6) {\texttt{db\_donacion}};
\node[font=\tiny, text=grisTexto, align=center, text width=2.0cm] at (2.4,1.55) {sin base propia: consume la bandeja de Donación};
\node[font=\tiny, text=grisTexto] at (-2.0,0.72) {trabajos de migración, respaldo y restauración: cada uno solo en la red de su base};

\node[draw=ringProbar!70, fill=ringProbar!8, rounded corners=4pt, minimum width=2.9cm, minimum height=2.6cm] (z5) at (5.9,4.4) {};
\node[anchor=north, font=\scriptsize\bfseries, text=vinotinto, align=center] at (z5.north) {Zona 5 · solo QA\\ Observabilidad};
\node[co] at (5.9,4.2) {Prometheus};
\node[co] at (5.9,3.45) {Grafana};

\node[font=\scriptsize, text=grisTexto, anchor=west] at (-1.8,8.25) {Navegador · HTTPS 443};
\draw[-{Stealth}, thick, draw=vinotinto!70] (-2.0,8.45) -- (-2.0,7.64);
\node[fr] at (-7.55,7.9) {F1 · TLS, cabeceras, tamaño de cuerpo};
\node[fr] at (-7.55,6.15) {F2 · token, compuerta por ruta y rol, límite de tasa};
\node[fr] at (-7.55,4.45) {F3 · revalidación del token y filtrado por jurisdicción};
\node[fr] at (-7.55,2.575) {F4 · red propia por base, rol de servicio sin DDL};
\node[fr] at (4.45,5.95) {F5 · solo puerto de gestión};
\end{tikzpicture}
\caption[Zonas de confianza y fronteras]{Zonas de confianza y fronteras. Cada zona tiene su red; solo la Zona 1 es alcanzable desde fuera. F1 la vigila el proxy de borde (ADR-015); F2, el gateway (ADR-007); F3, cada servicio (ADR-008, ADR-011); F4, la red propia y los roles de cada base ---un rol propietario para migrar y respaldar, y un rol de servicio sin definición de esquema ni modificación de tablas de solo anexado--- (ADR-005, ADR-014); F5, Prometheus y Grafana. La Zona 5 existe solo en QA: Prometheus raspa únicamente el puerto de gestión de cada servicio por la red de aplicación y Grafana se alcanza solo a través del borde. Las aristas permitidas entre componentes son las de la Tabla \ref{tab:aristas}.}
\label{fig:despliegue}
\end{figure}

\begin{longtable}{L{2.8cm} L{4.6cm} L{7.0cm}}
\caption{Zonas de confianza, componentes y red de contenedores.}\label{tab:zonas}\\
\toprule
\textbf{Zona} & \textbf{Componentes} & \textbf{Red y exposición} \\
\midrule
\endfirsthead
\toprule
\textbf{Zona} & \textbf{Componentes} & \textbf{Red y exposición} \\
\midrule
\endhead
\midrule
\contlinea{3} \\
\endfoot
\bottomrule
\endlastfoot
1 · Borde & Proxy de borde & \texttt{red\_borde}. Único componente con puertos publicados al anfitrión. \\
2 · Presentación & Aplicación web, API Gateway & \texttt{red\_borde}; el gateway, además, \texttt{red\_aplicacion}. Alcanzables solo desde el borde. \\
3 · Aplicación & Identidad, Institucional, Campañas, Donación, Notificaciones & \texttt{red\_aplicacion}; cada servicio con base, además, su red de datos. Alcanzables solo desde el gateway y entre sí por las aristas permitidas. \\
4 · Datos & Cuatro PostgreSQL y sus trabajos de migración & Una red por base: \texttt{red\_datos\_identidad}, \texttt{red\_datos\_institucional}, \texttt{red\_datos\_campana} y \texttt{red\_datos\_donacion}. Cada base, alcanzable solo desde su servicio y su trabajo de migración. \\
5 · Observabilidad (QA) & Prometheus, Grafana & \texttt{red\_observabilidad}; Prometheus, además, \texttt{red\_aplicacion}; Grafana, además, \texttt{red\_borde}. No recibe tráfico de usuario salvo Grafana a través del borde. \\
\end{longtable}

\begin{longtable}{L{0.6cm} L{2.9cm} L{3.3cm} L{5.4cm} L{1.7cm}}
\caption{Aristas permitidas. Toda comunicación no listada está prohibida.}\label{tab:aristas}\\
\toprule
\textbf{\#} & \textbf{Origen} & \textbf{Destino} & \textbf{Propósito} & \textbf{Protocolo} \\
\midrule
\endfirsthead
\toprule
\textbf{\#} & \textbf{Origen} & \textbf{Destino} & \textbf{Propósito} & \textbf{Protocolo} \\
\midrule
\endhead
\midrule
\contlinea{5} \\
\endfoot
\bottomrule
\endlastfoot
1 & Navegador & Proxy de borde & Toda petición externa & HTTPS 443 \\
2 & Proxy de borde & Aplicación web & Recursos estáticos & HTTP 8080 \\
3 & Proxy de borde & API Gateway & Rutas \texttt{/api/*} & HTTP 8080 \\
4 & Proxy de borde & Grafana (QA) & Rutas \texttt{/grafana/*} & HTTP 3000 \\
5 & API Gateway & Identidad, Institucional, Campañas, Donación & Rutas publicadas \texttt{/v1/*} & HTTP 8080 \\
6 & API Gateway & Identidad & Claves públicas, canje de su token de servicio y entrega de denegaciones & HTTP 8080 \\
7 & Institucional, Campañas, Donación & Identidad & Claves públicas; Donación canjea además su token de servicio & HTTP 8080 \\
8 & Donación & Institucional & Instituciones y territorios & HTTP 8080 \\
9 & Donación & Campañas & Campaña de una donación y campañas publicadas & HTTP 8080 \\
10 & Campañas & Institucional & Sede, institución organizadora y jurisdicción & HTTP 8080 \\
11 & Institucional & Donación & Indicadores agregados y serie de auditoría & HTTP 8080 \\
12 & Institucional & Campañas, Identidad & Series de auditoría & HTTP 8080 \\
13 & Notificaciones & Donación & Bandeja de salida & HTTP 8080 \\
14 & Notificaciones & Identidad & Tokens de servicio y contactos institucionales & HTTP 8080 \\
15 & Cada servicio con base & Su PostgreSQL & Persistencia & TCP 5432 \\
16 & Cada trabajo de la base: migración, respaldo y restauración & Su PostgreSQL & Esquema, datos sintéticos, respaldo y restauración & TCP 5432 \\
17 & Prometheus (QA) & API Gateway y los cinco servicios & Métricas en el puerto de gestión & HTTP 8081 \\
18 & Grafana (QA) & Prometheus & Consulta de series & HTTP 9090 \\
\end{longtable}

\noindent\textbf{Invariantes de la topología.} No existe ninguna arista entre la Zona 1 y las Zonas 3 o 4. No existe ninguna arista entre dos bases. Cada base es alcanzable solo desde su servicio propietario y sus trabajos de migración, respaldo y restauración. Ningún puerto de gestión (8081) se enruta por el gateway. Las rutas con prefijo \texttt{/internal/} no son enrutadas nunca por el gateway. Toda llamada entre servicios tiene un plazo de espera máximo de 3 segundos y degradación declarada. El Documento de Infraestructura V1 reproduce esta tabla como matriz de alcanzabilidad y comprueba los invariantes sobre la composición real.

\vspace{0.2cm}
\noindent\textbf{Orden de arranque.} El arranque se encadena por condición de salud y no por orden de declaración: cada base arranca y queda sana; su trabajo de migración termina con éxito; arranca Identidad; arrancan en paralelo Institucional, Campañas y Donación; arranca el trabajador de Notificaciones, con Donación e Identidad sanos; arranca el gateway, con los cuatro servicios que enruta sanos; arrancan la aplicación web y, por último, el proxy de borde. En QA, Prometheus y Grafana arrancan en paralelo al resto.

\subsection{Ambientes}

\begin{longtable}{L{2.2cm} L{4.2cm} L{7.8cm}}
\caption{Ambientes del sistema y su propósito.}\label{tab:ambientes}\\
\toprule
\textbf{Ambiente} & \textbf{Propósito} & \textbf{Condiciones} \\
\midrule
\endfirsthead
\toprule
\textbf{Ambiente} & \textbf{Propósito} & \textbf{Condiciones} \\
\midrule
\endhead
\midrule
\contlinea{3} \\
\endfoot
\bottomrule
\endlastfoot
Desarrollo & Máquina de cada integrante. & Doce contenedores levantados con un comando, sin Prometheus ni Grafana. Certificado de la autoridad interna del proxy de borde, cuyo certificado raíz cada integrante puede instalar en su equipo. Secretos en archivos generados por el guion de inicialización, excluidos del control de versiones y montados en \texttt{/run/secrets}. Datos sintéticos cargados por los trabajos de migración. \\
QA & Verificación de los escenarios de calidad y demostración al cliente en los hitos de las Semanas 10, 12, 14 y 16. & Máquina virtual institucional con los catorce contenedores, levantados con el mismo archivo de orquestación de Desarrollo y desplegados desde la rama \texttt{main}. Certificado de la autoridad interna del proxy de borde, cuyo certificado raíz se entrega al cliente evaluador junto con la dirección de la máquina. Secretos de Compose montados en \texttt{/run/secrets}, legibles solo por el usuario de despliegue. Es el ambiente donde se miden EC-06 a EC-08, EC-11, EC-15, EC-16, EC-20, EC-26 y EC-37. Su alcanzabilidad desde fuera de la red de la universidad es condición de viabilidad y está registrada como riesgo en el Documento de Herramientas. \\
Producción & Declarada y no construida. & El sistema no opera con datos reales de personas. Se declara para que su ausencia sea una decisión y no un olvido. Su topología de referencia conserva las cinco zonas y añade lo que exige operar con datos reales: certificado de autoridad pública por ACME, balanceador y réplicas del gateway y de los servicios, réplica de lectura por base, gestor externo de secretos con rotación automática, canal de correo real para las notificaciones y respaldo cifrado fuera del anfitrión. Con datos reales rigen como obligaciones legales la retención, el derecho de supresión y la notificación de incidentes de la Ley 1581 de 2012, y la ventana de revocación de quince minutos del token se revisa antes de operar. \\
\end{longtable}

\noindent\textbf{Pila de medición.} Las pruebas de carga y de sobrecarga nunca se ejecutan contra la instancia de QA que se muestra al cliente. Se ejecutan sobre una pila de medición: la misma composición de QA, levantada con otro nombre de proyecto sobre la máquina virtual institucional, con la instancia de demostración detenida y fuera de las ventanas de demostración. La carga la genera Gatling desde el equipo de un integrante, fuera de la máquina virtual, para que el generador no compita por los recursos que se miden. Así la medición se hace sobre el mismo anfitrión y la misma composición que se muestran, sin arriesgar la demostración.

\vspace{0.2cm}
\noindent La configuración específica de cada ambiente ---variables, secretos, límites de recurso--- corresponde al Documento de Infraestructura V1 de esta misma entrega. Este documento fija únicamente la regla que lo gobierna: \textbf{toda diferencia entre ambientes se expresa como variable de entorno, y ninguna se hornea en la imagen}. Es la condición de EC-23 y de RNF-11.

\subsection{Seguridad de la infraestructura}

\begin{longtable}{L{3.6cm} L{11.0cm}}
\caption{Controles de seguridad en la capa de infraestructura.}\label{tab:seginfra}\\
\toprule
\textbf{Control} & \textbf{Aplicación} \\
\midrule
\endfirsthead
\toprule
\textbf{Control} & \textbf{Aplicación} \\
\midrule
\endhead
\midrule
\contlinea{2} \\
\endfoot
\bottomrule
\endlastfoot
Superficie expuesta mínima & Solo el proxy de borde publica puertos al anfitrión. Ningún servicio, base ni pieza de observabilidad lo hace. \\
Borde & TLS 1.2 y 1.3, redirección a HTTPS, cabeceras de seguridad ---seguridad de transporte estricta, política de contenido restringida al propio origen, prohibición de enmarcado, tipo de contenido sin inferencia, sin referente y sin permisos de cámara, micrófono ni ubicación---, cuerpo máximo de 1~MiB en \texttt{/api/*} y registro de acceso en JSON con credenciales y cookies redactadas (ADR-015). \\
Autenticación & Token RS256 emitido solo por el Servicio de Identidad, validado por el gateway y de nuevo por cada servicio; rechazo ante la ausencia de claves públicas; token de acceso de quince minutos; token de renovación opaco, rotatorio, en cookie \texttt{HttpOnly}, \texttt{Secure} y \texttt{SameSite=Strict}, guardado como resumen; rotación de la clave de firma en cada sprint sin invalidar sesiones (ADR-008). \\
Procesos sin usuario & El gateway, el Servicio de Donación y el trabajador de Notificaciones se autentican con credencial de servicio propia, canjeada por un token de servicio de quince minutos. Cuando una llamada entre servicios ocurre dentro de una petición de usuario, se propaga el token del usuario y no el de servicio. \\
Límite de tasa & En el gateway, en dos escalas: 60 peticiones por minuto y por origen en las rutas públicas, y 300 por minuto y por sesión en las protegidas, con rechazo 429 antes de enrutar (EC-06). \\
Rutas internas & Las operaciones con prefijo \texttt{/internal/} y la publicación de claves públicas no se enrutan nunca por el gateway: solo son alcanzables dentro de la red de aplicación. \\
Roles separados por base & Cada base tiene su red y tres roles: el superusuario, usado solo por el guion de inicialización; el rol propietario, que define el esquema y solo lo usan los trabajos de migración, respaldo y restauración; y el rol de servicio, con que opera el servicio, sin definición de esquema y con solo inserción y lectura sobre las tablas de solo anexado ---los eventos de trazabilidad de las unidades y las cuatro series de la bitácora---. Es el control que hace detectable el \emph{killer} K-04 y que impone la inmutabilidad de la bitácora en la base, no solo en el código. \\
Ausencia de secretos en el repositorio & Ninguna credencial, cadena de conexión ni clave de firma se versiona. La integración continua ejecuta detección de secretos en cada Pull Request y lo bloquea si encuentra uno. \\
Imágenes & Construcción multietapa sin compilador, gestor de paquetes ni intérprete de órdenes en la etapa de ejecución: imagen \emph{chiseled} de ASP.NET Core para los servicios .NET y \emph{distroless} con un entorno de ejecución de Java reducido con \texttt{jlink} para los servicios Java, con sondas de salud que no dependen de un intérprete de órdenes; escaneo de vulnerabilidades en la integración continua con rechazo de las de severidad alta o crítica con corrección publicada. \\
Endurecimiento de contenedores & Usuario no privilegiado, sistema de archivos raíz de solo lectura, capacidades del núcleo retiradas, escalada de privilegios prohibida y límites de CPU y memoria declarados. \\
Registros sin datos protegidos & Ningún registro de aplicación, traza, registro de acceso del borde o etiqueta de métrica contiene datos personales, causa clínica ni contenido de token (EC-20). \\
Denegaciones en la bitácora & Toda denegación queda en la bitácora: las del gateway, en la serie del Servicio de Identidad; las de cada servicio, en la suya. Una denegación del gateway que no puede entregarse se contabiliza en una métrica propia (ADR-016). \\
\end{longtable}

\subsection{Observabilidad}

Nueve de los cuarenta y dos escenarios ---EC-06, EC-07, EC-08, EC-11, EC-15, EC-16, EC-20, EC-26 y EC-37--- tienen medidas de respuesta que solo pueden afirmarse con medición. Sin observabilidad no se incumplen: quedan sin verificar, que a efectos de una entrega es equivalente.

\begin{itemize}
\item \textbf{Métricas} recolectadas por Prometheus y presentadas en Grafana, solo en QA. Las series relevantes son disponibilidad del \emph{endpoint} de salud, latencia por percentil medida \emph{en el gateway}, tasa de error, uso de CPU y memoria por contenedor, accesos denegados por jurisdicción y denegaciones del gateway no entregadas a la bitácora.
\item \textbf{Registros} con identificador de correlación generado en el gateway cuando el cliente no lo aporta, y propagado en la cabecera de toda llamada entre servicios. El registro de acceso del proxy de borde es el primer eslabón de la traza.
\item \textbf{Sondeos de salud} que verifican la dependencia de base de datos y no solo que el proceso esté vivo; el del trabajador de Notificaciones verifica que alcanza la bandeja del Servicio de Donación.
\end{itemize}

\noindent\textbf{Secuencia de incorporación.} La instrumentación de métricas se deja lista en el Sprint 2 sin recolector, de modo que el Sprint 4 solo tenga que levantar dos contenedores e importar tableros. Concentrar contenerización, integración continua, gateway y observabilidad en un mismo sprint es el error de planeación más previsible de este proyecto, y se evita repartiéndolo.

\subsection{Respaldo y recuperación}

Respaldo diario de las cuatro bases, con retención suficiente para cubrir el sprint en curso. El procedimiento de restauración está documentado y se ensaya en cada sprint desde el Sprint 2; los ensayos de los Sprints 2 a 4 validan el procedimiento, y el del Sprint 5 es el que cierra EC-09, porque es el primero que se ejecuta sobre el sistema completo con datos de todos los módulos.

Cada servicio se recupera de forma independiente: son bases separadas y restaurar una no obliga a tocar las demás. Esa independencia es una consecuencia buscada de ADR-005 y ADR-014, no una casualidad.

% =====================================================================
%  10. TRAZABILIDAD
% =====================================================================
\section{Trazabilidad}
\label{sec:trazabilidad}

La trazabilidad de este documento es bidireccional. Hacia atrás, toda decisión arquitectónica se rastrea hasta un \emph{driver}, un requerimiento o una restricción. Hacia adelante, todo atributo de calidad cuenta con al menos un escenario, una táctica y un componente donde reside.

\begin{longtable}{L{2.0cm} L{2.6cm} L{4.2cm} L{2.6cm} L{2.0cm}}
\caption{Matriz de trazabilidad entre requerimientos, escenarios, tácticas y componentes.}\label{tab:matriztraza}\\
\toprule
\textbf{Origen} & \textbf{Driver} & \textbf{Escenarios} & \textbf{Componente} & \textbf{ADR} \\
\midrule
\endfirsthead
\toprule
\textbf{Origen} & \textbf{Driver} & \textbf{Escenarios} & \textbf{Componente} & \textbf{ADR} \\
\midrule
\endhead
\midrule
\contlinea{5} \\
\endfoot
\bottomrule
\endlastfoot
RNF-01, RI-02 & DQ-01 & EC-02; contorno en EC-05, EC-12, EC-20, EC-31 & Modelo de dominio, Donación & ADR-005 \\
RNF-02 & DQ-02 & EC-01; apoyo en EC-03, EC-06 & Gateway y cada servicio & ADR-007, ADR-008, ADR-011 \\
RNF-09 & DQ-03 & EC-05, EC-04, EC-33 & ST2 auditoría, cuatro series & ADR-016 \\
RNF-03 & --- & EC-17, EC-28, EC-29, EC-32 & Aplicación web & ADR-006 \\
RNF-04 & --- & EC-15, EC-16, EC-26, EC-37 & Gateway, Donación & ADR-003, ADR-007 \\
RNF-05 & DQ-06 & EC-24 & Modelo de dominio, Donación & ADR-005, ADR-012 \\
RNF-06 & DQ-05 & EC-08, EC-11, EC-09 & Cliente HTTP, bandeja de salida, trabajador de Notificaciones & ADR-003, ADR-014 \\
RD-01 & DQ-08 & EC-22, EC-38 & Contratos OpenAPI, Gateway & ADR-003, ADR-007, ADR-014 \\
RNF-07 & --- & Sin escenario asociado & Punto de integración documentado & --- \\
RNF-08 & --- & EC-13, EC-12, EC-25, EC-40 & Proceso programado, Donación & ADR-003 \\
RNF-10 & --- & EC-18 & Aplicación web & ADR-006 \\
ADR-006 & --- & EC-35 & Capa de componentes compartidos & ADR-006 \\
RNF-11 & DQ-07 & EC-23, EC-27 & Composición de contenedores & ADR-003, ADR-014 \\
Táctica de infraestructura & --- & EC-07 & Composición de contenedores & ADR-003 \\
RD-01 & --- & EC-08, EC-22, EC-34 & Fronteras entre servicios & ADR-003, ADR-014 \\
RD-02 & --- & --- & Reparto Java / .NET & ADR-009, ADR-014 \\
RD-03 & DQ-07 & EC-23, EC-27 & Composición de contenedores, proxy de borde & ADR-003, ADR-015 \\
RD-08 & DQ-02, DQ-03 & EC-01 a EC-06, EC-33 & Proxy de borde, gateway, cada servicio & ADR-007, ADR-008, ADR-015, ADR-016 \\
RF-03 a RF-06 & DF-01 & EC-04, EC-10, EC-12, EC-25, EC-36 & Servicio de Donación & ADR-003, ADR-005, ADR-014 \\
RF-09 & DF-02 & EC-01, EC-03, EC-42 & Servicio Institucional & ADR-014 \\
RF-10, RF-14 & DF-04 & EC-12, EC-08, EC-39, EC-33 & Proceso de escalamiento, Donación & ADR-013 \\
RF-11 & --- & EC-14, EC-30 & Catálogo de reconocimientos, M1 & --- \\
RF-12 & --- & EC-11, EC-30 & Bandeja de salida, trabajador de Notificaciones & ADR-014 \\
RF-15 & --- & Sin escenario asociado & M2, Campañas & --- \\
RF-18 & DF-05 & EC-13, EC-33, EC-40 & Proceso programado & ADR-003 \\
RF-19 & DF-04 & EC-08, EC-33 & Alertas y proceso de escalamiento, Donación & ADR-013 \\
RF-20 & --- & EC-40 & Proceso programado, Donación & ADR-003 \\
RF-21 & DF-06 & EC-05, EC-32 & ST2 auditoría, consulta compuesta & ADR-016 \\
RF-22 & DF-02 & EC-01, EC-03 & Servicio de Identidad & ADR-008 \\
RF-23, RF-24 & --- & EC-03 & Servicio de Identidad & ADR-008 \\
Ley 1581/2012 & DQ-01 & EC-02, EC-05, EC-20 & Modelo, registros, métricas & ADR-005 \\
Decreto 1571/1993 & --- & EC-14, EC-30 & Catálogo de reconocimientos, M1 & --- \\
Resolución 901/1996 & DQ-04 & EC-12, EC-13, EC-19, EC-39, EC-41 & Donación, aplicación web & ADR-003 \\
\end{longtable}

\noindent\textbf{Estado de la trazabilidad.} La matriz traza todo requerimiento no funcional, toda restricción de diseño con consecuencia estructural y todo requerimiento funcional que es \emph{driver} o que tiene escenario de calidad asociado; los demás requerimientos funcionales se trazan hasta sus historias en el SRS. Diez de los once requerimientos no funcionales del SRS quedan cubiertos por al menos un escenario; las cuarenta subcaracterísticas de ISO/IEC 25010:2023 quedan cubiertas por al menos un escenario; los cuarenta y dos escenarios tienen índice de prioridad, táctica y componente asignados; y toda decisión estructural remite a un ADR aprobado.

Se registran tres excepciones deliberadas. RD-02 no tiene escenario de calidad asociado porque es una restricción de composición del \emph{stack}, no un atributo verificable por estímulo y respuesta; su cumplimiento se comprueba por inspección de los repositorios. RNF-07 no tiene escenario porque la interoperabilidad con SIHEVI-INS está fuera de la línea base de la versión 1: el sistema documenta el punto de integración pero no ejecuta el reporte. RF-15 no tiene escenario porque está clasificado \emph{Could}, y su verificación se reduce a la prueba funcional de su historia.

% =====================================================================
%  11. RIESGOS Y DEUDA
% =====================================================================
\section{Riesgos arquitectónicos y deuda declarada}
\label{sec:riesgos}

Los riesgos de gestión ---capacidad del equipo, concurrencia académica, fecha inamovible--- viven en el Documento de Proyecto. Aquí se registran únicamente los que tienen consecuencia estructural.

\begin{longtable}{L{1.2cm} L{4.6cm} L{1.5cm} L{6.9cm}}
\caption{Riesgos arquitectónicos y su mitigación.}\label{tab:riesgos}\\
\toprule
\textbf{ID} & \textbf{Riesgo} & \textbf{Impacto} & \textbf{Mitigación y disparador de acción} \\
\midrule
\endfirsthead
\toprule
\textbf{ID} & \textbf{Riesgo} & \textbf{Impacto} & \textbf{Mitigación y disparador de acción} \\
\midrule
\endhead
\midrule
\contlinea{4} \\
\endfoot
\bottomrule
\endlastfoot
RA-01 & \textbf{La documentación va muy por delante del código.} Hay cuarenta y dos escenarios, un SRS completo, doce ADR vigentes y un catálogo de setenta y nueve entradas frente a un sistema cuyo código empieza en el Sprint 2. & Alto & Los Sprints 1 y 2 son fundacionales por diseño. El disparador de alarma es el cierre de cada sprint: si al cierre del Sprint 2 los doce contenedores de Desarrollo no levantan con un comando, el Sprint 3 no cabe y el recorte se decide en su Planning, no en la Semana 10. \\
RA-02 & \textbf{Concentración de verificación en el Sprint 3.} Diecisiete de los cuarenta y dos escenarios se vuelven verificables allí, doce de ellos de banda crítica o alta, en el mismo sprint en que se construyen el gateway, la autorización territorial y el primer flujo completo. & Alto & La estimación de los frentes de ese sprint ---gateway con sus cinco responsabilidades, emisión de token en el Servicio de Identidad y validación en dos plataformas, filtrado territorial en los cuatro servicios con rutas publicadas, primer flujo completo de M1 a M4, conexión del frontend con sus estados de red, pantallas de sesión y verificación de los diecisiete escenarios--- arroja entre 144 y 197 horas frente a las 126 de capacidad declarada. El sprint está sobrecomprometido entre un 14 y un 56\,\% antes de empezar. En consecuencia: no se compromete alcance funcional adicional, y el recorte se decide en el Planning del Sprint 3 y no dentro de él. Se recorta por categoría MoSCoW completa ---lo \emph{Could} y lo \emph{Should} enteros--- y nunca reduciendo uniformemente cada historia: recortar parejo produce muchas cosas a medias y ninguna demostrable. Los frentes no recortables son el gateway y la emisión de identidad, porque sostienen K-07. \\
RA-03 & \textbf{Curva de aprendizaje de .NET.} ADR-014 asigna a .NET cuatro componentes ---el gateway y los servicios de Identidad, Institucional y Campañas---, sobre la plataforma que el equipo domina menos, y el Servicio de Identidad emite la identidad de todo el sistema. & Medio & Plantilla común de servicio .NET, con autenticación, salud, métricas y correlación, compartida por los cuatro componentes; trabajo en pareja en el Servicio de Identidad; y el \emph{spike} de aprendizaje antes de comprometer funcionalidad sobre la plataforma. \\
RA-04 & \textbf{El gateway es punto único de fallo y de latencia} (PS-01). Concentra cuatro responsabilidades de seguridad y toda petición pasa por él. & Medio & Se acepta: la alternativa ---duplicar el control de acceso en cada servicio--- multiplica la superficie donde el control puede implementarse mal. La mitigación es el sondeo de salud con reinicio automático y la medición de latencia en el propio gateway. \\
RA-05 & \textbf{Puesta en marcha cerca de su límite.} Doce contenedores en Desarrollo y catorce en QA, más cuatro trabajos de migración, dejan menos margen a los diez minutos de EC-23 y a la memoria de la máquina de QA (ADR-014). & Medio & Imágenes precompiladas descargadas del registro en lugar de construidas en el arranque, límites de memoria declarados por contenedor, y medición del tiempo de arranque en las tres plataformas y del consumo en reposo al cierre de cada sprint. Si la medición supera diez minutos, se reabre ADR-014; el escenario no se relaja. \\
RA-06 & \textbf{Las pantallas de sesión no existen en el prototipo.} Inicio de sesión, credenciales incorrectas y sesión expirada están respaldadas desde el SRS V3.0 por RF-23 y RF-24, pero ninguna pantalla las realiza todavía. & Medio & Se construyen en el Sprint 3 junto con el Servicio de Identidad, porque son la condición para demostrar EC-03 y el criterio de cierre de ese sprint desde fuera de la red institucional. \\
RA-07 & \textbf{Ninguna pantalla del prototipo tiene estados de red.} No hay carga por latencia, error de servidor, reintento ni conflicto de concurrencia, porque hoy no hay red. & Bajo & Cada pantalla que consuma un \emph{endpoint} necesitará al menos tres estados que hoy no existen. Es trabajo real y se estima en la planeación del Sprint 3, no dentro de él. \\
RA-08 & \textbf{La consulta de auditoría atraviesa cuatro contextos.} Es la operación del sistema que más fronteras cruza en una sola acción de usuario (ADR-016). & Medio & Llamadas en paralelo con plazo de 3 segundos, respuesta parcial marcada con los contextos ausentes, y casos de EC-05 con cada contexto fuera de servicio. \\
RA-09 & \textbf{Denegaciones del gateway que no llegan a la bitácora.} La entrega al Servicio de Identidad es asíncrona para no retrasar el rechazo, y un evento puede perderse si Identidad no responde durante más de sesenta segundos. & Bajo & Cola en memoria de mil eventos con reintento, registro estructurado del evento no entregado con su identificador de correlación y métrica propia; el rechazo nunca depende de que su registro se complete. \\
\end{longtable}

\subsection{Deuda arquitectónica declarada}

Se registra lo que se decidió no hacer, para que sea una decisión y no un olvido.

\begin{itemize}
\item \textbf{No hay orquestación en clúster.} La restricción pide contenerización. Escalar horizontalmente el Servicio de Donación es posible detrás del gateway, pero se hace a mano con la orden de escalado de réplicas de la herramienta de composición.
\item \textbf{No hay intermediario de mensajería.} Las notificaciones se desacoplan con una bandeja de salida transaccional en el Servicio de Donación, consumida por contrato por un trabajador sin base propia. Es suficiente para EC-11 y evita operar una pieza de infraestructura que el equipo no conoce.
\item \textbf{No hay proveedor de identidad completo.} La emisión del token reside en el Servicio de Identidad, con la justificación registrada en ADR-008. No hay autoservicio de contraseña ni federación con proveedores externos.
\item \textbf{No hay revocación inmediata del token de acceso.} Una sesión revocada conserva su token de acceso vigente hasta quince minutos (ADR-008).
\item \textbf{No hay agregación centralizada de registros.} Los registros se consultan por contenedor y se correlacionan por identificador; su agregación queda en anillo \emph{Evaluar} del catálogo.
\item \textbf{No hay ambiente de Producción.} Está declarado con su topología de referencia y no se construye.
\item \textbf{No se contempla internacionalización.} El sistema es nacional y de una sola lengua. Probablemente no hace falta; se deja dicho para que sea una decisión.
\item \textbf{La accesibilidad está diseñada pero no auditada.} Foco visible, objetivos táctiles, iconografía propia y movimiento reducido respetado. La verificación formal contra WCAG 2.2 AA con lector de pantalla y sin ratón es EC-18 y está programada para el Sprint 5.
\end{itemize}

% =====================================================================
%  12. GOBIERNO DEL DOCUMENTO
% =====================================================================
\section{Gobierno y evolución del documento}

\subsection{Quién puede cambiar qué}

\begin{longtable}{L{4.4cm} L{4.0cm} L{6.2cm}}
\caption{Responsabilidad sobre los contenidos de este documento.}\label{tab:gobierno}\\
\toprule
\textbf{Contenido} & \textbf{Quién lo cambia} & \textbf{Condición} \\
\midrule
\endfirsthead
\toprule
\textbf{Contenido} & \textbf{Quién lo cambia} & \textbf{Condición} \\
\midrule
\endhead
\midrule
\contlinea{3} \\
\endfoot
\bottomrule
\endlastfoot
Estilo arquitectónico, fronteras entre servicios y propiedad de datos & Mesa de Arquitectura & Solo mediante ADR nuevo o reemplazo de uno existente. Nunca dentro de un Pull Request. \\
Índice de prioridad de un escenario y sus tres factores & Arquitecta de Software y Product Owner & Todo cambio de factor se registra con su justificación. El índice no se ajusta para acomodar la capacidad de un sprint. \\
Escenarios de calidad y medidas de respuesta & Arquitecta de Software, con visto bueno de QA & Toda medida nueva debe ser un número o una condición binaria, y toda subcaracterística de la norma debe conservar al menos un escenario. \\
Tácticas y componente donde residen & Arquitecta de Software & Debe conservarse la trazabilidad de la Tabla \ref{tab:matriztraza}. \\
Topología y ambientes & DevOps & Ninguna tecnología fuera del catálogo de herramientas. \\
Drivers y prioridades de calidad & Product Owner y Arquitecta & Un cambio de prioridad de atributo es un cambio de arquitectura y se documenta como tal. \\
\end{longtable}

\subsection{Evolución prevista}

\begin{longtable}{L{1.6cm} L{2.2cm} L{10.6cm}}
\caption{Evolución prevista de este documento a lo largo del proyecto.}\label{tab:evolucion}\\
\toprule
\textbf{Versión} & \textbf{Entrega} & \textbf{Qué incorpora} \\
\midrule
\endfirsthead
\toprule
\textbf{Versión} & \textbf{Entrega} & \textbf{Qué incorpora} \\
\midrule
\endhead
\midrule
\contlinea{3} \\
\endfoot
\bottomrule
\endlastfoot
V1 & Semana 6 & Drivers, killers, atributos con cobertura completa de la norma, cuarenta y dos escenarios priorizados, compromisos, ADR aprobados y las cuatro arquitecturas. \\
V2 & Semana 8 & Esta versión. ADR-007 a ADR-016: cinco servicios, zonas de confianza y aristas permitidas, proxy de borde, token con firma asimétrica, bitácora en cuatro contextos y escalamiento como proceso programado. Ambientes y seguridad de la infraestructura, cuya configuración detallada está en el Documento de Infraestructura V1. Escenarios corregidos con la revisión cruzada del catálogo. La prueba de concepto y su \emph{benchmarking} se documentan en el SDD V1. \\
V3 & Semana 10 & Resultado real de la verificación de los diecisiete escenarios del Sprint 3. Ajuste de las medidas de respuesta que la medición contradiga, con registro del ajuste. \\
V4 & Semana 12 & Escenarios de carga y de sobrecarga verificados con medición real sobre el volumen sintético declarado. Revisión de los objetivos de disponibilidad y de capacidad con datos, no con estimación. \\
V5 & Semana 14 & Ensayo de restauración y auditoría de accesibilidad. Cierre de la deuda declarada o su registro explícito como no resuelta. \\
\end{longtable}

\noindent Una versión posterior que corrija una afirmación de esta no es un defecto del documento: es el mecanismo funcionando. Lo que no puede ocurrir es que una medida de respuesta se ajuste al resultado obtenido sin dejar constancia de que se ajustó.

% =====================================================================
%  ANEXOS
% =====================================================================
\appendix

\section{Correspondencia de identificadores de módulo}
\label{sec:anexoA}

El catálogo de módulos fue reestructurado el 27 de agosto de 2026 conforme al criterio de contexto acotado y propiedad de datos. Documentación anterior a esa fecha usa la numeración previa. La correspondencia vigente es la del SRS V3.0, que se reproduce aquí para que este documento sea legible junto a artefactos antiguos.

\begin{longtable}{L{1.6cm} L{5.6cm} L{7.0cm}}
\caption{Correspondencia entre la numeración anterior de módulos y la vigente.}\label{tab:correspondencia}\\
\toprule
\textbf{Anterior} & \textbf{Denominación anterior} & \textbf{Vigente} \\
\midrule
\endfirsthead
\toprule
\textbf{Anterior} & \textbf{Denominación anterior} & \textbf{Vigente} \\
\midrule
\endhead
\midrule
\contlinea{3} \\
\endfoot
\bottomrule
\endlastfoot
M1 & Gestión de Donantes & M1, con el reconocimiento no monetario incorporado como característica. \\
M2 & Motivación y Gamificación & Suprimido como módulo. Su funcionalidad se especifica en RF-11, dentro de M1. \\
M3 & Gestión de Campañas & M2. \\
M4 & Trazabilidad y Ciclo de Vida & M3. \\
M5 & Inventario y Alertas & M4, con la exclusión automática de unidades vencidas incorporada. \\
M6 & Red de Transferencias & M5, con la movilización de donantes incorporada. \\
M7 & Administración Territorial & M6, acotado a la jerarquía y al registro institucional. La restricción de visibilidad se traslada a RNF-02 y a ST3. \\
M8 & Analítica e Indicadores & M7. \\
M9 & Extensibilidad de Dominio & Suprimido como módulo. Se especifica en RNF-05 y en RI-01. \\
\end{longtable}

\noindent\textbf{Nota de consistencia.} ADR-005 cita el módulo M9 y un requerimiento que ya no existen tras la reestructuración. ADR-012 deroga esa referencia sin reescribir la decisión, que sigue vigente: base relacional con esquema independiente por servicio.

\section{Cuestiones de arquitectura cerradas}
\label{sec:anexoB}

Este anexo consolida las cuestiones que la V1.0 dejó abiertas y la decisión que cierra cada una. Ninguna llega abierta al Sprint 3.

\begin{longtable}{L{1.3cm} L{4.6cm} L{7.0cm} L{1.5cm}}
\caption{Cuestiones de arquitectura de la V1.0 y decisión que las cierra.}\label{tab:consolidadopendientes}\\
\toprule
\textbf{ID} & \textbf{Cuestión} & \textbf{Decisión} & \textbf{Registro} \\
\midrule
\endfirsthead
\toprule
\textbf{ID} & \textbf{Cuestión} & \textbf{Decisión} & \textbf{Registro} \\
\midrule
\endhead
\midrule
\contlinea{4} \\
\endfoot
\bottomrule
\endlastfoot
P-01 & Selección del API Gateway. & YARP, con compuerta derivada de la matriz de autorización, límite de tasa en dos escalas y registro de denegaciones en el Servicio de Identidad. & ADR-007 \\
P-02 & Emisión y validación del token. & Emisor único en el Servicio de Identidad, firma RS256, token de acceso de quince minutos y sesión de renovación de ocho horas. Se descarta un proveedor de identidad completo por costo de operación y de aprendizaje. & ADR-008 \\
P-03 & Línea de soporte de la plataforma Java. & Java 25 con Eclipse Temurin y Spring Boot 4.1. & ADR-009 \\
P-04 & Orquestación del escalamiento. & Proceso programado del Servicio de Donación con alcance departamental y autoridad de sistema. & ADR-013 \\
P-05 & Qué ve el donante sobre el resultado de su propia donación. & El donante ve que donó, la fecha, el banco y su elegibilidad vigente, calculada solo a partir del intervalo desde su última donación; no ve el veredicto de tamizaje de sus unidades. Cada donación de su historial muestra un mismo mensaje neutro, idéntico para todos los donantes cualquiera que sea el destino de la unidad: el banco de sangre lo contactará por su canal clínico si necesita comunicarle algo sobre esa donación. Así, ni la presencia ni la ausencia de un texto permite inferir el resultado. El sistema no comunica resultados clínicos. Es coherente con RNI-01 y con el contrato de consulta del historial propio del DD. & Product Owner \\
P-06 & La reserva de cupo vive en M1 pero el cupo es un concepto de M2, en otro servicio. & El cupo y sus reservas pertenecen al Servicio de Campañas, dueño del concepto: la reserva se crea y expira allí, con su propio proceso programado, y el Servicio de Donación no guarda contadores de cupo. Una reserva no requiere transacción distribuida porque el contador y la reserva viven en la misma base. RF-15 pasa a M2. & SRS V3.0 \\
P-07 & Asignación de jurisdicción a un usuario. & RF-22, en el Servicio de Identidad, con efecto en el siguiente token de acceso emitido. & SRS V3.0, ADR-008 \\
P-08 & Vista de trazabilidad del perfil auditor. & Se construye sobre la consulta compuesta \texttt{GET /v1/auditoria} del Servicio Institucional. & ADR-016 \\
\end{longtable}

\vspace{0.4cm}
\noindent{\footnotesize\textcolor{grisTexto}{Este documento describe un sistema académico que opera exclusivamente sobre datos sintéticos. Los objetivos de disponibilidad, recuperación y capacidad aquí declarados corresponden a esa escala y no constituyen un compromiso de nivel de servicio de un sistema en producción.}}

\end{document}